\chapter{Geodesics and Distance}
\label{chap:pet-geodesics-and-distance}

Faithful transcription of Petersen, \emph{Riemannian Geometry} (GTM 171, 3rd ed.),
Chapter 5. The chapter develops the theory of geodesics — curves of vanishing
acceleration — via the ordinary-differential-equations existence and uniqueness
theory for the geodesic equation in coordinates, and studies how far geodesics
extend (geodesic completeness). It then builds the Riemannian distance function
from curve lengths, shows the resulting metric topology agrees with the manifold
topology, and develops the first variation formula of energy, characterizing
geodesics as the stationary curves for the energy functional and segments as
geodesics. The exponential map and Riemannian normal coordinates are introduced,
culminating in the Gauss Lemma and the theorem that sufficiently short geodesics
are the unique distance-minimizing segments. The technical starting point is a
characterization, via two properties analogous to torsion-freeness and metric
compatibility of the Levi-Civita connection, of mixed second partial derivatives
of maps into $M$, which underlies the definition of the acceleration of a curve;
this machinery, and the definition of a geodesic as a curve of vanishing
acceleration, is purely Riemannian, but degenerates in the pseudo-Riemannian
setting only at the point of defining arclength and distance, since the energy
functional still makes sense there. Riemannian isometries are studied
next: local isometries are shown to be uniquely determined by their value and
differential at one point, the fixed-point sets of isometries are shown to be
totally geodesic, complete simply connected constant-curvature manifolds are
classified (Killing–Hopf), the Myers–Steenrod theorem shows that distance-preserving
bijections are Riemannian isometries and that the isometry group is a Lie group,
and the slice theorem describes the local structure of isometric group actions.
The chapter closes with completeness: the Hopf–Rinow theorem, Tashiro's global
warped-product characterization of manifolds admitting a function with conformal
Hessian, and the theory of the segment domain, cut locus, and injectivity radius.

\section{Mixed Partials}

For a curve $c:\Omega\to M$ on an open $\Omega\subset\mathbb R^m$ (with
coordinates $t^i$, or $r,s,t,u,\dots$, on $\Omega$ to avoid clashing with the
coordinates $x^i$ reserved for $M$), the \emph{first partials} $\partial c/\partial t^i$
are simply the velocity fields of $t^i\mapsto c(t^1,\dots,t^i,\dots,t^m)$ with the
remaining coordinates fixed, valued in $TM$. To define \emph{second partials}
$\partial^2c/\partial t^i\partial t^j$ that again take values in $TM$ (rather than
in $TTM$), two natural requirements are imposed:
\begin{enumerate}
\item[(1)] equality of mixed partials,
$\dfrac{\partial^2c}{\partial t^i\partial t^j}=\dfrac{\partial^2c}{\partial t^j\partial t^i}$;
\item[(2)] the product rule,
$\dfrac{\partial}{\partial r^k}g\Big(\dfrac{\partial c}{\partial r^i},\dfrac{\partial c}{\partial r^j}\Big)
=g\Big(\dfrac{\partial^2c}{\partial r^k\partial r^i},\dfrac{\partial c}{\partial r^j}\Big)
+g\Big(\dfrac{\partial c}{\partial r^i},\dfrac{\partial^2c}{\partial r^k\partial r^j}\Big)$.
\end{enumerate}
Property (1) is analogous to torsion-freeness of the connection and (2) to
metric compatibility; as with the characterization of the Levi-Civita connection
by these same two properties, they force a Koszul-type formula for second
partials, established in the proof of the following lemma.
\source{petersen:page-0183}

\begin{lemma}[Lemma 5.1.1 — uniqueness of mixed partials]
  \label{lem:pet-ch5-mixed-partials-uniqueness}
  \lean{PetersenLib.mixedPartials_uniqueness}
  \leanok
  There is at most one way of defining mixed partials of maps $c:\Omega\to M$ so
  that properties (1) and (2) above hold.
  \source{petersen:page-0184}
\end{lemma}
\begin{proof}
  \leanok
  \uses{lem:pet-ch5-mixed-partials-uniqueness}
  First, (1) and (2) force the Koszul-type formula
  \[
    2g\Big(\frac{\partial^2c}{\partial r^i\partial r^j},\frac{\partial c}{\partial r^k}\Big)
    =\frac{\partial}{\partial r^i}g\Big(\frac{\partial c}{\partial r^j},\frac{\partial c}{\partial r^k}\Big)
    +\frac{\partial}{\partial r^j}g\Big(\frac{\partial c}{\partial r^k},\frac{\partial c}{\partial r^i}\Big)
    -\frac{\partial}{\partial r^k}g\Big(\frac{\partial c}{\partial r^i},\frac{\partial c}{\partial r^j}\Big).
  \]
  Indeed, expanding each of the three terms on the right by (2) and using (1) to
  identify $\partial^2c/\partial r^i\partial r^j=\partial^2c/\partial r^j\partial r^i$ and its
  cyclic relatives, the six resulting terms
  \[
    g\Big(\frac{\partial^2c}{\partial r^i\partial r^j},\frac{\partial c}{\partial r^k}\Big)
    +g\Big(\frac{\partial c}{\partial r^j},\frac{\partial^2c}{\partial r^i\partial r^k}\Big)
    +g\Big(\frac{\partial^2c}{\partial r^j\partial r^k},\frac{\partial c}{\partial r^i}\Big)
    +g\Big(\frac{\partial c}{\partial r^k},\frac{\partial^2c}{\partial r^j\partial r^i}\Big)
    -g\Big(\frac{\partial^2c}{\partial r^k\partial r^i},\frac{\partial c}{\partial r^j}\Big)
    -g\Big(\frac{\partial c}{\partial r^i},\frac{\partial^2c}{\partial r^k\partial r^j}\Big)
  \]
  pair off (via (1)) into $g(\partial c/\partial r^j,\partial^2c/\partial r^i\partial r^k)$
  cancelling $-g(\partial c/\partial r^i,\partial^2c/\partial r^k\partial r^j)$, and
  $g(\partial^2c/\partial r^j\partial r^k,\partial c/\partial r^i)$ cancelling
  $-g(\partial^2c/\partial r^k\partial r^i,\partial c/\partial r^j)$, leaving exactly
  $2g(\partial^2c/\partial r^i\partial r^j,\partial c/\partial r^k)$.

  Now fix $p$ in the image of $c$ and $v\in T_pM$ arbitrary. Adding an extra
  parameter $t^0$ to $c$ gives a map
  $\tilde c:(-\varepsilon,\varepsilon)\times\Omega\to M$ with
  $\partial\tilde c/\partial t^0|_p=v$; this extension exists independently of how
  mixed partials are defined. Applying the Koszul formula at $p$ with $k=0$,
  $r^0=t^0$, expresses $g(\partial^2c/\partial t^i\partial t^j,v)$, for every
  $v\in T_pM$, purely in terms of $g$ and first partials of $c,\tilde c$ — data
  independent of the definition of mixed partials. Since $g$ is nondegenerate,
  this determines $\partial^2c/\partial t^i\partial t^j|_p\in T_pM$ uniquely.
  \source{petersen:page-0184,petersen:page-0185}
\end{proof}

\begin{theorem}[Thm. 5.1.2 — existence of mixed partials]
  \label{thm:pet-ch5-mixed-partials-existence}
  \lean{PetersenLib.mixedPartials_existence}
  \leanok
  \uses{lem:pet-ch5-mixed-partials-uniqueness}
  It is possible to define mixed partials of maps $c:\Omega\to M$ in a coordinate
  system so that properties (1) and (2) of \cref{lem:pet-ch5-mixed-partials-uniqueness}
  hold.
  \source{petersen:page-0185}
\end{theorem}
\begin{proof}
  \leanok
  \uses{thm:pet-ch5-mixed-partials-existence}
  By \cref{lem:pet-ch5-mixed-partials-uniqueness}, if two maps $c_1,c_2:\Omega\to M$
  agree near a point of $\Omega$, the Koszul formula gives the same right-hand
  side for both, so there is no loss of generality in assuming the image of $c$
  lies in a single coordinate neighborhood; a local, coordinate-dependent
  definition satisfying (1) and (2) is then automatically coordinate independent.

  Let $c:\Omega\to U\subset M$, $U$ a coordinate neighborhood, with parameters
  $s,t$ and $c=(c^1,\dots,c^n)$ in the coordinates on $U$. The velocity in the $s$
  direction is $\partial c/\partial s=(\partial c^i/\partial s)\partial_i$. Define, for a
  vector field $X$ along $c$, $\partial X/\partial t|_p=\nabla_{c'(t)}X$ (where
  $c(t)=p$); then
  \[
    \frac{\partial}{\partial t}\frac{\partial c}{\partial s}
    =\frac{\partial}{\partial t}\Big(\frac{\partial c^i}{\partial s}\partial_i\Big)
    =\frac{\partial^2c^i}{\partial t\partial s}\partial_i+\frac{\partial c^i}{\partial s}\nabla_{\partial c/\partial t}\partial_i
    =\frac{\partial^2c^k}{\partial t\partial s}\partial_k
    +\frac{\partial c^i}{\partial s}\frac{\partial c^j}{\partial t}\Gamma^k_{ji}\partial_k.
  \]
  Accordingly define
  \[
    \frac{\partial^2c}{\partial t\partial s}
    :=\Big(\frac{\partial^2c^k}{\partial t\partial s}+\frac{\partial c^i}{\partial s}\frac{\partial c^j}{\partial t}\Gamma^k_{ji}\Big)\partial_k.
  \]
  Property (1) holds: $\partial^2c^k/\partial t\partial s$ is symmetric in $s,t$ by
  Clairaut's theorem on ordinary mixed partials, and $\Gamma^k_{ji}$ is symmetric
  in $i,j$, so the whole expression is symmetric under exchanging $s\leftrightarrow t$
  together with $i\leftrightarrow j$.

  For property (2), write $u$ for a third parameter and use the metric-compatibility
  identity for the Christoffel symbols, $\partial_kg_{ij}=\Gamma_{k,ij}+\Gamma_{k,ji}$
  (lowered indices), together with $g_{ij}\Gamma^i_{kl}=\Gamma_{j,kl}$. Then
  \begin{align*}
    \frac{\partial}{\partial t}g\Big(\frac{\partial c}{\partial s},\frac{\partial c}{\partial u}\Big)
    &=\frac{\partial}{\partial t}\Big(g_{ij}\frac{\partial c^i}{\partial s}\frac{\partial c^j}{\partial u}\Big)\\
    &=(\partial_kg_{ij})\frac{\partial c^k}{\partial t}\frac{\partial c^i}{\partial s}\frac{\partial c^j}{\partial u}
    +g_{ij}\frac{\partial^2c^i}{\partial t\partial s}\frac{\partial c^j}{\partial u}
    +g_{ij}\frac{\partial c^i}{\partial s}\frac{\partial^2c^j}{\partial t\partial u}\\
    &=(\Gamma_{k,ij}+\Gamma_{k,ji})\frac{\partial c^k}{\partial t}\frac{\partial c^i}{\partial s}\frac{\partial c^j}{\partial u}
    +g_{ij}\frac{\partial^2c^i}{\partial t\partial s}\frac{\partial c^j}{\partial u}
    +g_{ij}\frac{\partial c^i}{\partial s}\frac{\partial^2c^j}{\partial t\partial u}\\
    &=g_{ij}\Big(\frac{\partial^2c^i}{\partial t\partial s}+\Gamma^i_{kl}\frac{\partial c^k}{\partial t}\frac{\partial c^l}{\partial s}\Big)\frac{\partial c^j}{\partial u}
    +g_{ij}\frac{\partial c^i}{\partial s}\Big(\frac{\partial^2c^j}{\partial t\partial u}+\Gamma^j_{kl}\frac{\partial c^k}{\partial t}\frac{\partial c^l}{\partial u}\Big)\\
    &=g\Big(\frac{\partial^2c}{\partial t\partial s},\frac{\partial c}{\partial u}\Big)
    +g\Big(\frac{\partial c}{\partial s},\frac{\partial^2c}{\partial t\partial u}\Big),
  \end{align*}
  using the symmetry $\Gamma^k_{ji}=\Gamma^k_{ij}$ to match indices with the
  definition of $\partial^2c/\partial t\partial s$ and $\partial^2c/\partial t\partial u$
  above.
  \source{petersen:page-0185,petersen:page-0186}
\end{proof}

In case $M\subset\tilde M$ it is often convenient to calculate mixed partials in
$\tilde M$ first and then project them onto $M$. For $v\in T_p\tilde M$, $p\in M$,
write $v=v^\top+v^\perp$ for the decomposition into tangential $T_pM$ and normal
$T_p^\perp M$ components.

\begin{proposition}[Prop. 5.1.3 — mixed partials in submanifolds]
  \label{prop:pet-ch5-mixed-partials-submanifold}
  \lean{PetersenLib.mixedPartialSubmanifoldProjection}
  \uses{thm:pet-ch5-mixed-partials-existence}
  If $c:\Omega\to M\subset\tilde M$ and $\frac{\partial^2c}{\partial r\partial\theta}\in T_p\tilde M$
  is the mixed partial computed in $\tilde M$, then
  \[
    \left(\frac{\partial^2c}{\partial r\partial\theta}\right)^\top\in T_pM
  \]
  is the mixed partial computed in $M$.
  \source{petersen:page-0187}
\end{proposition}
\begin{proof}
  \uses{prop:pet-ch5-mixed-partials-submanifold}
  Let $\tilde g$ be the metric on $\tilde M$ and $g$ its restriction to $M$. The
  mixed partial in $\tilde M$ satisfies, for the coordinate curve
  $c:\Omega\to M\subset\tilde M$ with parameters $r,\theta,k$,
  \[
    2\tilde g\!\left(\frac{\partial^2c}{\partial r\partial\theta},\frac{\partial c}{\partial k}\right)
    =\frac{\partial}{\partial r}\tilde g\!\left(\frac{\partial c}{\partial\theta},\frac{\partial c}{\partial k}\right)
    +\frac{\partial}{\partial\theta}\tilde g\!\left(\frac{\partial c}{\partial k},\frac{\partial c}{\partial r}\right)
    -\frac{\partial}{\partial k}\tilde g\!\left(\frac{\partial c}{\partial r},\frac{\partial c}{\partial\theta}\right).
  \]
  Since $\partial c/\partial r,\partial c/\partial\theta,\partial c/\partial k\in TM$, the right-hand side
  depends only on $g=\tilde g|_{TM}$, and replacing the left-hand mixed partial by
  its tangential projection does not change the pairing against $\partial c/\partial k\in TM$:
  \[
    2\tilde g\!\left(\frac{\partial^2c}{\partial r\partial\theta},\frac{\partial c}{\partial k}\right)
    =2\tilde g\!\left(\left(\frac{\partial^2c}{\partial r\partial\theta}\right)^{\!\top},\frac{\partial c}{\partial k}\right).
  \]
  As $\partial c/\partial k$ ranges over $T_pM$ this identifies
  $(\partial^2c/\partial r\partial\theta)^\top$ as exactly the vector characterized by the
  mixed-partial defining identity within $M$, i.e.\ the mixed partial of $c$ computed
  intrinsically in $M$.
  \source{petersen:page-0187}
\end{proof}

\section{Geodesics}

\begin{definition}[acceleration of a curve]
  \label{def:pet-ch5-acceleration}
  \lean{PetersenLib.curveAcceleration}
  \leanok
  For a $C^\infty$ curve $c:I\to M$ the \emph{acceleration} is $\ddot c=\frac{d^2c}{dt^2}$.
  In local coordinates,
  \[
    \ddot c=\frac{d^2c^k}{dt^2}\partial_k+\frac{dc^l}{dt}\frac{dc^j}{dt}\Gamma^k_{lj}\partial_k.
  \]
  \source{petersen:page-0187,petersen:page-0188}
\end{definition}

\begin{definition}[geodesic; arclength parametrization]
  \label{def:pet-ch5-geodesic}
  \lean{PetersenLib.IsGeodesic,PetersenLib.geodesicCoordinateEquation,PetersenLib.IsChartGeodesicOn}
  \leanok
  \uses{def:pet-ch5-acceleration}
  A $C^\infty$ curve $c:I\to M$ with vanishing acceleration, $\ddot c=0$, is a
  \emph{geodesic}. If $c$ is a geodesic its speed $|\dot c|=\sqrt{g(\dot c,\dot c)}$
  is constant, since $\frac{d}{dt}g(\dot c,\dot c)=2g(\ddot c,\dot c)=0$; equivalently
  $c$ is parametrized proportionally to arc length. If $|\dot c|\equiv1$, $c$ is
  \emph{parametrized by arclength}. In local coordinates on $U\subset M$, $c:I\to U$
  is a geodesic iff its components satisfy the second-order system
  \[
    \ddot c^k(t)=-\dot c^l(t)\dot c^j(t)\,\Gamma^k_{lj}|_{c(t)},\qquad k=1,\dots,n,
  \]
  giving the initial-value problem $c(0)=q$, $\dot c(0)=\dot c^i(0)\partial_i|_q$ for
  specified $q\in U$ and initial velocity.
  \source{petersen:page-0188}
\end{definition}

\begin{remark}[Remark 5.2.1 — distance functions produce geodesics]
  \label{rem:pet-ch5-distance-function-geodesics}
  \lean{PetersenLib.distanceFunctionIntegralCurvesAreGeodesics}
  \uses{def:pet-ch5-geodesic}
  If $r:U\to\mathbb R$ is a distance function then $\nabla_{\partial_r}\partial_r=0$, where
  $\partial_r=\nabla r$; consequently the integral curves of $\nabla r$ are geodesics.
  The theory of geodesics is developed independently of distance functions and is
  ultimately used to establish their existence.
  \source{petersen:page-0188}
\end{remark}

\begin{theorem}[Thm. 5.2.2 — local uniqueness of geodesics]
  \label{thm:pet-ch5-local-uniqueness}
  \lean{PetersenLib.geodesic_local_uniqueness}
  \leanok
  \uses{def:pet-ch5-geodesic}
  Let $I_1,I_2$ be intervals with $t_0\in I_1\cap I_2$. If $c_1:I_1\to U$ and
  $c_2:I_2\to U$ are geodesics in a chart $U\subset M$ with $c_1(t_0)=c_2(t_0)$ and
  $\dot c_1(t_0)=\dot c_2(t_0)$, then $c_1|_{I_1\cap I_2}=c_2|_{I_1\cap I_2}$.
  \source{petersen:page-0189}
\end{theorem}
\begin{proof}
  \leanok
  \uses{thm:pet-ch5-local-uniqueness}
  In the coordinates of $U$ write $u_i=\varphi\circ c_i$; the geodesic equation
  is the second-order system $\ddot u=-\Gamma(\dot u,\dot u)\circ u$, i.e.\ the
  first-order system $z'=F(z)$ for $z=(u,\dot u)$ with
  $F(x,w)=(w,-\Gamma(w,w)(x))$. The Christoffel symbols are smooth on $U$, so
  $F$ is $C^1$ on the open set $\varphi(U)\times\mathbb R^n$ and in particular
  locally Lipschitz. The set $A=\{t\in I_1\cap I_2\mid z_1(t)=z_2(t)\}$
  contains $t_0$, is closed in $I_1\cap I_2$ by continuity, and is open since
  at any $t_1\in A$ the Picard--Lindel\"of/Gr\"onwall uniqueness theorem for
  locally Lipschitz fields forces $z_1=z_2$ near $t_1$. As $I_1\cap I_2$ is an
  interval, $A=I_1\cap I_2$, and applying $\varphi^{-1}$ gives $c_1=c_2$ there.
  \source{petersen:page-0189}
\end{proof}

\begin{theorem}[Thm. 5.2.3 — local existence of geodesics]
  \label{thm:pet-ch5-local-existence}
  \lean{PetersenLib.geodesic_local_existence}
  \uses{def:pet-ch5-geodesic}
  For each $p\in U$ and $v\in\mathbb R^n$ there are a neighborhood $V_1$ of $p$, a
  neighborhood $V_2$ of $v$, and $\varepsilon>0$ such that for each $q\in V_1$,
  $w\in V_2$, there is a geodesic $c_{q,w}:(-\varepsilon,\varepsilon)\to U$ with $c(0)=q$,
  $\dot c(0)=w^i\partial_i|_q$. Moreover $(q,w,t)\mapsto c_{q,w}(t)$ is $C^\infty$ on
  $V_1\times V_2\times(-\varepsilon,\varepsilon)$.
  \source{petersen:page-0189}
\end{theorem}
\begin{proof}
  \uses{thm:pet-ch5-local-existence}
  As in the proof of \cref{thm:pet-ch5-local-uniqueness}, the geodesic
  equation in the coordinates of $U$ is the first-order system $z'=F(z)$,
  $F(x,w)=(w,-\Gamma(w,w)(x))$, with $F$ smooth on
  $\varphi(U)\times\mathbb R^n$. Picard--Lindel\"of with parameters gives, for
  the initial condition $z_0=(\varphi(p),v)$, radii $r,\varepsilon>0$ and a
  local flow $Z(z,t)$ defined for $|t|<\varepsilon$ and all initial conditions
  $z$ in the ball of radius $r$ about $z_0$; the neighborhoods are
  $V_1=\varphi^{-1}(B(\varphi(p),r/2))$ and $V_2=B(v,r/2)$, and
  $c_{q,w}=\varphi^{-1}(Z((\varphi(q),w),\cdot)_1)$ solves the geodesic
  equation with $c_{q,w}(0)=q$, $\dot c_{q,w}(0)=w$. Smooth dependence of
  solutions of ODEs with smooth right-hand side on their initial conditions
  makes $(q,w,t)\mapsto c_{q,w}(t)$ jointly $C^\infty$ on
  $V_1\times V_2\times(-\varepsilon,\varepsilon)$.
  \source{petersen:page-0189}
\end{proof}

\begin{lemma}[Lemma 5.2.4 — global uniqueness of geodesics]
  \label{lem:pet-ch5-global-uniqueness}
  \lean{PetersenLib.geodesic_global_uniqueness}
  \leanok
  \uses{thm:pet-ch5-local-uniqueness}
  Let $I_1,I_2$ be open intervals with $t_0\in I_1\cap I_2$. If $c_1:I_1\to M$ and
  $c_2:I_2\to M$ are geodesics with $c_1(t_0)=c_2(t_0)$ and $\dot c_1(t_0)=\dot c_2(t_0)$,
  then $c_1|_{I_1\cap I_2}=c_2|_{I_1\cap I_2}$.
  \source{petersen:page-0189}
\end{lemma}
\begin{proof}
  \uses{lem:pet-ch5-global-uniqueness}
  Let $A=\{t\in I_1\cap I_2\mid c_1(t)=c_2(t),\ \dot c_1(t)=\dot c_2(t)\}$. Then
  $t_0\in A$, and $A$ is closed by continuity of $c_1,c_2,\dot c_1,\dot c_2$. It is
  open: if $t_1\in A$, choose a coordinate chart $U$ around $c_1(t_1)=c_2(t_1)$; then
  $(t_1-\varepsilon,t_1+\varepsilon)\subset I_1\cap I_2$ for some $\varepsilon>0$ with
  $c_i|_{(t_1-\varepsilon,t_1+\varepsilon)}$ valued in $U$ for $i=1,2$, and
  \cref{thm:pet-ch5-local-uniqueness} gives $c_1=c_2$ on that interval, so
  $(t_1-\varepsilon,t_1+\varepsilon)\subset A$. As $I_1\cap I_2$ is connected, $A=I_1\cap I_2$.
  \source{petersen:page-0189}
\end{proof}

\begin{lemma}[uniform short-time existence over a compact set]
  \label{lem:pet-ch5-uniform-time-existence}
  \lean{PetersenLib.geodesic_uniformShortTimeExistence}
  \leanok
  \uses{thm:pet-ch5-local-existence}
  If $\tilde K\subset TM$ is compact, there is $\varepsilon>0$ such that for each
  $(q,v)\in\tilde K$ there is a geodesic $c:(-\varepsilon,\varepsilon)\to M$ with
  $c(0)=q$, $\dot c(0)=v$.
  \source{petersen:page-0190}
\end{lemma}
\begin{proof}
  \uses{thm:pet-ch5-local-existence}
  By \cref{thm:pet-ch5-local-existence}, every $(q,v)\in TM$ has a neighborhood
  $V_1\times V_2\subset TM$ and $\varepsilon_{(q,v)}>0$ such that all initial data in
  $V_1\times V_2$ produce geodesics defined on $(-\varepsilon_{(q,v)},\varepsilon_{(q,v)})$.
  These neighborhoods cover the compact set $\tilde K$, so finitely many
  $V_1^{(1)}\times V_2^{(1)},\dots,V_1^{(m)}\times V_2^{(m)}$ suffice; taking
  $\varepsilon=\min_i\varepsilon_{(q_i,v_i)}>0$ gives a single time on which all geodesics
  with initial data in $\tilde K$ exist, since every $(q,v)\in\tilde K$ lies in some
  $V_1^{(i)}\times V_2^{(i)}$.
  \source{petersen:page-0190}
\end{proof}

\begin{proposition}[maximal geodesics leave every compact set]
  \label{prop:pet-ch5-leaves-compact-set}
  \lean{PetersenLib.maximalGeodesic_leavesCompactSet}
  \leanok
  \uses{lem:pet-ch5-uniform-time-existence,def:pet-ch5-geodesic}
  Let $c:I=(a,b)\to M$ be a geodesic on a maximal interval of definition, with
  $b<\infty$. Then for every compact $K\subset M$ there is $t_K<b$ such that
  $t_K<t<b$ implies $c(t)\notin K$.
  \source{petersen:page-0190}
\end{proposition}
\begin{proof}
  \uses{lem:pet-ch5-uniform-time-existence,lem:pet-ch5-global-uniqueness}
  Suppose not: there is a compact $K$ and a sequence $t_1<t_2<\cdots$ in $I$ with
  $t_j\to b$ and $c(t_j)\in K$ for all $j$. Since $|\dot c|$ is constant, the
  velocities $\dot c(t_j)$ all lie in the set
  $\hat K=\{v_q\mid q\in K,\ v\in T_qM,\ |v|\le|\dot c|\}\subset TM$, which is
  compact: over a compact set $C$ inside a single chart, continuity and positive
  definiteness of the metric give a uniform bound $g_q(v,v)\ge\lambda\,|v|_{\mathrm{eucl}}^2$
  ($\lambda>0$ the minimum of the Gram pairing over $\varphi(C)\times S^{n-1}$), so the
  bounded-velocity set over $C$ is carried by the bundle trivialization onto a closed
  subset of $C\times\overline{B}(0,\sqrt{c/\lambda})$; finitely many such chart pieces
  cover $K$. By
  \cref{lem:pet-ch5-uniform-time-existence} there is $\varepsilon>0$ such that for
  each $v\in\hat K$ there is a geodesic $(-\varepsilon,\varepsilon)\to M$ with that initial
  velocity. Choose $t_j$ with $b-t_j<\varepsilon/2$; by
  \cref{lem:pet-ch5-global-uniqueness} the geodesic through $c(t_j)$ with velocity
  $\dot c(t_j)$ agrees with $c$ on the overlap of the intervals and patches together
  with $c$ to a geodesic on $(a, t_j+\varepsilon)$ with $t_j+\varepsilon>b$, contradicting
  maximality of $I$.
  \source{petersen:page-0190}
\end{proof}

\begin{corollary}[Corollary 5.2.5 — compact manifolds are geodesically complete]
  \label{cor:pet-ch5-compact-manifold-complete}
  \lean{PetersenLib.compactManifold_geodesicallyComplete}
  \leanok
  \uses{prop:pet-ch5-leaves-compact-set}
  If $M$ is a compact Riemannian manifold, then for each $p\in M$, $v\in T_pM$
  there is a geodesic $c:\mathbb R\to M$ with $c(0)=p$, $\dot c(0)=v$: geodesics exist
  for all time.
  \source{petersen:page-0190}
\end{corollary}
\begin{proof}
  \uses{prop:pet-ch5-leaves-compact-set,lem:pet-ch5-global-uniqueness,thm:pet-ch5-local-existence}
  Let $I$ be the union of all open intervals around $0$ carrying a geodesic with
  $c(0)=p$, $\dot c(0)=v$; it is open, an interval, and nonempty by local existence
  (\cref{thm:pet-ch5-local-existence}). Any two such geodesics agree on the overlap
  of their intervals by \cref{lem:pet-ch5-global-uniqueness}, so they patch to a
  well-defined maximal geodesic $c:I=(a,b)\to M$. If $b<\infty$, any geodesic
  extension of $c|_{(0,b)}$ past $b$ would glue with $c$ into a strictly larger
  admissible interval, contradicting the definition of $I$; so
  \cref{prop:pet-ch5-leaves-compact-set} applies to the compact set $K=M$ and forces
  $c(t)\notin M$ for $t$ near $b$, which is absurd. Hence $b=\infty$, and $a=-\infty$
  by applying the same argument to the reversed geodesic $t\mapsto c(-t)$, which has
  initial velocity $-v$.
  \source{petersen:page-0190}
\end{proof}

\begin{definition}[geodesically complete]
  \label{def:pet-ch5-geodesically-complete}
  \lean{PetersenLib.IsGeodesicallyComplete}
  \leanok
  \uses{def:pet-ch5-geodesic}
  A Riemannian manifold on which every geodesic is defined for all time (i.e.\ on
  all of $\mathbb R$) is called \emph{geodesically complete}.
  \source{petersen:page-0190}
\end{definition}

\begin{lemma}[Lemma 5.2.6 — uniform extension on a compact interval]
  \label{lem:pet-ch5-uniform-neighborhood}
  \lean{PetersenLib.geodesic_uniformExtensionOnCompactInterval}
  \leanok
  \uses{thm:pet-ch5-local-existence}
  Suppose $c:[a,b]\to M$ is a geodesic on a compact interval. There is a
  neighborhood $V$ in $TM$ of $\dot c(a)$ such that if $v\in V$, then there is a
  geodesic $c_v:[a,b]\to M$ with $\dot c_v(a)=v$.
  \source{petersen:page-0190,petersen:page-0191}
\end{lemma}
\begin{proof}
  \uses{thm:pet-ch5-local-existence,lem:pet-ch5-global-uniqueness}
  A compactness argument subdivides $a=b_0<b_1<\dots<b_k=b$ so that each $b_i$ has a
  neighborhood $V_i\subset TM$ such that every geodesic with initial velocity in
  $V_i$ (based at the corresponding point) is defined on $[b_i,b_{i+1}]$, by
  \cref{thm:pet-ch5-local-existence}. Using continuity of $(t,v)\mapsto c_v(t)$,
  choose $U_0\subset V_0$ a neighborhood of $\dot c(b_0)$ with $\dot c_v(b_1)\in V_1$
  for $v\in U_0$; inductively choose $U_{i}\subset U_{i-1}$ with $\dot c_v(b_{i+1})\in
  V_{i+1}$ for $v\in U_i$. After $k$ steps, $V=U_{k-1}$ is the desired neighborhood.
  \source{petersen:page-0190,petersen:page-0191}
\end{proof}

Geodesics in Euclidean space are straight lines; this yields simple examples of the
above ideas on open subsets of $\mathbb R^2$.

\begin{example}[Example 5.2.7 — jump in maximal interval]
  \label{ex:pet-ch5-punctured-plane-geodesic}
  \lean{PetersenLib.example_punctured_plane_geodesic_interval}
  \uses{def:pet-ch5-geodesically-complete}
  In the punctured plane $\mathbb R^2-\{(0,0)\}$ the unit-speed geodesic from
  $(-1,0)$ with tangent $(1,0)$ is defined only on $(-\infty,1)$, but nearby
  geodesics from $(-1,0)$ with tangent $(1+\varepsilon_1,\varepsilon_2)$,
  $\varepsilon_2\ne0$ small, are defined on all of $\mathbb R$: maximal intervals can
  jump up in size but, by \cref{lem:pet-ch5-uniform-neighborhood}, not down.
  \source{petersen:page-0191}
\end{example}

\begin{example}[Example 5.2.8 — endpoints escaping to infinity]
  \label{ex:pet-ch5-region-geodesic-jump}
  \lean{PetersenLib.example_region_geodesic_endpoint_limit}
  \uses{lem:pet-ch5-uniform-neighborhood}
  On $\{(x,y)\mid -1<xy\}$, $t\mapsto(t,0)$ is a geodesic on all of $\mathbb R$ that
  is a limit as $\varepsilon\to0$ of unit-speed geodesics $t\mapsto(t,-\varepsilon)$,
  each defined only on a finite interval whose endpoints tend to infinity, as
  required by \cref{lem:pet-ch5-uniform-neighborhood}.
  \source{petersen:page-0191}
\end{example}

\begin{example}[Example 5.2.9 — great circles on spheres]
  \label{ex:pet-ch5-sphere-great-circles}
  \lean{PetersenLib.sphereGeodesicsAreGreatCircles}
  \uses{def:pet-ch5-geodesic,prop:pet-ch5-mixed-partials-submanifold}
  On $S^n(R)\subset\mathbb R^{n+1}$, the acceleration of $c:I\to S^n(R)$ equals the
  Euclidean acceleration in $\mathbb R^{n+1}$ projected onto $S^n(R)$
  (\cref{prop:pet-ch5-mixed-partials-submanifold}), so $c$ is a geodesic iff $\ddot c$
  is normal to $S^n(R)$, i.e.\ proportional to $c$. The great circles
  $c(t)=p\cos(at)+v\sin(at)$ with $p,v\in\mathbb R^{n+1}$, $|p|=|v|=R$, $p\perp v$,
  satisfy this, with $c(0)=p$, $\dot c(0)=av$: for every initial value problem there
  is such a geodesic.
  \source{petersen:page-0191}
\end{example}
\begin{proof}
  \uses{ex:pet-ch5-sphere-great-circles}
  Differentiating twice, $\ddot c(t)=-a^2\big(p\cos(at)+v\sin(at)\big)=-a^2c(t)$,
  which is normal to $S^n(R)$ at $c(t)$ since $c(t)\cdot c(t)=R^2$ makes $c(t)$
  itself a normal vector to the sphere there. Hence the Euclidean acceleration is
  already normal to $S^n(R)$, so its tangential projection — the intrinsic
  acceleration by \cref{prop:pet-ch5-mixed-partials-submanifold} — vanishes, and $c$
  is a geodesic.
  \source{petersen:page-0191}
\end{proof}

\begin{example}[Example 5.2.10 — hyperbolas on hyperbolic space]
  \label{ex:pet-ch5-hyperbolic-hyperbolas}
  \lean{PetersenLib.hyperbolicSpaceGeodesicsAreHyperbolas}
  \uses{def:pet-ch5-geodesic,prop:pet-ch5-mixed-partials-submanifold}
  On $H^n(R)\subset\mathbb R^{n,1}$, the acceleration is again the Minkowski
  projection of the ambient (Euclidean-formula) acceleration, and
  $T_pH^n(R)=\{v\mid v\perp p\}$ in the Minkowski sense. The curves
  $c(t)=p\cosh(at)+v\sinh(at)$, $p,v\in\mathbb R^{n,1}$, $|p|^2=-R^2$, $|v|^2=R^2$,
  $p\perp v$ (Minkowski), satisfy $\ddot c=a^2c$, proportional to $c$ hence normal to
  $H^n(R)$, so these hyperbolas are the geodesics.
  \source{petersen:page-0192}
\end{example}

\begin{example}[Example 5.2.11 — left-invariant geodesics on Lie groups]
  \label{ex:pet-ch5-lie-group-geodesics}
  \lean{PetersenLib.leftInvariantGeodesicFields}
  \uses{def:pet-ch5-geodesic}
  On a Lie group $G$ with left-invariant metric, the integral curves of a
  left-invariant vector field $X$ are geodesics iff $\nabla_XX\equiv0$; the
  upper-half-plane Lie group model does not satisfy this, but for a bi-invariant
  metric and left-invariant $X$, $\nabla_XX=\tfrac12[X,X]=0$; all compact Lie groups
  admit bi-invariant metrics.
  \source{petersen:page-0192}
\end{example}

\section{The Metric Structure of a Riemannian Manifold}

\begin{definition}[piecewise $C^\infty$ curve]
  \label{def:pet-ch5-piecewise-smooth-curve}
  \lean{PetersenLib.IsPiecewiseSmoothCurve}
  \leanok
  A curve $c:[a,b]\to M$ is \emph{piecewise $C^\infty$} if it is continuous and there
  is a partition $a=a_1<a_2<\dots<a_k=b$ such that $c|_{[a_i,a_{i+1}]}$ is $C^\infty$
  for $i=1,\dots,k-1$.
  \source{petersen:page-0193}
\end{definition}

\begin{definition}[length of a curve]
  \label{def:pet-ch5-curve-length}
  \lean{PetersenLib.curveLength,PetersenLib.ContMDiffOn.intervalIntegrable_sqrt_curveSpeedSq,PetersenLib.IsPiecewiseSmoothCurve.intervalIntegrable_sqrt_curveSpeedSq,PetersenLib.curveLength_comp_mul_add,PetersenLib.curveLength_comp_const_sub}
  \leanok
  \uses{def:pet-ch5-piecewise-smooth-curve}
  For $c:[a,b]\to M$ piecewise $C^\infty$, the \emph{length} is
  \[
    L(c)=\int_a^b|\dot c(t)|\,dt=\int_a^b\sqrt{g(\dot c(t),\dot c(t))}\,dt,
  \]
  a well-defined finite nonnegative number, invariant under reparametrization by the
  chain and substitution rules. $c$ is \emph{parametrized by arc length} if
  $L(c|_{[a,t]})=t-a$ for all $t\in[a,b]$, equivalently $|\dot c(t)|=1$ for all $t$.
  \source{petersen:page-0193,petersen:page-0194}
\end{definition}

\begin{proposition}[arclength reparametrization of regular curves]
  \label{prop:pet-ch5-arclength-reparametrization}
  \lean{PetersenLib.regularCurve_arclengthReparametrization,PetersenLib.IsPiecewiseRegularCurve,PetersenLib.piecewiseRegularCurve_arclengthReparametrization,PetersenLib.curveLength_eq_of_unitSpeed_offFinset,PetersenLib.contDiffAt_curveSpeedSq,PetersenLib.curveSpeedSq_nonneg}
  \leanok
  \uses{def:pet-ch5-curve-length}
  A regular curve $c:[a,b]\to M$ (velocity never vanishing) admits a reparametrization
  $c\circ\varphi^{-1}:[0,L(c)]\to M$ by an arclength parameter, piecewise smooth with
  unit speed everywhere.
  \source{petersen:page-0194}
\end{proposition}
\begin{proof}
  \leanok
  \uses{prop:pet-ch5-arclength-reparametrization}
  Set $s=\varphi(t)=\int_a^t|\dot c(\tau)|\,d\tau$. Since $c$ is regular,
  $\varphi'(t)=|\dot c(t)|>0$ on each smooth piece, so $\varphi:[a,b]\to[0,L(c)]$ is
  strictly increasing and piecewise smooth, hence invertible with piecewise smooth
  inverse. The curve $\bar c=c\circ\varphi^{-1}$ satisfies
  $\dot{\bar c}(s)=\dot c(\varphi^{-1}(s))\cdot(\varphi^{-1})'(s)
  =\dot c(t)/\varphi'(t)$ at $t=\varphi^{-1}(s)$, so $|\dot{\bar c}(s)|=1$.
  \source{petersen:page-0194}
\end{proof}

\begin{definition}[Riemannian distance]
  \label{def:pet-ch5-distance}
  \lean{PetersenLib.riemannianDistance,PetersenLib.riemannianDistance_comm,PetersenLib.riemannianDistance_triangle,PetersenLib.riemannianDistance_triangle_of_connected,PetersenLib.exists_isPiecewiseSmoothCurve_connecting}
  \leanok
  \uses{def:pet-ch5-curve-length}
  For $p,q\in M$ let $\Omega_{p,q}=\{c:[0,1]\to M\mid c\text{ piecewise }C^\infty,\
  c(0)=p,\ c(1)=q\}$ and define the \emph{distance} $d(p,q)=|pq|=\inf\{L(c)\mid
  c\in\Omega_{p,q}\}$. Immediately $|pq|=|qp|$ and $|pq|\le|pr|+|rq|$; that $|pq|=0$
  only for $p=q$ is established in \cref{thm:pet-ch5-metric-topology}. When the
  metric needs to be specified one writes $|pq|_g$.
  \source{petersen:page-0194}
\end{definition}

\begin{definition}[metric balls and set distance]
  \label{def:pet-ch5-metric-balls}
  \lean{PetersenLib.metricBall,PetersenLib.metricClosedBall,PetersenLib.setDistance}
  \leanok
  \uses{def:pet-ch5-distance}
  $B(p,r)=\{x\in M\mid|px|<r\}$, $\bar B(p,r)=\{x\in M\mid|px|\le r\}$. For
  $A,B\subset M$, $d(A,B)=|AB|=\inf\{|pq|\mid p\in A,q\in B\}$, and
  $B(A,r)=\{x\in M\mid|Ax|<r\}$, $\bar B(A,r)=D(A,r)=\{x\in M\mid|Ax|\le r\}$.
  \source{petersen:page-0194}
\end{definition}

\begin{example}[Example 5.3.1 — infimum not realized]
  \label{ex:pet-ch5-infimum-not-realized}
  \lean{PetersenLib.example_infimum_not_realized}
  \uses{def:pet-ch5-distance}
  In $\mathbb R^2-\{(0,0)\}$ with the induced metric, $|(-1,0)(1,0)|=2$, but no curve
  of length $2$ from $(-1,0)$ to $(1,0)$ avoids $(0,0)$, so this distance is not
  realized by any curve; in the sense developed later, $\mathbb R^2-\{(0,0)\}$ is
  incomplete.
  \source{petersen:page-0194,petersen:page-0195}
\end{example}

\begin{definition}[segment]
  \label{def:pet-ch5-segment}
  \lean{PetersenLib.IsSegment,PetersenLib.segmentNotation}
  \leanok
  \uses{def:pet-ch5-distance,def:pet-ch5-curve-length}
  A curve $\sigma\in\Omega_{p,q}$ is a \emph{segment} if $L(\sigma)=|pq|$ and $\sigma$
  is parametrized proportionally to arc length. The notation $\overline{pq}$ denotes
  a specific segment on $[0,|pq|]$ with $\overline{pq}(0)=p$, $\overline{pq}(|pq|)=q$.
  \source{petersen:page-0195}
\end{definition}

\begin{lemma}[Lemma 5.3.2 (length bound) — distance functions bound length below]
  \label{lem:pet-ch5-distance-function-length-bound}
  \lean{PetersenLib.distanceFunction_curveLength_ge}
  \leanok
  \uses{def:pet-ch5-curve-length,def:pet-ch3-distance-function}
  Let $r:U\to\mathbb R$ be a smooth distance function on an open set
  $U\subset(M,g)$. For every $C^\infty$ curve $c:[a,b]\to M$ whose image lies in
  $U$,
  \[
    r(c(b))-r(c(a))\le L(c)|_a^b.
  \]
  \source{petersen:page-0195}
\end{lemma}
\begin{proof}
  \leanok
  \uses{lem:pet-ch5-distance-function-length-bound}
  By Cauchy–Schwarz and the eikonal equation $|\nabla r|=1$,
  $dr(v)=g(\nabla r,v)\le|\nabla r|\,|v|=|v|$ for every tangent vector $v$; and
  the chain rule gives the pointwise identity $(r\circ c)'(t)=dr(\dot c(t))
  =g(\nabla r,\dot c)$. Hence
  \[
    L(c)=\int_a^b|\dot c|\,dt\ge\int_a^b g(\nabla r,\dot c)\,dt
      =\int_a^b (r\circ c)'\,dt=r(c(b))-r(c(a)),
  \]
  the last step by the fundamental theorem of calculus.
  \source{petersen:page-0195}
\end{proof}

\begin{lemma}[Lemma 5.3.2 — distance functions and segments]
  \label{lem:pet-ch5-distance-function-segments}
  \lean{PetersenLib.distanceFunction_integralCurvesAreSegments,PetersenLib.distanceFunction_integralCurve_curveLength}
  \leanok
  \uses{def:pet-ch5-segment,lem:pet-ch5-distance-function-length-bound}
  If $r:U\to\mathbb R$ is a smooth distance function on open $U\subset(M,g)$, the
  integral curves of $\nabla r$ are \emph{segments in $U$}: an integral curve $c$
  of $\nabla r$ on $[a,b]$ realises $L(c)|_a^b=r(c(b))-r(c(a))=b-a$ and minimises
  length among curves in $U$ with the same endpoints.
  \source{petersen:page-0195}
\end{lemma}
\begin{proof}
  \leanok
  \uses{lem:pet-ch5-distance-function-segments}
  Along an integral curve $\dot c=\nabla r|_c$, so $|\dot c|=|\nabla r|=1$ and
  $L(c)|_a^b=\int_a^b 1\,dt=b-a$; likewise $dr(\dot c)=g(\nabla r,\nabla r)=1$, so
  $r(c(b))-r(c(a))=\int_a^b 1\,dt=b-a$, whence $L(c)|_a^b=r(c(b))-r(c(a))$. For
  any other curve $\tilde c$ in $U$ with the same endpoints,
  \cref{lem:pet-ch5-distance-function-length-bound} gives $L(\tilde c)\ge
  r(c(b))-r(c(a))=L(c)$, so $c$ is length-minimising in $U$.
  \source{petersen:page-0195}
\end{proof}

\begin{remark}[rigidity of length-minimisers]
  \label{rem:pet-ch5-distance-function-rigidity}
  \notready
  \uses{lem:pet-ch5-distance-function-segments}
  Moreover if $c\in\Omega_{p,q}(U)$ satisfies $L(c)=r(c(1))-r(c(0))$, the
  Cauchy–Schwarz equality case in $dr(\dot c)\le|\dot c|$ forces
  $\dot c=|\dot c|\,\nabla r$, so $c$ reparametrises the integral curve $\sigma$
  of $\nabla r$ through $p$ via $\varphi(s)=\int_0^s|\dot c|\,dt$:
  $c=\sigma\circ\varphi$.
  \source{petersen:page-0195}
\end{remark}

\begin{example}[Example 5.3.3 — Euclidean segments]
  \label{ex:pet-ch5-euclidean-segments}
  \lean{PetersenLib.euclideanSegmentsAreStraightLines}
  \uses{lem:pet-ch5-distance-function-segments}
  In $\mathbb R^n$, constant-speed straight segments $t\mapsto p+tv$ are segments, by
  \cref{lem:pet-ch5-distance-function-segments} applied to $r(x)=v\cdot x$ for a unit
  vector $v$. Each pair of points is joined by a segment unique up to
  reparametrization.
  \source{petersen:page-0195}
\end{example}

\begin{example}[Example 5.3.4 — circle segments]
  \label{ex:pet-ch5-circle-segment}
  \lean{PetersenLib.circleSegmentsBoundedByPi}
  \uses{lem:pet-ch5-distance-function-segments}
  For $M=S^1$, $U=S^1-\{(1,0)\}$, the distance function $r(\theta)=\theta$,
  $\theta\in(0,2\pi)$, shows any $c(\theta)=(\cos\theta,\sin\theta)$, $\theta\in
  I\subset(0,2\pi)$, is a segment in $U$; but if the length of $I$ exceeds $\pi$,
  $c$ is not a segment in $S^1$.
  \source{petersen:page-0195}
\end{example}

\begin{example}[Example 5.3.5 — segments on the round sphere]
  \label{ex:pet-ch5-sphere-segments}
  \lean{PetersenLib.sphereSegmentsCharacterization}
  \uses{lem:pet-ch5-distance-function-segments}
  On $S^n(1)\subset\mathbb R^{n+1}$ with the warped product structure
  $ds_n^2=dr^2+\sin^2(r)ds_{n-1}^2$ centered at $p$: antipodal $p,q$ correspond to
  $r=0,\pi$; otherwise $p=(a,x_0)$, $q=(b,x_0)\in(0,\pi)\times S^{n-1}$, and $r$ is
  smooth on $U\simeq(0,\pi)\times S^{n-1}$. In either case any minimizing curve has
  length $\ge\pi$ resp.\ $\ge|r(p)-r(q)|$, realized by great-circle arcs.
  \source{petersen:page-0195,petersen:page-0196}
\end{example}
\begin{proof}
  \uses{ex:pet-ch5-sphere-segments}
  If $p,q$ are antipodal, any curve $c:[0,b]\to S^n$ from $p$ to $q$ may be shortened
  to lie in $U$ except at its endpoints, with $c(0)=-c(b)$ at $r=0,\pi$. By
  \cref{lem:pet-ch5-distance-function-segments},
  $L(c|_{[\epsilon,b-\epsilon]})\ge|r(c(b-\epsilon))-r(c(\epsilon))|$, so letting
  $\epsilon\to0$, $L(c)\ge\pi$; great-circle arcs of length $\pi$ realize this. When
  $p,q$ are not antipodal they lie on a unique integral curve of $\nabla r$, part of a
  great circle in $U$, which by the lemma is the shortest curve in $U$; any curve
  leaving $U$ passes through $r=0$ or $r=\pi$ and by the same triangle-inequality
  argument as above has length $\ge\min\{r(p)+r(q),2\pi-r(p)-r(q)\}\ge|r(p)-r(q)|$.
  \source{petersen:page-0195,petersen:page-0196}
\end{proof}

\begin{example}[Example 5.3.6 — segments in hyperbolic space]
  \label{ex:pet-ch5-hyperbolic-segments}
  \lean{PetersenLib.hyperbolicSpaceAllGeodesicsAreSegments}
  \uses{ex:pet-ch5-sphere-segments}
  The same strategy as for the sphere shows that all geodesics in hyperbolic space
  are segments.
  \source{petersen:page-0196}
\end{example}

\begin{example}[Example 5.3.7 — punctured plane again]
  \label{ex:pet-ch5-punctured-plane-no-segment}
  \lean{PetersenLib.example_punctured_plane_no_segment}
  \uses{ex:pet-ch5-infimum-not-realized}
  In $\mathbb R^2-\{(0,0)\}$, as already noted, not every pair of points is joined by
  a segment.
  \source{petersen:page-0196}
\end{example}

\begin{theorem}[Thm. 5.3.8 — metric and manifold topologies agree]
  \label{thm:pet-ch5-metric-topology}
  \lean{PetersenLib.riemannianDistance_inducesManifoldTopology,PetersenLib.riemannianMetricSpace,PetersenLib.isOpen_iff_riemannianDistance,PetersenLib.eq_of_riemannianDistance_eq_zero,PetersenLib.exists_riemannianDistance_le_sqrt_mul,PetersenLib.exists_sqrt_mul_le_riemannianDistance,PetersenLib.sqrt_mul_norm_sub_le_curveLength,PetersenLib.sqrt_mul_le_curveLength_of_exit}
  \leanok
  \uses{def:pet-ch5-distance}
  The metric topology on $M$ induced by $|\cdot|$ coincides with the manifold
  topology.
  \source{petersen:page-0196}
\end{theorem}
\begin{proof}
  \uses{thm:pet-ch5-metric-topology}
  \leanok
  Fix $p\in M$ and a coordinate chart $U$ around $p$ with $x^i(p)=0$ and
  $g_{ij}(p)=\delta_{ij}$. Let $g_0(\partial_i,\partial_j)=\delta_{ij}$ be the flat metric
  in these coordinates, so $g_0=g$ at $p$; shrink $U$ so that
  $U=B^{g_0}(p,\varepsilon)=\{x\in U\mid\sqrt{\textstyle\sum_i(x^i)^2}<\varepsilon\}$.
  By continuity of $g$ there are continuous $\lambda,\mu:U\to(0,\infty)$, with
  $\lambda(x),\mu(x)\to1$ as $x\to p$, such that
  $\lambda(x)|v|_{g_0}\le|v|_g\le\mu(x)|v|_{g_0}$ for all $v\in T_xM$.

  Let $c:[0,1]\to M$ be a curve from $p$ to $x\in U$. If $c$ is a Euclidean straight
  line it lies in $U$ and, with $\mu_{\max}=\max_t\mu(c(t))$,
  \[
    |px|_{g_0}=L_{g_0}(c)\ge\frac1{\mu_{\max}}\int_0^1|\dot c|_g\,dt
    =\frac1{\mu_{\max}}L_g(c)\ge\frac1{\mu_{\max}}|px|_g.
  \]
  If $c$ is a general curve lying entirely in $U$, with $\lambda_{\min}=\min_t
  \lambda(c(t))$,
  \[
    L_g(c)=\int_0^1|\dot c|_g\,dt\ge\lambda_{\min}\int_0^1|\dot c|_{g_0}\,dt
    \ge\lambda_{\min}|px|_{g_0}.
  \]
  If $c$ leaves $U$, let $t_0$ be the first exit time; then
  $L_g(c)\ge\int_0^{t_0}|\dot c|_g\,dt\ge\lambda_{\min}\int_0^{t_0}|\dot c|_{g_0}\,dt
  \ge\lambda_{\min}\varepsilon\ge\lambda_{\min}|px|_{g_0}$.

  Shrinking $U$ further so $\min\lambda\ge\lambda_0>0$ and $\max\mu\le\mu_0<\infty$
  on $U$, taking the infimum over $c$ in each of the three cases gives
  $|px|_g\le\mu_0|px|_{g_0}$ and $\lambda_0|px|_{g_0}\le\inf L_g(c)=|px|_g$. Thus the
  Euclidean and Riemannian distances are comparable near $p$, so the metric and
  manifold topologies agree there, and $|pq|=0$ forces $p=q$. Finally
  $\lim_{x\to p}|px|_g/|px|_{g_0}=1$ since $\lambda(x),\mu(x)\to1$.
  \source{petersen:page-0196,petersen:page-0197,petersen:page-0198}
\end{proof}

\begin{corollary}[Corollary 5.3.9 — compact manifolds are metrically complete]
  \label{cor:pet-ch5-compact-complete-metric}
  \lean{PetersenLib.compactManifold_metricallyComplete}
  \leanok
  \uses{thm:pet-ch5-metric-topology}
  If $M$ is compact and $g$ a Riemannian metric on $M$, then $(M,|\cdot|_g)$ is a
  complete metric space.
  \source{petersen:page-0198}
\end{corollary}
\begin{proof}
  \uses{cor:pet-ch5-compact-complete-metric}
  \leanok
  By \cref{thm:pet-ch5-metric-topology} the metric topology equals the (compact)
  manifold topology, so $(M,|\cdot|_g)$ is a compact metric space, and compact metric
  spaces are complete.
  \source{petersen:page-0198}
\end{proof}

\begin{corollary}[Corollary 5.3.10 — approximation by constant-speed curves]
  \label{cor:pet-ch5-regular-curve-approx}
  \lean{PetersenLib.approximateByConstantSpeedCurve,PetersenLib.IsConstantSpeedCurve}
  \leanok
  \uses{thm:pet-ch5-metric-topology,prop:pet-ch5-arclength-reparametrization}
  For $c\in\Omega_{p,q}$ and $\epsilon>0$ there is a constant-speed
  $\bar c\in\Omega_{p,q}$ with $L(\bar c)\le(1+\epsilon)L(c)$.
  \source{petersen:page-0198}
\end{corollary}
\begin{proof}
  \leanok
  \uses{cor:pet-ch5-regular-curve-approx}
  It suffices to find a regular curve with the stated property, since regular curves
  admit exact-length arclength reparametrizations
  (\cref{prop:pet-ch5-arclength-reparametrization}). In a chart with Euclidean
  comparison metric $g_0$ as in the proof of \cref{thm:pet-ch5-metric-topology}, with
  $\lambda_0L_{g_0}(c)\le L_g(c)\le\mu_0L_{g_0}(c)$ for curves in the chart,
  approximate $c$ by a regular (e.g.\ polygonal, smoothed) curve $\bar c$ with
  $L_{g_0}(\bar c)\le L_{g_0}(c)$; then
  \[
    L_g(\bar c)\le\mu_0L_{g_0}(\bar c)\le\mu_0L_{g_0}(c)\le\frac{\mu_0}{\lambda_0}L_g(c).
  \]
  Shrinking the chart makes $\mu_0/\lambda_0<1+\epsilon$; compactness of $c([0,1])$
  covers it by finitely many such charts, giving the desired global regular curve
  $\bar c$.
  \source{petersen:page-0198}
\end{proof}

\begin{remark}[Remark 5.3.11 — other classes of curves]
  \label{rem:pet-ch5-absolutely-continuous-curves}
  \lean{PetersenLib.remark_absolutelyContinuousCurves}
  \uses{cor:pet-ch5-regular-curve-approx}
  The distance theory can equivalently be developed using absolutely continuous
  curves; by \cref{cor:pet-ch5-regular-curve-approx} one could also have restricted
  to piecewise smooth constant-speed curves throughout.
  \source{petersen:page-0198}
\end{remark}

\begin{definition}[functional distance]
  \label{def:pet-ch5-functional-distance}
  \lean{PetersenLib.functionalDistance,PetersenLib.HasGradientLeOne,PetersenLib.functionalDistance_nonneg}
  \leanok
  \uses{def:pet-ch5-distance}
  The \emph{functional distance} is
  \[
    d_F(p,q)=\sup\{|f(p)-f(q)|\mid f:M\to\mathbb R,\ |\nabla f|\le1\text{ on }M\}.
  \]
  Always $d_F\le d$, and $d_F$ also generates the manifold topology; once smooth
  distance functions are established the two distances agree for $p,q$ sufficiently
  close.
  \source{petersen:page-0199}
\end{definition}

\section{First Variation of Energy}

\begin{definition}[energy functional]
  \label{def:pet-ch5-energy-functional}
  \lean{PetersenLib.energyFunctional,PetersenLib.energyFunctional_nonneg,PetersenLib.IsPiecewiseSmoothCurve.intervalIntegrable_curveSpeedSq}
  \leanok
  \uses{def:pet-ch5-curve-length}
  The \emph{arclength functional} is $L(c)=\int_0^1|\dot c|\,dt$, $c\in\Omega_{p,q}$;
  its minima (if they exist) have minimal length but not necessarily the correct
  (constant-speed) parametrization, and $L$ is invariant under reparametrization. The
  \emph{energy functional}
  \[
    E(c)=\frac12\int_0^1|\dot c|^2\,dt,\qquad c\in\Omega_{p,q},
  \]
  depends on parametrization and measures kinetic energy.
  \source{petersen:page-0199}
\end{definition}

\begin{proposition}[Prop. 5.4.1 — $E$- and $L$-minimizers coincide]
  \label{prop:pet-ch5-energy-length-minimizers}
  \lean{PetersenLib.energyMinimizer_of_constantSpeed_lengthMinimizer,PetersenLib.lengthMinimizer_of_energyMinimizer,PetersenLib.curveLength_sq_le_two_mul_energyFunctional,PetersenLib.IsConstantSpeedCurve.two_mul_energyFunctional}
  \leanok
  \uses{def:pet-ch5-energy-functional,cor:pet-ch5-regular-curve-approx}
  If $\sigma\in\Omega_{p,q}$ has constant speed and minimizes $L$, then $\sigma$
  minimizes $E$. Conversely, if $\sigma$ minimizes $E$, then $\sigma$ minimizes $L$.
  \source{petersen:page-0199}
\end{proposition}
\begin{proof}
  \leanok
  \uses{prop:pet-ch5-energy-length-minimizers}
  Cauchy--Schwarz gives $L(c)=\int_0^1|\dot c|\cdot1\,dt\le\sqrt{\int_0^1|\dot c|^2dt}
  \sqrt{\int_0^11^2dt}=\sqrt{2E(c)}$, with equality iff $|\dot c|$ is constant.

  If $\sigma$ has constant speed and minimizes $L$, then for any $c\in\Omega_{p,q}$,
  $E(\sigma)=\tfrac12L(\sigma)^2\le\tfrac12L(c)^2\le E(c)$, so $\sigma$ minimizes $E$.

  Conversely let $\sigma$ minimize $E$ and $c\in\Omega_{p,q}$ arbitrary. If $c$ lacks
  constant speed, \cref{cor:pet-ch5-regular-curve-approx} gives constant-speed
  $c_\epsilon$ with $L(c_\epsilon)\le(1+\epsilon)L(c)$ for any $\epsilon>0$. Then
  \[
    L(\sigma)\le\sqrt{2E(\sigma)}\le\sqrt{2E(c_\epsilon)}=L(c_\epsilon)\le(1+\epsilon)L(c),
  \]
  and letting $\epsilon\to0$ gives $L(\sigma)\le L(c)$.
  \source{petersen:page-0199,petersen:page-0200}
\end{proof}

\begin{definition}[variation of a curve]
  \label{def:pet-ch5-variation}
  \lean{PetersenLib.CurveVariation,PetersenLib.IsProperVariation,PetersenLib.CurveVariation.curve,PetersenLib.CurveVariation.isPiecewiseSmoothCurve}
  \leanok
  A \emph{variation} of $c:I\to M$ is $\bar c:(-\varepsilon,\varepsilon)\times[a,b]\to
  M$ with $\bar c(0,t)=c(t)$ for $t\in[a,b]$; it is \emph{piecewise smooth} if
  continuous and $[a,b]$ splits into $[a_i,a_{i+1}]$, $i=0,\dots,m-1$, on which
  $\bar c$ is smooth. Write $c_s(t)=\bar c(s,t)$. The \emph{velocity field}
  $\partial\bar c/\partial t$ is well defined on each $[a_i,a_{i+1}]$, with one-sided
  values $\partial\bar c/\partial t^\pm(s,a_i)$ at the breaks. The \emph{variational
  field} is $\partial\bar c/\partial s$, well defined and continuous everywhere,
  smooth on each $(-\varepsilon,\varepsilon)\times[a_i,a_{i+1}]$. When $a=0$, $b=1$,
  $\bar c(s,0)=p$, $\bar c(s,1)=q$ for all $s$, all $c_s\in\Omega_{p,q}$ and the
  variation is called \emph{proper}.
  \source{petersen:page-0200}
\end{definition}

\begin{lemma}[Lemma 5.4.2 — first variation formula]
  \label{lem:pet-ch5-first-variation-formula}
  \lean{PetersenLib.firstVariationOfEnergy,PetersenLib.hasDerivAt_pieceEnergy,PetersenLib.hasDerivAt_windowEnergy,PetersenLib.hasDerivAt_windowEnergy_chart,PetersenLib.chartReading_acceleration_transfer,PetersenLib.firstVariation_integrand_chart,PetersenLib.hasDerivAt_halfSpeedSq_variation}
  \leanok
  \uses{def:pet-ch5-variation,def:pet-ch5-energy-functional}
  For a piecewise smooth variation $\bar c:(-\varepsilon,\varepsilon)\times[a,b]\to M$,
  \begin{align*}
    \frac{dE(c_s)}{ds}
    &=-\int_a^bg\!\left(\frac{\partial^2\bar c}{\partial t^2},\frac{\partial\bar c}{\partial s}\right)dt
    +g\!\left(\frac{\partial\bar c}{\partial t},\frac{\partial\bar c}{\partial s}\right)\Big|_{(s,b)}
    -g\!\left(\frac{\partial\bar c}{\partial t},\frac{\partial\bar c}{\partial s}\right)\Big|_{(s,a)}\\
    &\quad+\sum_{i=1}^{m-1}g\!\left(\frac{\partial\bar c}{\partial t^-}-\frac{\partial\bar c}{\partial t^+},
      \frac{\partial\bar c}{\partial s}\right)\Big|_{(s,a_i)}.
  \end{align*}
  (The Lean statement is at the base parameter $s=0$ --- the case every
  application uses; the general $s$ follows by re-basing the variation --- and
  records the boundary contribution in the equivalent per-piece form
  $\sum_i\big(g(\partial_s\bar c,\partial_t^-\bar c)|_{(0,a_{i+1})}
  -g(\partial_s\bar c,\partial_t^+\bar c)|_{(0,a_i)}\big)$, from which the
  displayed endpoint-plus-break-term shape arises by regrouping the finite
  sum.)
  \source{petersen:page-0200,petersen:page-0201}
\end{lemma}
\begin{proof}
  \leanok
  \uses{lem:pet-ch5-first-variation-formula}
  It suffices to treat smooth variations, since the piecewise case follows by
  splitting $E(c_s)=\sum_i\int_{a_i}^{a_{i+1}}|\partial\bar c/\partial t|^2dt$ and
  applying the formula on each smooth piece. For a smooth variation,
  \begin{align*}
    \frac{dE(c_s)}{ds}
    &=\frac{d}{ds}\frac12\int_a^bg\!\left(\frac{\partial\bar c}{\partial t},\frac{\partial\bar c}{\partial t}\right)dt
    =\frac12\int_a^b\frac{\partial}{\partial s}g\!\left(\frac{\partial\bar c}{\partial t},\frac{\partial\bar c}{\partial t}\right)dt\\
    &=\int_a^bg\!\left(\frac{\partial^2\bar c}{\partial s\partial t},\frac{\partial\bar c}{\partial t}\right)dt
    =\int_a^bg\!\left(\frac{\partial^2\bar c}{\partial t\partial s},\frac{\partial\bar c}{\partial t}\right)dt\\
    &=\int_a^b\frac{\partial}{\partial t}g\!\left(\frac{\partial\bar c}{\partial s},\frac{\partial\bar c}{\partial t}\right)dt
      -\int_a^bg\!\left(\frac{\partial\bar c}{\partial s},\frac{\partial^2\bar c}{\partial t^2}\right)dt\\
    &=-\int_a^bg\!\left(\frac{\partial\bar c}{\partial s},\frac{\partial^2\bar c}{\partial t^2}\right)dt
      +g\!\left(\frac{\partial\bar c}{\partial s},\frac{\partial\bar c}{\partial t}\right)\Big|_{(s,b)}
      -g\!\left(\frac{\partial\bar c}{\partial s},\frac{\partial\bar c}{\partial t}\right)\Big|_{(s,a)},
  \end{align*}
  using symmetry of mixed partials (\cref{prop:pet-ch5-mixed-partials-submanifold}
  applies in the ambient chart) and metric-compatibility of $\nabla$ to differentiate
  under the integral, then the fundamental theorem of calculus.
  \source{petersen:page-0201}
\end{proof}

\begin{theorem}[Thm. 5.4.3 — local minima of energy are geodesics]
  \label{thm:pet-ch5-energy-minima-geodesics}
  \lean{PetersenLib.energyLocalMinimum_isGeodesic}
  \leanok
  \uses{def:pet-ch5-energy-functional,def:pet-ch5-geodesic}
  If $c\in\Omega_{p,q}$ is a local minimum for $E:\Omega_{p,q}\to[0,\infty)$, then
  $c$ is a smooth geodesic.
  \source{petersen:page-0201}
\end{theorem}
\begin{proof}
  \leanok
  \uses{lem:pet-ch5-first-variation-formula,def:pet-ch5-geodesic}
  Being a local minimum forces $c$ to be a stationary point, $dE(c_s)/ds=0$, for
  every proper variation. Fix a smooth field $w$ supported in a chart window and
  deform $c$ inside the chart, $\bar c(s,t)=\varphi^{-1}(\varphi(c(t))+s\,w(t))$; for
  small $s$ this is a proper piecewise smooth variation whose variational field
  $V=\partial_s\bar c|_{s=0}$ is the chart transport of $w$ and can be made to vanish
  off the window. By \cref{lem:pet-ch5-first-variation-formula} at $s=0$,
  \[
    0=\frac{dE(c_s)}{ds}\Big|_{s=0}
    =-\int_a^bg(\ddot c,V)\,dt+\sum_{i=1}^{m-1}g\!\left(\frac{dc}{dt^-}(a_i)-\frac{dc}{dt^+}(a_i),V(a_i)\right).
  \]
  Concentrating $w$ by a bump $\lambda\ge 0$ at an interior time $\tau$ of a smooth
  piece, in the direction of the chart acceleration $\ddot c(\tau)$, makes the break
  sum vanish and leaves $0=-\int_a^b\lambda(t)\,g(\ddot c,\ddot c)\,dt$; continuity and
  positive-definiteness of $g$ force $\ddot c(\tau)=0$. Thus $\ddot c=0$ on every
  smooth piece and $c$ is a broken geodesic. Next take a bump equal to $1$ at a break
  $a_i$ in the direction of the velocity jump $\frac{dc}{dt^-}(a_i)-\frac{dc}{dt^+}(a_i)$;
  the integral term already vanishes, so
  \[
    0=\sum_{i=1}^{m-1}\left|\frac{dc}{dt^-}(a_i)-\frac{dc}{dt^+}(a_i)\right|^2,
  \]
  forcing $\frac{dc}{dt^-}(a_i)=\frac{dc}{dt^+}(a_i)$ at every break. Finally, at each
  break position and velocity agree and both one-sided pieces satisfy the geodesic ODE
  $\ddot u=-\Gamma(\dot u,\dot u)(u)$, whose right-hand side is a continuous function of
  position and velocity; hence the one-sided second derivatives agree (both equal
  $-\Gamma(\dot u,\dot u)(u(a_i))$) and the chart reading is $C^2$ across $a_i$. So the
  geodesic equation holds at the breaks as well, and $c$ is a smooth geodesic.
  \source{petersen:page-0201,petersen:page-0202,petersen:page-0203}
\end{proof}

\begin{corollary}[characterization of segments]
  \label{cor:pet-ch5-segments-are-geodesics}
  \lean{PetersenLib.segment_isGeodesic}
  \leanok
  \uses{thm:pet-ch5-energy-minima-geodesics,prop:pet-ch5-energy-length-minimizers,def:pet-ch5-segment}
  Any piecewise smooth segment is a geodesic.
  \source{petersen:page-0203}
\end{corollary}
\begin{proof}
  \leanok
  \uses{thm:pet-ch5-energy-minima-geodesics,prop:pet-ch5-energy-length-minimizers,def:pet-ch5-segment}
  A segment $\sigma$ is a constant-speed length minimizer, hence by
  \cref{prop:pet-ch5-energy-length-minimizers} an energy minimizer, hence a local
  energy minimum, hence a smooth geodesic by
  \cref{thm:pet-ch5-energy-minima-geodesics}. (The constant-speed property follows
  from the arclength-proportional parametrization by differentiating the length
  integral $\int_a^t|\dot\sigma|\,dt=k(t-a)$ at each interior time.)
  \source{petersen:page-0203}
\end{proof}

\begin{remark}[energy minima need not be minimal geodesics]
  \label{rem:pet-ch5-sphere-not-locally-minimizing}
  \lean{PetersenLib.longSphereGeodesicsNotLocallyMinimizing}
  \uses{thm:pet-ch5-energy-minima-geodesics}
  Local minima of $E$ characterize geodesics, but not every geodesic is a local
  minimum. In $\mathbb R^n$ all geodesics are integral curves of globally defined
  distance functions $u(x)=v\cdot x$; on the unit sphere, however, no geodesic of
  length $>\pi$ is locally minimizing, since it forms part of a great circle whose
  complementary arc has length $<\pi$, giving a shorter competitor.
  \source{petersen:page-0203}
\end{remark}

\section{Riemannian Coordinates}

\subsection{The Exponential Map}

\begin{definition}[exponential map; homogeneity property]
  \label{def:pet-ch5-exponential-map}
  \lean{PetersenLib.expMap,PetersenLib.geodesicHomogeneity}
  \leanok
  \uses{def:pet-ch5-geodesic}
  For $v\in T_pM$ let $c_v$ be the unique geodesic with $c_v(0)=p$, $\dot c_v(0)=v$,
  and $[0,L_v)$ the nonnegative part of its maximal domain. Uniqueness of geodesics
  gives the \emph{homogeneity property} $c_{\alpha v}(t)=c_v(\alpha t)$ for
  $\alpha>0$, $t<L_{\alpha v}$, so $L_{\alpha v}=\alpha^{-1}L_v$. Let $O_p\subset
  T_pM$ be the vectors $v$ with $1<L_v$. The \emph{exponential map} at $p$,
  $\exp_p:O_p\to M$, is $\exp_p(v)=c_v(1)$; equivalently $\exp_p(tv)=c_v(t)$, so
  $\exp_p(v)$ is reached from $p$ by going ``distance'' $|v|$ in direction $v/|v|$.
  Collecting the $\exp_p$ over $p\in M$ gives $\exp:O=\bigcup_pO_p\to M$. By
  \cref{lem:pet-ch5-uniform-neighborhood,thm:pet-ch5-local-existence}, $O$ is open in
  $TM$ and $\exp$ is smooth; likewise each $O_p\subset T_pM$ is open and $\exp_p$ is
  smooth.
  \source{petersen:page-0204}
\end{definition}

\begin{proposition}[Prop. 5.5.1 — $\exp_p$ and $E$ are local diffeomorphisms]
  \label{prop:pet-ch5-exp-diffeomorphism-properties}
  \lean{PetersenLib.expMap_localDiffeomorphism,PetersenLib.expDiagonalMap_localDiffeomorphism}
  \uses{def:pet-ch5-exponential-map}
  \begin{enumerate}
  \item For $p\in M$, $D\exp_p:T_0(T_pM)\to T_pM$ is nonsingular at the origin, so
  $\exp_p$ is a local diffeomorphism near $0\in T_pM$.
  \item Define $E:O\to M\times M$ by $E(v)=(\pi(v),\exp v)$, $\pi(v)$ the base point
  of $v$. For each $p\in M$, $DE:T_{0_p}(TM)\to T_{(p,p)}(M\times M)$ is nonsingular,
  so $E$ is a diffeomorphism from a neighborhood of the zero section of $TM$ onto a
  neighborhood of the diagonal in $M\times M$.
  \end{enumerate}
  \source{petersen:page-0204,petersen:page-0205}
\end{proposition}
\begin{proof}
  \uses{prop:pet-ch5-exp-diffeomorphism-properties}
  (1) Let $I_0:T_pM\to T_0T_pM$ be the canonical isomorphism
  $I_0(v)=\frac{d}{dt}(tv)|_{t=0}$. By the homogeneity property $c_v(t)=c_{tv}(1)$,
  \[
    D\exp_p(I_0(v))=\frac{d}{dt}\exp_p(tv)\Big|_{t=0}
    =\frac{d}{dt}c_{tv}(1)\Big|_{t=0}=\frac{d}{dt}c_v(t)\Big|_{t=0}=\dot c_v(0)=v,
  \]
  so $D\exp_p\circ I_0=\mathrm{id}_{T_pM}$, and $D\exp_p$ is nonsingular. The inverse
  function theorem gives the local diffeomorphism claim.

  (2) Identify $T_{(p,p)}(M\times M)\cong T_pM\times T_pM$ and, via $I_{0_p}$,
  $T_{(p,0_p)}(TM)\cong T_pM\times T_pM$, points of $TM$ near $(p,0_p)$ written
  $(q,v)$ with $v\in T_qM$, so $E(q,v)=(q,\exp_q(v))$. Varying $q$ alone (with $v=0$)
  is the identity in the first coordinate; varying $v$ alone (with $q=p$ fixed) is
  $\exp_p(v)$ in the second coordinate, which by part (1) has differential the
  identity at $v=0$ (via $I_{0_p}$). Hence, as a linear map $T_pM\times T_pM\to
  T_pM\times T_pM$, $DE|_{(p,0_p)}$ has block form $\begin{bmatrix}I&0\\ *&I
  \end{bmatrix}$, which is nonsingular. The inverse function theorem gives local
  diffeomorphisms of neighborhoods of $(p,0_p)$ onto neighborhoods of $(p,p)$; since
  $E$ maps the zero section of $TM$ diffeomorphically onto the diagonal (which is a
  properly embedded submanifold), these local diffeomorphisms fit together into a
  single diffeomorphism from a neighborhood of the zero section of $TM$ onto a
  neighborhood of the diagonal in $M\times M$.
  \source{petersen:page-0205,petersen:page-0206}
\end{proof}

\begin{definition}[injectivity radius]
  \label{def:pet-ch5-injectivity-radius}
  \lean{PetersenLib.injectivityRadius}
  \leanok
  \uses{prop:pet-ch5-exp-diffeomorphism-properties}
  The \emph{injectivity radius} at $p$, $\operatorname{inj}_p$, is the largest
  $\varepsilon>0$ such that $\exp_p:B(0,\varepsilon)\to M$ is defined and a
  diffeomorphism onto its image.
  \source{petersen:page-0206}
\end{definition}

\begin{remark}[Riemannian normal coordinates]
  \label{rem:pet-ch5-riemannian-normal-coordinates}
  \lean{PetersenLib.riemannianNormalCoordinates}
  \uses{def:pet-ch5-injectivity-radius}
  Identifying $T_pM$ with $\mathbb R^n$ via a linear isomorphism and using that
  $\exp_p$ is a diffeomorphism near $0$ gives \emph{exponential} or \emph{Riemannian
  normal coordinates} at $p$, unique up to the choice of identification; requiring a
  linear isometry makes them unique up to $O(n)$. Part (2) of
  \cref{prop:pet-ch5-exp-diffeomorphism-properties} says that points $q_1,q_2$
  sufficiently near $p$ are joined by a unique short geodesic $\exp_{q_1}(tv)$,
  $|v|$ small.
  \source{petersen:page-0206}
\end{remark}

\begin{corollary}[uniform injectivity radius on a compact set]
  \label{cor:pet-ch5-uniform-injectivity-radius-compact}
  \lean{PetersenLib.compactSet_uniformInjectivityRadius}
  \uses{prop:pet-ch5-exp-diffeomorphism-properties}
  Let $K\subset(M,g)$ be compact. There is $\varepsilon>0$ such that for every
  $p\in K$, $\exp_p:B(0,\varepsilon)\to M$ is defined and a diffeomorphism onto its
  image.
  \source{petersen:page-0206}
\end{corollary}
\begin{proof}
  \uses{cor:pet-ch5-uniform-injectivity-radius-compact}
  By compactness it suffices to find, for a fixed $x\in M$, $\varepsilon>0$ working
  for all $p$ in a neighborhood of $x$. Choose a neighborhood $U$ of $(x,0_x)\in TM$
  on which $E$ is a diffeomorphism onto its image (part (2) of
  \cref{prop:pet-ch5-exp-diffeomorphism-properties}); a neighborhood $x\in V\subset
  M$ and a trivialization $F:V\times T_xM\to\pi^{-1}(V)\subset TM$, linear on each
  fiber $\{p\}\times T_xM\to T_pM$; and $\delta>0$ with
  $F(V\times B(0_x,\delta))\subset U$. Continuity of the pulled-back metric on
  $\{p\}\times T_xM$ (via $F$) as $p$ ranges over $V$ gives $\varepsilon>0$ and a
  neighborhood $x\in W\subset V$ with $B(0_p,\varepsilon)\subset F(\{p\}\times
  B(0_x,\delta))$ for $p\in W$. Restricting $E$ to $B(0_p,\varepsilon)\subset T_pM$
  is then a diffeomorphism onto its image $\{p\}\times M$, i.e.\ $\exp_p$ restricted
  to $B(0,\varepsilon)$ is a diffeomorphism onto its image, for all $p\in W$.
  Compactness of $K$ finishes the proof.
  \source{petersen:page-0206}
\end{proof}

\begin{definition}[normal exponential map]
  \label{def:pet-ch5-normal-exponential-map}
  \lean{PetersenLib.normalExpMap}
  \uses{def:pet-ch5-exponential-map}
  Let $N\subset M$ be a properly embedded submanifold; its normal bundle is
  $T^\perp N=\{v\in T_pM\mid p\in N,\ v\in(T_pN)^\perp\}$, so $T_pM=T_pN\oplus
  (T_pN)^\perp$ orthogonally for $p\in N$. The \emph{normal exponential map}
  $\exp^\perp:O\cap TN^\perp\to M$ restricts $\exp$ to $TN^\perp$.
  \source{petersen:page-0206,petersen:page-0207}
\end{definition}

\begin{corollary}[Corollary 5.5.3 — tubular neighborhoods]
  \label{cor:pet-ch5-tubular-neighborhood}
  \lean{PetersenLib.tubularNeighborhoodTheorem}
  \uses{def:pet-ch5-normal-exponential-map,prop:pet-ch5-exp-diffeomorphism-properties}
  $D\exp^\perp$ is nonsingular at $0_p$ for all $p\in N$, and there is an open
  neighborhood $U$ of the zero section in $TN^\perp$ on which $\exp^\perp$ is a
  diffeomorphism onto its image $\exp^\perp(U)\subset M$, a \emph{tubular
  neighborhood} of $N$.
  \source{petersen:page-0207}
\end{corollary}
\begin{proof}
  \uses{cor:pet-ch5-tubular-neighborhood}
  As in part (2) of \cref{prop:pet-ch5-exp-diffeomorphism-properties}, identify
  $T_{(p,0_p)}(TN^\perp)\cong T_pN^\perp\oplus T_{0_p}(T_p^\perp N)\cong T_pN\oplus
  T_p^\perp N\oplus T_p^\perp N$ (the first summand from varying the base point along
  $N$, the remaining two from the fiber). The map $\Xi(v)=(\pi(v),\exp^\perp(v))$
  restricts $DE$ of \cref{prop:pet-ch5-exp-diffeomorphism-properties} to this
  subspace of $T(TM)$; since $DE|_{(p,0_p)}$ is the nonsingular block matrix
  $\begin{bmatrix}I&0\\ *&I\end{bmatrix}$ on the full $T_pM\oplus T_pM$, its
  restriction to the subspace $T_pN\oplus T_p^\perp N\subset T_pM$ (embedded
  diagonally in the domain factor and mapped isomorphically by the identity block) is
  again nonsingular. Hence $D\exp^\perp$ is nonsingular at every $0_p$, $p\in N$, and
  the inverse function theorem gives local diffeomorphisms of neighborhoods of each
  $0_p$ in $TN^\perp$ onto neighborhoods of $p$ in $M$. As $N$ is properly embedded
  these patch, exactly as in the zero-section argument of
  \cref{prop:pet-ch5-exp-diffeomorphism-properties}, to a single diffeomorphism from
  a neighborhood $U$ of the zero section of $TN^\perp$ onto its image in $M$.
  \source{petersen:page-0207}
\end{proof}

\subsection{Short Geodesics Are Segments}

\begin{theorem}[Thm. 5.5.4 — short geodesics are segments]
  \label{thm:pet-ch5-short-geodesics-segments}
  \lean{PetersenLib.shortGeodesicsAreUniqueSegments}
  \uses{def:pet-ch5-exponential-map,def:pet-ch5-segment}
  Let $(M,g)$ be Riemannian, $p\in M$, $\varepsilon>0$ such that
  $\exp_p:B(0,\varepsilon)\to U\subset M$ is a diffeomorphism onto its image $U$.
  Then $U=B(p,\varepsilon)$, and for each $v\in B(0,\varepsilon)$ the geodesic
  $\exp_p(tv)$, $t\in[0,1]$, is the unique segment of speed $|v|$ from $p$ to
  $\exp_p(v)$ in $M$.
  \source{petersen:page-0207}
\end{theorem}

\begin{lemma}[radial geodesics realize distance and are segments]
  \label{lem:pet-ch5-radial-geodesic-segment}
  \lean{PetersenLib.expMap_riemannianDistance_eq, PetersenLib.exists_expMap_isSegment}
  \leanok
  \uses{def:pet-ch5-exponential-map,def:pet-ch5-segment}
  For every $p\in M$ there is $\varepsilon>0$ such that for all $v\in T_pM$ with
  $|v|_g<\varepsilon$ one has $v\in\mathcal O_p$ and
  \[
    d(p,\exp_p v)=|v|_g=\sqrt{g_p(v,v)},
  \]
  and the radial geodesic $\gamma(t)=\exp_p(tv)$, $t\in[0,1]$, is a segment from
  $p$ to $\exp_p v$ of speed $|v|_g$. This is the metric core (the existence half)
  of \cref{thm:pet-ch5-short-geodesics-segments}; the uniqueness of the segment and
  the identity $\exp_p(B(0,\varepsilon))=B(p,\varepsilon)$ are its remaining
  clauses.
  \source{petersen:page-0207}
\end{lemma}
\begin{proof}
  \leanok
  \uses{lem:pet-ch5-gauss-lemma}
  \emph{Lower bound.} On a small normal ball, the Gauss-radius comparison shows
  that every piecewise-$C^\infty$ curve from $p$ to $\exp_p v$ has length at least
  $|v|_g$; taking the infimum over such curves gives $|v|_g\le d(p,\exp_p v)$.

  \emph{Upper bound.} The radial geodesic $\gamma(t)=\exp_p(tv)$ is a competitor
  for the same infimum. Read in the chart at $p$, the pair
  $z(t)=\bigl(\varphi_p(\gamma(t)),\,(D\exp_p)_{tv}(v)\bigr)$ of position and
  velocity solves the first-order geodesic-spray equation $z'(t)=S(z(t))$ with
  $S$ smooth; since $\gamma$ is $C^1$, this ODE bootstraps $\gamma$ to
  $C^\infty$. Hence $\gamma$ is an admissible piecewise-$C^\infty$ curve, and its
  length is $\int_0^1|\dot\gamma|_g=|v|_g$ (geodesics have constant speed
  $|\dot\gamma(0)|_g=|v|_g$), so $d(p,\exp_p v)\le|v|_g$. Combining,
  $d(p,\exp_p v)=|v|_g$.

  \emph{Segment.} The curve $\gamma$ is piecewise-$C^\infty$, its length $|v|_g$
  equals $d(p,\exp_p v)=d(\gamma(0),\gamma(1))$, and it is parametrized
  proportionally to arc length: for $t\in[0,1]$ the restriction $\gamma|_{[0,t]}$
  is, after the affine reparametrization $s\mapsto ts$, the radial geodesic of
  $tv$ (homogeneity $\exp_p(tsv)=\exp_p(s\cdot tv)$), whose length is
  $|tv|_g=t\,|v|_g$; thus $L(\gamma)|_0^t=|v|_g\,t$. So $\gamma$ is a segment.
  \source{petersen:page-0207}
\end{proof}

\begin{definition}[radial distance function in exponential coordinates]
  \label{def:pet-ch5-radial-distance-function}
  \lean{PetersenLib.radialDistanceFunction}
  \uses{thm:pet-ch5-short-geodesics-segments}
  On $U=\exp_p(B(0,\varepsilon))$ let $r(x)=|\exp_p^{-1}(x)|$, the Euclidean
  distance from the origin in $B(0,\varepsilon)\subset T_pM$, continuously extended
  by $r|_{\partial U}=\varepsilon$. In Cartesian coordinates on $T_pM$,
  $\nabla r=\partial_r=\tfrac1rx^i\partial_i$.
  \source{petersen:page-0207}
\end{definition}

\begin{lemma}[Lemma 5.5.5 — the Gauss Lemma]
  \label{lem:pet-ch5-gauss-lemma}
  \lean{PetersenLib.gaussLemma}
  \leanok
  \uses{def:pet-ch5-radial-distance-function}
  On $(U,g)$ the function $r$ has gradient $\nabla r=\partial_r$, where
  $\partial_r=D\exp_p(\partial_r)$ (the pushforward of the flat radial field).
  Equivalently, and as formalized: on a normal ball $B(0,\rho)\subset T_pM$,
  read in the chart at $p$, for all $v$ with $\|v\|<\rho$ and all $w$,
  \[
    \big\langle D\exp_p|_v(v),\,D\exp_p|_v(w)\big\rangle_{\exp_p(v)}=g_p(v,w),
  \]
  where $D\exp_p|_v(v)$ is the radial pushforward $\partial_r$; this is exactly
  $dr(w)=g(\partial_r,w)$ for all $w$, i.e. $\nabla r=\partial_r$.
  \source{petersen:page-0207}
\end{lemma}
\begin{proof}
  \leanok
  \uses{lem:pet-ch5-gauss-lemma}
  Select an orthonormal basis of $T_pM$, giving Cartesian coordinates $x^1,\dots,x^n$
  transported to $U$ via $\exp_p$, with $r^2=\sum_i(x^i)^2$,
  $\partial_r=\tfrac1rx^i\partial_i$, $dr=\tfrac1r\sum_ix^idx^i$. We must show
  $dr(v)=g(\partial_r,v)$ for all $v$, and it suffices to check this for $v=\partial_r$
  and for the rotational fields $J=-x^i\partial_j+x^j\partial_i$ ($i<j$; the angular
  field in dimension $2$), since any $v$ is a linear combination of these.

  For $v=\partial_r$: the right side is $g(\partial_r,\partial_r)=1$ since the integral
  curves of $\partial_r$ (pushed-forward radial rays) are unit-speed geodesics
  $\exp_p(tv/|v|)$; the left side $dr(\partial_r)=1$ directly from the formula for
  $dr$.

  For $v=J$: directly $dr(J)=\tfrac1r\sum_ix^i(-x^i\partial_j+x^j\partial_i)(\cdot)=0$
  by antisymmetry (the coefficient of $dx^i$ in $J$ contracted with $x^i dx^i$
  cancels against that of $dx^j$). For the right side, $J$ is a Jacobi field for the
  radial geodesic flow, $[\partial_r,J]=L_{\partial_r}J=0$ (it is the pushforward of a
  linear vector field on $T_pM$ commuting with scaling), and since $\nabla_{\partial_r}
  \partial_r=0$ (geodesic equation),
  \[
    \partial_rg(\partial_r,J)=g(\nabla_{\partial_r}\partial_r,J)+g(\partial_r,\nabla_{\partial_r}J)
    =g(\partial_r,\nabla_JJ_{\text{sym}})=g(\partial_r,\nabla_{\partial_r}\partial_r)_{{}}=0,
  \]
  more precisely $g(\partial_r,\nabla_{\partial_r}J)=g(\partial_r,\nabla_J\partial_r)=
  \tfrac12\,\partial_Jg(\partial_r,\partial_r)=0$ since $|\partial_r|\equiv1$; hence
  $g(\partial_r,J)$ is constant along each radial geodesic emanating from $p$. To show
  this constant is $0$, estimate near $p$ ($r\to0$):
  \[
    |g(\partial_r,J)|\le|\partial_r|\,|J|=|J|\le|x^i|\,|\partial_i|+|x^j|\,|\partial_j|
    \le r(x)\big(|\partial_i|+|\partial_j|\big).
  \]
  Continuity of $D\exp_p$ at $0$ keeps $\partial_i,\partial_j$ bounded near $p$, so
  $|g(\partial_r,J)|\to0$ as $r\to0$; being constant along the radial geodesic, it
  must vanish identically. Thus $dr(v)=g(\partial_r,v)$ for all $v$.
  \source{petersen:page-0207,petersen:page-0208,petersen:page-0209}
\end{proof}
\begin{proof}[Proof of \cref{thm:pet-ch5-short-geodesics-segments}]
  \uses{thm:pet-ch5-short-geodesics-segments,lem:pet-ch5-gauss-lemma,lem:pet-ch5-distance-function-segments}
  In $B(0,\varepsilon)-\{0\}$ the integral curves of $\partial_r$ are the unit-speed
  rays $s\mapsto sv/|v|$; on $U$ they are forced (via $\exp_p$) to be the unit-speed
  geodesics $s\mapsto\exp_p(sv/|v|)$. By \cref{lem:pet-ch5-gauss-lemma}, $r$ is a
  smooth distance function on $U-\{p\}$, so \cref{lem:pet-ch5-distance-function-segments}
  applies: its integral curves are segments.

  First, $U\subset B(p,\varepsilon)$ since the short geodesic from $p$ to any
  $q\in U$ has length $<\varepsilon$. To see it is the \emph{only} segment, let
  $c:[0,b]\to M$ join $p$ to $q\in U$; discard any initial portion after the last
  time $c=p$, so assume $c(t)\ne p$ for $t>0$. Let $b_0\in(0,b]$ be the first time
  $c(t_0)\notin U$ (or $b_0=b$ if none exists); $c|_{[0,b_0)}$ lies in $U-\{p\}$, and
  by \cref{lem:pet-ch5-distance-function-segments},
  \[
    L(c|_{[0,b_0)})=\int_0^{b_0}|\dot c|\,dt\ge\int_0^{b_0}dr(\dot c)\,dt=r(c(b_0)),
  \]
  using $r(p)=0$. If $c(b_0)\in\partial U$ then $L(c)\ge\varepsilon>L=r(q)$, so $c$ is
  not a segment. If $b_0=b$, $L(c)\ge r(c(b))=r(q)=L$, with equality iff $\dot c$ is
  everywhere proportional to $\nabla r$, i.e.\ $c$ reparametrizes the radial geodesic
  $\exp_p(tv)$: so the short geodesic is a segment and any competing minimizer of the
  same length is a reparametrization of it.

  Finally $B(p,\varepsilon)\subset U$: if $q\in B(p,\varepsilon)$, a curve of length
  $<\varepsilon$ joins $p,q$, and the above argument (applied with $b_0=b$, since
  such a curve cannot reach $\partial U$ without exceeding length $\varepsilon$)
  shows it lies in $U$, so $q\in U$.
  \source{petersen:page-0208}
\end{proof}

\begin{corollary}[Corollary 5.5.6 — exponential image of balls]
  \label{cor:pet-ch5-metric-ball-image}
  \lean{PetersenLib.expMap_ballImage}
  \uses{thm:pet-ch5-short-geodesics-segments}
  If $\exp_p:B(0,\varepsilon)\to B(p,\varepsilon)$ is defined and a diffeomorphism,
  then for each $\delta<\varepsilon$, $\exp_p(B(0,\delta))=B(p,\delta)$ and
  $\exp_p(\bar B(0,\delta))=\bar B(p,\delta)$.
  \source{petersen:page-0209}
\end{corollary}
\begin{proof}
  \uses{cor:pet-ch5-metric-ball-image}
  By \cref{thm:pet-ch5-short-geodesics-segments} applied with radius $\delta$
  (restricting the diffeomorphism), $\exp_p(tv)$ for $v\in B(0,\delta)$, $t\in[0,1]$,
  is the unique segment of speed $|v|<\delta$ from $p$ to $\exp_p(v)$, so
  $|p\exp_p(v)|=|v|<\delta$: this shows $\exp_p(B(0,\delta))\subset B(p,\delta)$.
  Conversely if $q\in B(p,\delta)\subset B(p,\varepsilon)=U$, the theorem gives a
  unique $v\in B(0,\varepsilon)$ with $\exp_p(v)=q$ and $|v|=|pq|<\delta$, so
  $v\in B(0,\delta)$. This proves the open case; the closed case follows by the same
  argument with non-strict inequalities, using continuity of $\exp_p$ and $r$.
  \source{petersen:page-0209}
\end{proof}

\subsection{Properties of Exponential Coordinates}

\begin{remark}[radial isometry reformulation of the Gauss Lemma]
  \label{rem:pet-ch5-radial-isometry}
  \lean{PetersenLib.radialIsometryCondition}
  \leanok
  \uses{lem:pet-ch5-gauss-lemma}
  The Gauss Lemma says $\exp_p:B(0,\varepsilon)\to B(p,\varepsilon)$ is a
  \emph{radial isometry}: $g(D\exp_p(\partial_r),D\exp_p(v))=g_p(\partial_r,v)$.
  Indeed $\tfrac1rg_{ij}v^iv^j=g(\tfrac1rx^i\partial_i,v^j\partial_j)=
  \tfrac1r\delta_{ij}x^iv^j=dr(v)$, so this is equivalent to $\nabla r=\partial_r=
  \tfrac1rx^i\partial_i$, i.e.\ to $g_{ij}v^j=\delta_{ij}v^j$ for all $v$.
  \source{petersen:page-0210}
\end{remark}

\begin{lemma}[Lemma 5.5.7 — first-order flatness of exponential coordinates]
  \label{lem:pet-ch5-normal-coords-first-order}
  \lean{PetersenLib.expCoordinates_firstOrderFlatness}
  \uses{rem:pet-ch5-radial-isometry}
  In exponential coordinates, $g_{ij}=\delta_{ij}+O(r^2)$.
  \source{petersen:page-0210}
\end{lemma}
\begin{proof}
  \uses{lem:pet-ch5-normal-coords-first-order}
  Differentiate $g_{ij}v^j=\delta_{ij}v^j$, i.e.\ $\sum_jg_{ij}x^j=\sum_j\delta_{ij}x^j$
  (valid for all $x$ in exponential coordinates), once in $x^k$:
  \[
    \delta_{ik}=\partial_k\Big(\sum_jg_{ij}x^j\Big)=(\partial_kg_{ij})x^j+g_{ik}.
  \]
  At $p$ ($x^j=0$) this gives $g_{ik}|_p=\delta_{ik}$. Differentiating again in
  $x^l$:
  \[
    0=\partial_l\big((\partial_kg_{ij})x^j\big)+\partial_lg_{ik}
    =(\partial_l\partial_kg_{ij})x^j+\partial_kg_{il}+\partial_lg_{ik},
  \]
  and evaluating at $p$: $\partial_kg_{il}|_p+\partial_lg_{ik}|_p=0$. Applying this
  three times with indices permuted,
  \[
    2\partial_kg_{ij}|_p=(\partial_kg_{ij}+\partial_jg_{ik})+(\partial_kg_{ij}+\partial_ig_{jk})
    -(\partial_ig_{jk}+\partial_jg_{ik})\Big|_p=0,
  \]
  so $\partial_kg_{ij}|_p=0$ for all $i,j,k$. With $g_{ij}|_p=\delta_{ij}$ and
  $\partial_kg_{ij}|_p=0$, Taylor's theorem gives $g_{ij}=\delta_{ij}+O(r^2)$.
  \source{petersen:page-0210,petersen:page-0211}
\end{proof}

\begin{remark}[asymptotics of the Hessian of $r$]
  \label{rem:pet-ch5-hessian-r-asymptotics}
  \lean{PetersenLib.hessianRadialFunction_asymptotics}
  \uses{lem:pet-ch5-normal-coords-first-order}
  Writing the metric in polar form $g=dr^2+g_r$ (restriction of $g$ to the level
  spheres of $r$) near $p$, and comparing with the Euclidean
  $\delta_{ij}=dr^2+r^2ds_{n-1}^2$, first-order agreement
  (\cref{lem:pet-ch5-normal-coords-first-order}) gives $\lim_{r\to0}g_r=0$ and
  $\lim_{r\to0}\big(\partial_rg_r-\partial_r(r^2ds_{n-1}^2)\big)=0$. Since
  $\partial_rg_r=2\operatorname{Hess}r$, this yields
  $\lim_{r\to0}\big(\operatorname{Hess}r-\tfrac1rg_r\big)=0$.
  \source{petersen:page-0211}
\end{remark}

\begin{theorem}[Thm. 5.5.8 (Riemann, 1854) — local rigidity of constant curvature]
  \label{thm:pet-ch5-riemann-constant-curvature-local}
  \lean{PetersenLib.constantCurvature_locallyIsometricToSpaceForm}
  \uses{rem:pet-ch5-hessian-r-asymptotics}
  If a Riemannian $n$-manifold $(M,g)$ has constant sectional curvature $k$, every
  point of $M$ has a neighborhood isometric to an open subset of the space form
  $S^n_k$.
  \source{petersen:page-0211}
\end{theorem}
\begin{proof}
  \uses{thm:pet-ch5-riemann-constant-curvature-local}
  Work in exponential coordinates around $p$. The constant-curvature assumption
  makes the radial curvature equation read
  $\nabla_{\partial_r}\operatorname{Hess}r+(\operatorname{Hess}r)^2=-kg_r$, with
  $\lim_{r\to0}\operatorname{Hess}r=0$ by \cref{rem:pet-ch5-hessian-r-asymptotics}.
  Regarding $g_r$ as given, this ODE (in $r$, along each radial geodesic) has a
  unique solution for $\operatorname{Hess}r$; but $\tfrac{\mathrm{sn}_k'(r)}{\mathrm{sn}_k(r)}g_r$
  also solves it, since (using $\nabla_{\partial_r}g_r+\nabla_{\partial_r}dr^2=0$
  and the derivative identities for $\mathrm{sn}_k$):
  \begin{align*}
    \nabla_{\partial_r}\!\left(\frac{\mathrm{sn}_k'}{\mathrm{sn}_k}g_r\right)
    +\left(\frac{\mathrm{sn}_k'}{\mathrm{sn}_k}\right)^2g_r
    &=\left(\frac{\mathrm{sn}_k'}{\mathrm{sn}_k}\right)'g_r
    +\left(\frac{\mathrm{sn}_k'}{\mathrm{sn}_k}\right)^2g_r
    +\frac{\mathrm{sn}_k'}{\mathrm{sn}_k}\nabla_{\partial_r}g_r\\
    &=-kg_r-\frac{\mathrm{sn}_k'}{\mathrm{sn}_k}\nabla_{\partial_r}dr^2=-kg_r,
  \end{align*}
  since $\nabla_{\partial_r}dr^2$ vanishes on radial and angular directions alike (a
  direct computation with $\operatorname{Hess}r(\partial_r,\cdot)=0$). Uniqueness of
  the solution with the stated initial behavior forces
  $\operatorname{Hess}r=\tfrac{\mathrm{sn}_k'(r)}{\mathrm{sn}_k(r)}g_r$. Combined with
  $g=dr^2+g_r$, the function
  \[
    f_k(r)=\begin{cases}\tfrac12r^2,&k=0\\ \tfrac1k-\tfrac1k\mathrm{cs}_k(r),&k\ne0\end{cases}
  \]
  satisfies $\operatorname{Hess}f_k=(1-kf_k)g$ (direct computation from
  $f_k'=\mathrm{sn}_k$, $f_k''=\mathrm{cs}_k=1-kf_k$ at $r=0$-normalization, propagated
  by the Hessian formula above). The classification of such metrics (space forms) then
  identifies a neighborhood of $p$ isometrically with an open subset of $S^n_k$.
  \source{petersen:page-0211,petersen:page-0212}
\end{proof}

\begin{remark}[Remark 5.5.9 — on the uniqueness used above]
  \label{rem:pet-ch5-uniqueness-remark-5-5-9}
  \lean{PetersenLib.remark_hessianODE_uniquenessSubtlety}
  \uses{thm:pet-ch5-riemann-constant-curvature-local}
  Neither of the systems $L_{\partial_r}g_r=2\operatorname{Hess}r$,
  $\nabla_{\partial_r}\operatorname{Hess}r+(\operatorname{Hess}r)^2=-kg_r$, nor the
  variant with $L_{\partial_r}\operatorname{Hess}r-(\operatorname{Hess}r)^2=-kg_r$,
  has unique solutions with $g_r,\operatorname{Hess}r\to0$ at $r=0$ in general — the
  trivial solution $g_r\equiv0$, $\operatorname{Hess}r\equiv0$ always solves them.
  Nor is it clear a priori that $\operatorname{Hess}r=\tfrac{\mathrm{sn}_k'}{\mathrm{sn}_k}g_r$
  solves the second system without knowing in advance that $L_{\partial_r}g_r=
  2\tfrac{\mathrm{sn}_k'}{\mathrm{sn}_k}g_r$; and the initial value problem
  $L_{\partial_r}g_r=2\tfrac{\mathrm{sn}_k'}{\mathrm{sn}_k}g_r$, $\lim_{r\to0}g_r=0$ has
  infinitely many solutions $\lambda\,\mathrm{sn}_k^2(r)\,ds_{n-1}^2$, $\lambda\in
  \mathbb R$, of which only one gives a smooth metric at $p$.
  \source{petersen:page-0212,petersen:page-0213}
\end{remark}

\section{Riemannian Isometries}

\subsection{Local Isometries}

\begin{definition}[local Riemannian isometry]
  \label{def:pet-ch5-local-isometry}
  \lean{PetersenLib.IsLocalRiemannianIsometry}
  \leanok
  A map $F:(M,g_M)\to(N,g_N)$ is a \emph{local Riemannian isometry} if for each
  $p\in M$, $DF_p:T_pM\to T_{F(p)}N$ is a linear isometry. A local coordinate
  system $\varphi:U\to\Omega\subset\mathbb R^n$, with $U$ given the induced metric
  $g$ and $\Omega$ the coordinate representation $(\varphi^{-1})^*g=g_{ij}dx^idx^j$,
  is a trivial example.
  \source{petersen:page-0213}
\end{definition}

\begin{proposition}[Prop. 5.6.1 — basic properties of local isometries]
  \label{prop:pet-ch5-local-isometry-properties}
  \lean{PetersenLib.localIsometry_mapsGeodesicsToGeodesics,PetersenLib.localIsometry_expNaturality,PetersenLib.localIsometry_distanceDecreasing,PetersenLib.localIsometry_distancePreserving}
  \uses{def:pet-ch5-local-isometry,def:pet-ch5-geodesic,def:pet-ch5-exponential-map,def:pet-ch5-distance}
  Let $F:(M,g_M)\to(N,g_N)$ be a local Riemannian isometry.
  \begin{enumerate}
  \item[(1)] $F$ maps geodesics to geodesics.
  \item[(2)] $F\circ\exp_p(v)=\exp_{F(p)}\circ DF_p(v)$ whenever $\exp_p(v)$ is
  defined.
  \item[(3)] $F$ is distance decreasing.
  \item[(4)] If $F$ is also a bijection, $F$ is distance preserving.
  \end{enumerate}
  \source{petersen:page-0213}
\end{proposition}
\begin{proof}
  \uses{prop:pet-ch5-local-isometry-properties}
  (1) The coordinate geodesic equation depends only on the metric and its first
  derivatives; a local isometry pulls back $g_N$ to $g_M$ and is a local
  diffeomorphism, so in matched coordinates the geodesic equations for $M$ and $N$
  coincide, and $F$ carries geodesics to geodesics.

  (2) If $\exp_p(v)$ is defined, $t\mapsto\exp_p(tv)$ is a geodesic, so by (1)
  $t\mapsto F(\exp_p(tv))$ is a geodesic in $N$, with initial velocity
  $\frac{d}{dt}F(\exp_p(tv))|_{t=0}=DF(v)$ by the chain rule. By uniqueness of
  geodesics (\cref{thm:pet-ch5-local-uniqueness}), $F(\exp_p(tv))=\exp_{F(p)}(tDF(v))$
  for all valid $t$; setting $t=1$ gives the claim.

  (3) $F$ preserves lengths of curves since $DF$ preserves the norms of tangent
  vectors, so $L_N(F\circ c)=L_M(c)$ for any curve $c$; taking the infimum over
  $c\in\Omega_{p,q}$ and using that $F\circ c\in\Omega_{F(p),F(q)}$ (a subset of
  competitors) gives $|F(p)F(q)|_N\le|pq|_M$.

  (4) Both $F$ and $F^{-1}$ are distance decreasing by (3), so
  $|F(p)F(q)|_N\le|pq|_M$ and $|pq|_M\le|F(p)F(q)|_N$, giving equality.
  \source{petersen:page-0213}
\end{proof}

\begin{proposition}[Prop. 5.6.2 — uniqueness of Riemannian isometries]
  \label{prop:pet-ch5-isometry-uniqueness}
  \lean{PetersenLib.localIsometry_uniquelyDeterminedByOnePoint}
  \uses{prop:pet-ch5-local-isometry-properties}
  Let $F,G:(M,g_M)\to(N,g_N)$ be local Riemannian isometries, $M$ connected. If
  $F(p)=G(p)$ and $DF_p=DG_p$ for some $p\in M$, then $F=G$ on $M$.
  \source{petersen:page-0214}
\end{proposition}
\begin{proof}
  \uses{prop:pet-ch5-isometry-uniqueness}
  Let $A=\{x\in M\mid F(x)=G(x),\ DF_x=DG_x\}$. Then $p\in A$ and $A$ is closed by
  continuity. If $x\in A$, part (2) of \cref{prop:pet-ch5-local-isometry-properties}
  gives, for $v\in T_xM$ with $\exp_x(v)$ defined,
  \[
    F\circ\exp_x(v)=\exp_{F(x)}\circ DF_x(v)=\exp_{G(x)}\circ DG_x(v)=G\circ\exp_x(v),
  \]
  and since $\exp_x$ maps onto a neighborhood of $x$, a neighborhood of $x$ lies in
  $A$: $A$ is open. By connectedness $A=M$.
  \source{petersen:page-0214}
\end{proof}

\begin{proposition}[Prop. 5.6.3 — completeness is preserved by Riemannian coverings]
  \label{prop:pet-ch5-covering-completeness}
  \lean{PetersenLib.riemannianCovering_completenessEquivalence}
  \uses{def:pet-ch5-local-isometry,def:pet-ch5-geodesically-complete}
  Let $F:(M,g_M)\to(N,g_N)$ be a Riemannian covering map. $(M,g_M)$ is geodesically
  complete iff $(N,g_N)$ is.
  \source{petersen:page-0214}
\end{proposition}
\begin{proof}
  \uses{prop:pet-ch5-covering-completeness}
  Let $c:(-\varepsilon,\varepsilon)\to N$ be a geodesic, $c(0)=p$, $\dot c(0)=v$. For
  $\tilde p\in F^{-1}(p)$ there is a unique lift $\tilde c:(-\varepsilon,\varepsilon)
  \to M$ with $F\circ\tilde c=c$, $\tilde c(0)=\tilde p$; since $F$ is a local
  isometry with locally-defined isometric inverse, $\tilde c$ is also a geodesic. If
  $N$ is geodesically complete, $c$ (and hence $\tilde c$) extends to all of
  $\mathbb R$ for every choice of $p,v$; as every geodesic of $M$ arises this way,
  $M$ is geodesically complete. Conversely if $M$ is complete, any $\tilde c$
  extends to all time, and $F\circ\tilde c$ is a geodesic in $N$ defined for all
  time extending $c$; so $N$ is complete.
  \source{petersen:page-0214}
\end{proof}

\begin{lemma}[Lemma 5.6.4 — completeness forces a covering map]
  \label{lem:pet-ch5-local-isometry-covering}
  \lean{PetersenLib.completeLocalIsometry_isCoveringMap}
  \uses{def:pet-ch5-local-isometry,def:pet-ch5-geodesically-complete,cor:pet-ch5-uniform-injectivity-radius-compact}
  If $F:(M,g_M)\to(N,g_N)$ is a local Riemannian isometry, $M$ geodesically
  complete, $N$ connected, then $F$ is a Riemannian covering map.
  \source{petersen:page-0214}
\end{lemma}
\begin{proof}
  \uses{lem:pet-ch5-local-isometry-covering,prop:pet-ch5-local-isometry-properties}
  Fix $q\in N$ with $\exp_q:B(0,\varepsilon)\to B(q,\varepsilon)$ a diffeomorphism.
  We claim $F^{-1}(B(q,\varepsilon))$ is evenly covered by the $B(p,\varepsilon)$,
  $F(p)=q$. Completeness of $M$ gives $\exp_p:B(0,\varepsilon)\to B(p,\varepsilon)$
  defined, and by part (2) of \cref{prop:pet-ch5-local-isometry-properties},
  $F\circ\exp_p(v)=\exp_q(DF_p(v))$ for $v\in B(0,\varepsilon)$; as $\exp_q$ and
  $DF_p$ (an isometry) are diffeomorphisms on this ball, $F\circ\exp_p:B(0,\varepsilon)
  \to B(q,\varepsilon)$ is a diffeomorphism, hence so are $\exp_p:B(0,\varepsilon)\to
  B(p,\varepsilon)$ and $F:B(p,\varepsilon)\to B(q,\varepsilon)$ individually.

  Next, $F^{-1}(B(q,\varepsilon))=\bigcup_{F(p)=q}B(p,\varepsilon)$: given
  $x\in F^{-1}(B(q,\varepsilon))$, join $q$ to $F(x)$ by the unique geodesic
  $c(t)=\exp_q(tv)$, $v\in B(0,\varepsilon)$; lift to a geodesic $\sigma:[0,1]\to M$
  with $\sigma(1)=x$, $DF_x(\dot\sigma(1))=\dot c(1)$ (using completeness to run the
  lift backward from $x$). Since $F\circ\sigma$ is a geodesic agreeing with $c$ at
  $t=1$ in position and velocity, $F\circ\sigma=c$ by
  \cref{thm:pet-ch5-local-uniqueness}, so $F(\sigma(0))=c(0)=q$, and
  $x\in B(\sigma(0),\varepsilon)$.

  Finally $F$ is surjective: $F(M)$ is open (as $F$ is a local diffeomorphism), and
  the above shows it is closed — for $q_i\in F(M)\to q$, use
  \cref{cor:pet-ch5-uniform-injectivity-radius-compact} to find $\varepsilon>0$
  working simultaneously for $x\in\{q,q_1,q_2,\dots\}$; for $k$ large, $q\in
  B(q_k,\varepsilon)\subset F(M)$. As $N$ is connected, $F(M)=N$.
  \source{petersen:page-0214,petersen:page-0215}
\end{proof}

\begin{definition}[fixed-point set; totally geodesic submanifold]
  \label{def:pet-ch5-fixed-point-set}
  \lean{PetersenLib.fixedPointSet,PetersenLib.IsTotallyGeodesic}
  \uses{def:pet-ch5-exponential-map}
  For a set $S$ of isometries of $(M,g)$, $\operatorname{Fix}(S)=\{x\in M\mid
  F(x)=x\ \forall F\in S\}$. A submanifold $N\subset(M,g)$ is \emph{totally
  geodesic} if for each $p\in N$ a neighborhood of $0\in T_pN\subset T_pM$ is
  mapped into $N$ by $\exp_p$ (for $M$): geodesics of $N$ are geodesics of $M$, and
  any $M$-geodesic tangent to $N$ at a point stays in $N$ for a short time.
  \source{petersen:page-0215}
\end{definition}

\begin{proposition}[Prop. 5.6.5 — fixed-point sets are totally geodesic]
  \label{prop:pet-ch5-fixed-point-totally-geodesic}
  \lean{PetersenLib.fixedPointSetComponent_totallyGeodesic}
  \uses{def:pet-ch5-fixed-point-set,cor:pet-ch5-metric-ball-image}
  If $S\subset\operatorname{Iso}(M,g)$, each connected component of
  $\operatorname{Fix}(S)$ is a totally geodesic submanifold.
  \source{petersen:page-0215}
\end{proposition}
\begin{proof}
  \uses{prop:pet-ch5-fixed-point-totally-geodesic}
  Let $p\in\operatorname{Fix}(S)$ and $V\subset T_pM$ the subspace fixed by every
  $DF_p:T_pM\to T_pM$, $F\in S$ (each $F\in S$ fixes $p$, so $DF_p$ acts on
  $T_pM$). If $v\in V$, the geodesic $t\mapsto\exp_p(tv)$ is fixed pointwise by
  every $F\in S$, since $F$ preserves the initial position $p$ and velocity $v$; so
  $\exp_p(tv)\in\operatorname{Fix}(S)$ whenever defined, giving $\exp_p:V\to
  \operatorname{Fix}(S)$.

  Let $\varepsilon>0$ with $\exp_p:B(0,\varepsilon)\to B(p,\varepsilon)$ a
  diffeomorphism (\cref{def:pet-ch5-injectivity-radius}). If $q\in\operatorname{Fix}(S)
  \cap B(p,\varepsilon)$, the unique geodesic $c:[0,1]\to B(p,\varepsilon)$ from $p$
  to $q$ has both endpoints fixed by each $F\in S$; then $F\circ c$ is a geodesic
  from $p$ to $q$ of the same length lying in $B(p,\varepsilon)$
  (\cref{cor:pet-ch5-metric-ball-image}), so by uniqueness $F\circ c=c$, and hence
  $c$ lies in $\operatorname{Fix}(S)\cap B(p,\varepsilon)$. This shows
  $\exp_p:V\cap B(0,\varepsilon)\to\operatorname{Fix}(S)\cap B(p,\varepsilon)$ is a
  bijection, exhibiting $\operatorname{Fix}(S)$ near $p$ as the image under $\exp_p$
  of a linear subspace, i.e.\ totally geodesic.
  \source{petersen:page-0215,petersen:page-0216}
\end{proof}

\begin{remark}[Remark 5.6.6 — even codimension of orientation-preserving fixed sets]
  \label{rem:pet-ch5-fixed-point-codimension}
  \lean{PetersenLib.orientationPreservingFixedSet_evenCodimension}
  \uses{prop:pet-ch5-fixed-point-totally-geodesic}
  If $DF_p:T_pM\to T_pM$ is orientation preserving for $p\in\operatorname{Fix}(F)$,
  the Zariski tangent space at $p$ has even codimension, since the $+1$-eigenspace
  of an element of $SO(n)$ has even codimension; so each component of
  $\operatorname{Fix}(F)$ has even codimension.
  \source{petersen:page-0216}
\end{remark}

\subsection{Constant Curvature Revisited}

\begin{theorem}[Thm. 5.6.7 — the monodromy map]
  \label{thm:pet-ch5-monodromy-theorem}
  \lean{PetersenLib.monodromyMap}
  \uses{thm:pet-ch5-riemann-constant-curvature-local,prop:pet-ch5-isometry-uniqueness}
  Let $(M,g)$ be simply connected of dimension $n$ and constant curvature $k$, and
  $L:T_pM\to T_qS^n_k$ a linear isometry. There is a unique local Riemannian
  isometry $F:M\to S^n_k$ (the \emph{monodromy map}) with $DF_p=L$. If $(M,g)$ is
  geodesically complete, $F$ is a diffeomorphism.
  \source{petersen:page-0216}
\end{theorem}
\begin{proof}
  \uses{thm:pet-ch5-monodromy-theorem,lem:pet-ch5-local-isometry-covering,prop:pet-ch5-isometry-uniqueness}
  By \cref{thm:pet-ch5-riemann-constant-curvature-local}, every $x\in M$ has a small
  ball $B(x,r)$ isometric to a ball $B(\bar x,r)\subset S^n_k$; composing with an
  isometry of $S^n_k$, for $q\in B(x,r)$, $\bar q\in S^n_k$, and a linear isometry
  $L:T_qU\to T_{\bar q}S^n_k$, there is a unique isometric embedding
  $F:B(x,r)\to S^n_k$ with $F(q)=\bar q$, $DF|_q=L$. Choose $r$ small enough that
  metric balls of that radius in $S^n_k$ are convex (automatic if $k\le0$; if
  $k>0$, radius $<\pi/\sqrt k$), so that in $M$ small metric balls are either
  disjoint or have connected intersection.

  Fix $p\in M$, $\bar p\in S^n_k$, linear isometry $L:T_pM\to T_{\bar p}S^n_k$. For
  $x\in M$ and $c\in\Omega_{p,x}$, cover $c$ by balls $B(p_i,r)$, $p_0=p$,
  $p_m=x$, consecutive balls meeting in a connected intersection. Build
  $F_0:B(p_0,r)\to S^n_k$ with $F_0(p)=\bar p$, $DF_0|_p=L$, then successively
  $F_i:B(p_i,r)\to S^n_k$ agreeing with $F_{i-1}$ on the (connected) overlap — such
  agreement is determined by matching value and differential at one point, by
  \cref{prop:pet-ch5-isometry-uniqueness}. Set $G(c)=F_m(x)$.

  $G$ is independent of the choice of covering: for two coverings, the points of
  $[0,1]$ (parametrizing $c$) where the two constructions agree form a nonempty
  (contains $0$), open, and closed subset (by the local uniqueness argument as in
  \cref{prop:pet-ch5-isometry-uniqueness}), hence all of $[0,1]$. If
  $\bar c\in\Omega_{p,x}$ is sufficiently close to $c$ it lies inside the same
  covering, so $G(c)=G(\bar c)$: $G$ is locally constant on $\Omega_{p,x}$, hence
  constant on homotopy classes, and since $M$ is simply connected, constant on all
  of $\Omega_{p,x}$. Thus $F(x):=G(c)$ is well defined independent of $c$, and $F$
  is a local Riemannian isometry with $DF_p=L$.

  If $M$ is geodesically complete, \cref{lem:pet-ch5-local-isometry-covering} shows
  $F$ is a covering map; as $S^n_k$ is simply connected, $F$ is a diffeomorphism.
  \source{petersen:page-0216,petersen:page-0217,petersen:page-0218}
\end{proof}

\begin{example}[Example 5.6.8 — immersions realize monodromy maps]
  \label{ex:pet-ch5-immersion-monodromy}
  \lean{PetersenLib.example_immersion_monodromy}
  \uses{thm:pet-ch5-monodromy-theorem}
  For an immersion $M^n\to S^n_k$ with pulled-back metric, the monodromy map $F$ of
  \cref{thm:pet-ch5-monodromy-theorem} is (a version of) this immersion; such maps
  can fold wildly for $n\ge2$ and need not resemble covering maps.
  \source{petersen:page-0216}
\end{example}

\begin{example}[Example 5.6.9 — a diffeomorphism of a punctured space form]
  \label{ex:pet-ch5-folding-immersion}
  \lean{PetersenLib.example_puncturedSpaceForm_diffeomorphism}
  \uses{thm:pet-ch5-monodromy-theorem}
  If $U\subset S^n_k$ is an open disc with smooth boundary, there is a
  diffeomorphism $F:M=S^n_k-\{p\}\to S^n_k-U$; the pulled-back metric on $M$ has
  constant curvature but looks poor near the missing point.
  \source{petersen:page-0217}
\end{example}

\begin{example}[Example 5.6.10 — $\mathbb{RP}^n$ admits no immersion into $S^n$]
  \label{ex:pet-ch5-rp-n-no-immersion}
  \lean{PetersenLib.example_projectiveSpace_noImmersionIntoSphere}
  \uses{thm:pet-ch5-monodromy-theorem}
  $M=\mathbb{RP}^n=(\mathbb R^n-\{0\})/(\text{antipodal map})$ is not simply
  connected, hence admits no immersion into $S^n$ realizing the monodromy
  construction.
  \source{petersen:page-0217}
\end{example}

\begin{example}[Example 5.6.11 — monodromy map of a universal cover]
  \label{ex:pet-ch5-universal-cover-monodromy}
  \lean{PetersenLib.example_universalCover_monodromy}
  \uses{thm:pet-ch5-monodromy-theorem}
  If $M$ is the universal cover of $S^2-\{\pm p\}$, the monodromy map is not
  one-to-one; it is exactly the covering map $M\to S^2-\{\pm p\}$.
  \source{petersen:page-0217}
\end{example}

\begin{corollary}[Corollary 5.6.12 — closed simply connected constant curvature]
  \label{cor:pet-ch5-closed-simply-connected-constant-curvature}
  \lean{PetersenLib.closedSimplyConnectedConstantCurvature_isSphere}
  \uses{thm:pet-ch5-monodromy-theorem}
  If $M$ is closed and simply connected with constant curvature $k$, then $k>0$
  and $M=S^n$; in particular $S^p\times S^q$ and $\mathbb{CP}^n$ admit no
  constant-curvature metric.
  \source{petersen:page-0217}
\end{corollary}
\begin{proof}
  \uses{cor:pet-ch5-closed-simply-connected-constant-curvature}
  $M$ closed is compact, hence geodesically complete
  (\cref{cor:pet-ch5-compact-manifold-complete}), so by
  \cref{thm:pet-ch5-monodromy-theorem} the monodromy map $F:M\to S^n_k$ is a
  diffeomorphism. As $M$ is compact, $S^n_k$ must be compact, which forces $k>0$
  (for $k\le0$, $S^n_k$ is $\mathbb R^n$ or $H^n$, noncompact); then $S^n_k=S^n$ up
  to rescaling, so $M\cong S^n$.
  \source{petersen:page-0217}
\end{proof}

\begin{corollary}[Corollary 5.6.13 — complete noncompact constant curvature]
  \label{cor:pet-ch5-noncompact-constant-curvature}
  \lean{PetersenLib.completeNoncompactConstantCurvature_universalCoverIsRn}
  \uses{thm:pet-ch5-monodromy-theorem}
  If $M$ is geodesically complete, noncompact, with constant curvature $k$, then
  $k\le0$ and the universal cover of $M$ is diffeomorphic to $\mathbb R^n$; in
  particular $S^2\times\mathbb R^2$ and $S^n\times\mathbb R$ admit no geodesically
  complete constant-curvature metric.
  \source{petersen:page-0217}
\end{corollary}
\begin{proof}
  \uses{cor:pet-ch5-noncompact-constant-curvature}
  The universal cover $\tilde M$ is simply connected and, by
  \cref{prop:pet-ch5-covering-completeness}, geodesically complete with constant
  curvature $k$; by \cref{thm:pet-ch5-monodromy-theorem}, the monodromy map
  $\tilde M\to S^n_k$ is a diffeomorphism. If $k>0$, $S^n_k$ is compact, forcing
  $\tilde M$ (hence $M$, a quotient of a compact space by a covering group action)
  compact, contrary to hypothesis. So $k\le0$ and $S^n_k$ is $\mathbb R^n$ or $H^n$,
  both diffeomorphic to $\mathbb R^n$.
  \source{petersen:page-0217}
\end{proof}

\begin{corollary}[Corollary 5.6.14 (Killing, 1893; Hopf, 1926) — classification of complete constant curvature spaces]
  \label{cor:pet-ch5-classification-constant-curvature}
  \lean{PetersenLib.killingHopfClassification}
  \uses{thm:pet-ch5-monodromy-theorem,prop:pet-ch5-covering-completeness}
  If $(M,g)$ is connected, geodesically complete, with constant curvature $k$, its
  universal cover is isometric to $S^n_k$.
  \source{petersen:page-0218}
\end{corollary}
\begin{proof}
  \uses{cor:pet-ch5-classification-constant-curvature}
  The universal cover $\tilde M$ is simply connected and, by
  \cref{prop:pet-ch5-covering-completeness}, geodesically complete with the same
  constant curvature $k$. By \cref{thm:pet-ch5-monodromy-theorem} the monodromy map
  $\tilde M\to S^n_k$ (built from any base point and any linear isometry of tangent
  spaces there) is a local isometry that is also a diffeomorphism, hence a global
  isometry.
  \source{petersen:page-0218}
\end{proof}

\subsection{Metric Characterization of Maps}

\begin{remark}[recovering the metric from the distance function]
  \label{rem:pet-ch5-polarization-metric-recovery}
  \lean{PetersenLib.metricRecoveredFromDistance}
  \uses{def:pet-ch5-distance}
  By polarization, $g(v,w)=\tfrac12(|v+w|^2-|v|^2-|w|^2)$, so it suffices to
  recover $|v|$ from $|\cdot|_g$; for a curve $\alpha$ with $\dot\alpha(0)=v$,
  $|v|=\lim_{t\to0}|\alpha(t)\alpha(0)|/t$. Thus $g$ is determined by
  $(M,|\cdot|_g)$ together with the differentiable structure, motivating synthetic
  (metric-space-only) characterizations of Riemannian maps.
  \source{petersen:page-0218,petersen:page-0219}
\end{remark}

\begin{definition}[distance coordinates]
  \label{def:pet-ch5-distance-coordinates}
  \lean{PetersenLib.distanceCoordinates}
  \uses{cor:pet-ch5-uniform-injectivity-radius-compact}
  Fix $p\in M$ and $U\ni p$ with $B(q,\operatorname{inj}(q))\supset U$ for all
  $q\in U$; then $r_q(x)=|xq|$ is smooth on $U-\{q\}$ for each $q\in U$. Choosing
  $q_1,\dots,q_n\in U-\{p\}$, $n=\dim M$, with $\nabla r_{q_1}(p),\dots,\nabla
  r_{q_n}(p)$ linearly independent, the inverse function theorem makes
  $\varphi=(r_{q_1},\dots,r_{q_n})$ coordinates near $p$.
  \source{petersen:page-0219}
\end{definition}

\begin{theorem}[Thm. 5.6.15 (Myers and Steenrod, 1939) — distance-preserving bijections are isometries]
  \label{thm:pet-ch5-myers-steenrod-isometry}
  \lean{PetersenLib.distancePreservingBijection_isRiemannianIsometry}
  \uses{def:pet-ch5-distance-coordinates,def:pet-ch5-local-isometry}
  If $F:M\to N$ is a bijection between Riemannian manifolds with
  $|F(p)F(q)|_{g_N}=|pq|_{g_M}$ for all $p,q\in M$, then $F$ is a Riemannian
  isometry.
  \source{petersen:page-0219}
\end{theorem}
\begin{proof}
  \uses{thm:pet-ch5-myers-steenrod-isometry}
  \emph{Smoothness.} Fix $p\in M$, $q=F(p)$. Introduce distance coordinates
  $(r_{q_1},\dots,r_{q_n})$ near $q$ (\cref{def:pet-ch5-distance-coordinates}) and
  choose $p_i$ with $F(p_i)=q_i$ (possible by distance-preservation and surjectivity
  of $F$, choosing $q_i$ realized as $F(p_i)$). Then
  $r_{q_i}\circ F(x)=|F(x)q_i|=|F(x)F(p_i)|=|xp_i|$; since $|pp_i|=|qq_i|$ the
  functions $x\mapsto|xp_i|$ are themselves smooth near $p$ (being of the same
  distance-coordinate type at $p$), so $(r_{q_1},\dots,r_{q_n})\circ F$ is smooth at
  $p$, and since $\varphi=(r_{q_1},\dots,r_{q_n})$ is a coordinate chart near $q$,
  $F$ is smooth at $p$.

  \emph{Isometry.} It suffices to check $|DF(v)|=|v|$ for $v\in T_pM$. Let
  $c(t)=\exp_p(tv)$; for small $t$, $c$ is a constant-speed segment, so for small
  $t,s$,
  \[
    |t-s|\cdot|v|=|c(t)c(s)|_{g_M}=|F(c(t))F(c(s))|_{g_N},
  \]
  whence
  \[
    |DF(v)|=\left|\frac{d(F\circ c)}{dt}\right|_{t=0}
    =\lim_{t\to0}\frac{|F(c(t))F(c(0))|_{g_N}}{|t|}
    =\lim_{t\to0}\frac{|c(t)c(0)|_{g_M}}{|t|}=|\dot c(0)|=|v|.
  \]
  \source{petersen:page-0219,petersen:page-0220}
\end{proof}

\begin{definition}[submetry]
  \label{def:pet-ch5-submetry}
  \lean{PetersenLib.IsSubmetry}
  \uses{def:pet-ch5-metric-balls}
  $F:(\tilde M,g_{\tilde M})\to(M,g_M)$ is a \emph{submetry} if for every
  $\tilde p\in\tilde M$ there is $r>0$ such that $F(B(\tilde p,\varepsilon))=
  B(F(\tilde p),\varepsilon)$ for all $\varepsilon\le r$. Submetries are locally
  distance nonincreasing, hence continuous, and compositions of submetries (or of
  Riemannian submersions) are again submetries (resp.\ Riemannian submersions).
  \source{petersen:page-0220}
\end{definition}

\begin{theorem}[Thm. 5.6.16 (Berestovskii, 1995) — surjective submetries are $C^1$ submersions]
  \label{thm:pet-ch5-berestovskii-submetry}
  \lean{PetersenLib.surjectiveSubmetry_isC1RiemannianSubmersion}
  \uses{def:pet-ch5-submetry,def:pet-ch5-distance-coordinates}
  If $F:(\tilde M,g_{\tilde M})\to(M,g_M)$ is a surjective submetry, then $F$ is a
  $C^1$ Riemannian submersion.
  \source{petersen:page-0220}
\end{theorem}
\begin{proof}
  \uses{thm:pet-ch5-berestovskii-submetry}
  Write $\tilde p\in F^{-1}(p)$. Choose $r<\operatorname{inj}_p,\operatorname{inj}_{\tilde p}$
  so all relevant geodesic segments are unique. The submetry property gives: if
  $|pq|<r$, then for each $\tilde p\in F^{-1}(p)$ there is a unique
  $\bar q\in F^{-1}(q)$ with $|\tilde p\bar q|=|pq|$, and $\tilde p\mapsto\bar q$ is
  continuous. Define horizontal lifts of unit vectors by
  $\overline{\tilde p\bar q}=\overline{pq}/|pq|$; this is well defined since if
  $\overline{p q_1}=\overline{pq_2}$ then $q_1,q_2$ lie on the same segment from $p$,
  forcing $\bar q_1,\bar q_2$ likewise.

  Select distance coordinates $(r_1,\dots,r_k)$ around $p$
  (\cref{def:pet-ch5-distance-coordinates}); each $r_i$ is a Riemannian submersion,
  hence a submetry, so each $r_i\circ F$ is a submetry (composition of submetries).
  $F$ is $C^1$ iff all $r_i\circ F$ are $C^1$, so it suffices to treat scalar
  submetries $r:U\subset M\to(a,b)$. There, $\nabla r$ is the horizontal lift of
  $\partial_r$ on $(a,b)$, and continuity of $\nabla r$ follows from continuity of
  $\tilde p\mapsto\bar q$ established above.
  \source{petersen:page-0220}
\end{proof}

\begin{remark}[Remark 5.6.17 — submetries are $C^{1,1}$, sharply]
  \label{rem:pet-ch5-submetry-c11}
  \lean{PetersenLib.remark_submetriesAreC11Sharp}
  \uses{thm:pet-ch5-berestovskii-submetry}
  Submetries are in fact $C^{1,1}$ (locally Lipschitz derivative), via local
  Lipschitz continuity of $\tilde p\mapsto\bar q$; this cannot be improved in
  general, e.g.\ for $K=[0,1]^2\subset\mathbb R^2$ the level sets of $r(x)=|xK|$ are
  rounded squares with quarter-circle corners, not $C^2$.
  \source{petersen:page-0220,petersen:page-0221}
\end{remark}

\subsection{The Slice Theorem}

\begin{definition}[proper action; orbit; isotropy group]
  \label{def:pet-ch5-proper-action}
  \lean{PetersenLib.IsProperAction,PetersenLib.orbit,PetersenLib.actionIsotropyGroup}
  An action of a topological group $H$ on $M$ is \emph{proper} if
  $H\times M\to M\times M$, $(h,p)\mapsto(hp,p)$, is proper. The \emph{orbit}
  through $p$ is $Hp=\{hp\mid h\in H\}$; $H\backslash M$ carries the quotient
  topology, making $M\to H\backslash M$ continuous and open, second countable, and
  Hausdorff when the action is proper. The \emph{isotropy group} at $p$ is
  $H_p=\{h\in H\mid hp=p\}$; along an orbit $H_{hp}=hH_ph^{-1}$, $H_p$ is compact
  for proper actions, and $H/H_p\to Hp$ is a homeomorphism. $H$ acts \emph{freely}
  if $H_p=\{e\}$ for all $p$.
  \source{petersen:page-0221}
\end{definition}

\begin{example}[Example 5.6.18 — a nonproper isometric action]
  \label{ex:pet-ch5-nonproper-action}
  \lean{PetersenLib.example_nonproperIsometricAction}
  \uses{def:pet-ch5-proper-action}
  $\operatorname{Iso}(M)$ acts properly on $M$ (Arzelà–Ascoli), but a subgroup
  $H\subset\operatorname{Iso}(M)$ need not act properly unless closed: the action
  $\theta\cdot(z_1,z_2)=(e^{i\theta}z_1,e^{i\alpha\theta}z_2)$ of $\mathbb R$ on
  $S^1\times S^1$ is proper iff $\alpha$ is rational.
  \source{petersen:page-0221}
\end{example}

\begin{theorem}[Thm. 5.6.19 (Myers and Steenrod, 1939) — orbits are submanifolds; $\operatorname{Iso}(M,g)$ is a Lie group]
  \label{thm:pet-ch5-myers-steenrod-lie-group}
  \lean{PetersenLib.closedIsometrySubgroup_orbitsAreSubmanifolds,PetersenLib.isometryGroup_isLieGroup}
  \uses{def:pet-ch5-proper-action,cor:pet-ch5-tubular-neighborhood}
  If $H$ is a closed subgroup of $\operatorname{Iso}(M,g)$, the orbits of the
  action are submanifolds; in particular $\operatorname{Iso}(M,g)$ itself is a Lie
  group.
  \source{petersen:page-0221}
\end{theorem}
\begin{proof}
  \uses{thm:pet-ch5-myers-steenrod-lie-group,cor:pet-ch5-tubular-neighborhood,cor:pet-ch5-uniform-injectivity-radius-compact,prop:pet-ch5-exp-diffeomorphism-properties}
  Work locally, using that $M$ looks Euclidean via exponential coordinates around a
  point (\cref{prop:pet-ch5-exp-diffeomorphism-properties}) and around a tube around
  a submanifold (\cref{cor:pet-ch5-tubular-neighborhood}), so metric balls have
  smooth boundary.

  \emph{Tangent and normal vectors to an orbit.} Say $v\in T_pM$ is \emph{tangent}
  to $Hp$ if $v=\lim\dot c_i(0)$ for geodesic segments $c_i$ from $p$ to
  $p_i=h_ip\in Hp$ with $p_i\to p$; compactness of $H_p$ lets one arrange
  $h_i\to e$, so $Dh_i|_p\to\mathrm{id}$. The set $T_pHp$ of such tangent vectors is a
  linear subspace: it is scale-invariant (reparametrize the geodesic segments), and
  for $v=\lim v_i$, $w=\lim w_i$ with $v_i,w_i$ initial velocities of segments from
  $p$ to $p_i=h_ip$, $q_i\in Hp$ respectively, local Euclidean-ness near $p$ makes
  the (suitably parametrized) geodesic velocity from $p_i$ to $q_i$ close to $w-v$;
  applying $h_i^{-1}$ (moving $p_i$ to $p$, using $h_i\to e$) gives
  $\lim\frac{d(h_i^{-1}\circ c_i)}{dt}(0)=w-v\in T_pHp$. The group action carries
  tangent spaces to tangent spaces, $T_{hp}Hp=Dh(T_pHp)$, and since locally
  $Dh=\exp_p^{-1}\circ h\circ\exp_p$, these subspaces vary continuously along $Hp$.

  Say $v\in T_pM$ is \emph{normal} to $Hp$ if proportional to $\overline{pq}$ with
  $|qp|=|q\,Hp|$; then $B(q,|qp|)\cap Hp=\varnothing$, forcing tangent–normal angles
  $\ge\pi/2$, and since the tangent vectors form a subspace, tangent and normal
  vectors are in fact perpendicular. Consequently, if $O$ is open with smooth
  boundary, $O\cap Hp=\varnothing$, then $q\in\partial O\cap Hp$ gives
  $T_qHp\subset T_q\partial O$.

  \emph{Local graph structure.} Let $N^k$ be a small submanifold with
  $T_pN=T_pHp$; use the normal exponential map to get coordinates $(x,y)$ on a
  tubular neighborhood $N\times B$, $B=B(0,\epsilon)\subset\mathbb R^{n-k}$
  (\cref{cor:pet-ch5-tubular-neighborhood}). For $N,\epsilon$ small, $N\times\{y\}$
  is almost perpendicular to $\{x\}\times B$ at $(x,y)$, and since $Th_pHp$ varies
  continuously with $h$ it has trivial intersection with the tangent spaces of
  $\{x\}\times B$.

  Consider the projection $(x,y)\mapsto x$ from a neighborhood of $p\in Hp$ to a
  neighborhood of $p\in N$. The image is closed in $N$ (as $Hp$ is closed); let
  $N'$ be the complement. If $p_k=h_kp$ and $q_k\in Hp$ project to the same point in
  $N$, then $\overline{p_kq_k}$ is tangent to $B$; but if $p_k,q_k\to p$, then
  $\overline{p_kq_k}$ subconverges to a vector orthogonal to $T_pHp$, and
  $\overline{h_i^{-1}p_k\,h_i^{-1}q_k}$ likewise, a contradiction with tangency to
  $B$. Hence, if $p\in\partial N'$, choosing open $O_i'\subset N'$ with smooth
  boundary meeting $\partial N'$ and shrinking to $p$, the sets $O_i=O_i'\times B$
  avoid $Hp$, so boundary points $p_i=(x_i,y_i)\in\partial O_i\cap Hp\to(p,0)$ give
  $T_pHp\subset T_p\partial O_i$; but $\dim T_pHp=k$ and
  $\dim T_p(\{x_i\}\times B)=n-k$ inside the $(n-1)$-dimensional
  $T_p\partial O_i=T_{x_i}O_i'\oplus T_yB$ forces nontrivial intersection, while
  $T_p(\{x_i\}\times B)\to T_p^\perp N$ is by continuity almost perpendicular to
  $T_pHp$ — a contradiction. So $p\notin\partial N'$, and by shrinking $N$,
  $Hp\cap(N\times B)$ is a continuous graph over $N$.

  This graph is $C^1$: tangent vectors to the orbit are uniquely determined by
  their (almost-orthogonal) projection onto $TN$, so smooth curves in $N$ lift to
  curves in the orbit with continuously varying, uniquely determined velocity, and
  this lifted velocity is exactly the graph's derivative.

  \emph{Smoothness bootstrap and Lie group structure.} $\operatorname{Iso}(M,g)$
  acts properly and freely on the open subset $O\subset M^{n+1}$ of
  $(p_0,\dots,p_n)$ with $\overline{p_0p_i}$, $i=1,\dots,n$, linearly independent,
  identifying $\operatorname{Iso}(M,g)$ with a $C^1$ submanifold of $O$ via the
  orbit map. The formulas $Th_pHp=Dh(T_pHp)$, $Dh=\exp_{hp}^{-1}\circ h\circ
  \exp_p$ show that $THp\subset TM$ is as smooth as $H$: if $H$ is $C^k$, $k\ge1$,
  so is $THp$, hence $Hp$ is $C^{k+1}$ by the graph argument above, and the same
  orbit-map construction then shows $H$ itself is $C^{k+1}$. Iterating from
  $C^1$ upgrades $H=\operatorname{Iso}(M,g)$ to a smooth ($C^\infty$) Lie group; the
  same bootstrap applies to any closed subgroup $H$, whose orbits are then smooth
  submanifolds.
  \source{petersen:page-0221,petersen:page-0222,petersen:page-0223}
\end{proof}

\begin{proposition}[Prop. 5.6.20 — kernel of the orbit map]
  \label{prop:pet-ch5-orbit-map-differential}
  \lean{PetersenLib.orbitMapDifferentialKernel}
  \uses{thm:pet-ch5-myers-steenrod-lie-group}
  Let $H$ be a closed subgroup of $\operatorname{Iso}(M,g)$ with Lie algebra
  $\mathfrak h=T_eH$, and $\mathfrak h_p=\{v\in\mathfrak h\mid\exp(tv)\in H_p\ \forall t\}$
  the subalgebra of $H_p$. For the orbit map $O_p(h)=hp$,
  $\ker\big((DO_p)|_v\big)=\mathfrak h_p$, and more generally
  $\ker\big((DO_p)|_s\big)=DL_s(\mathfrak h_p)$.
  \source{petersen:page-0224}
\end{proposition}
\begin{proof}
  \uses{prop:pet-ch5-orbit-map-differential}
  The general statement follows from the case at $e$ by the chain rule and
  $(h_1j)p=h_1(jp)$. At $e$: $(DO_p)|_v(v)=\frac{d}{dt}(\exp(tv)\cdot p)|_{t=0}$, and
  by the one-parameter group property,
  \[
    \frac{d}{dt}(\exp(tv)\cdot p)\Big|_{t=t_0}
    =D(\exp(t_0v))\left(\frac{d}{ds}(\exp(sv)\cdot p)\Big|_{s=0}\right).
  \]
  If $(DO_p)|_v(v)=0$ (i.e.\ the $s=0$ derivative vanishes), this forces
  $\exp(tv)\cdot p=p$ for all $t$ (the derivative of $t\mapsto\exp(tv)p$ vanishes
  identically), so $v\in\mathfrak h_p$. Conversely if $v\in\mathfrak h_p$ then
  $\exp(tv)p\equiv p$, so its derivative at $t=0$, namely $(DO_p)|_v(v)$, vanishes.
  \source{petersen:page-0224}
\end{proof}

\begin{theorem}[Thm. 5.6.21 — the free slice theorem]
  \label{thm:pet-ch5-free-slice-theorem}
  \lean{PetersenLib.freeSliceTheorem}
  \uses{def:pet-ch5-proper-action,cor:pet-ch5-tubular-neighborhood,thm:pet-ch5-myers-steenrod-lie-group}
  If $H\subset\operatorname{Iso}(M,g)$ is closed and acts freely, $H\backslash M$
  has a smooth manifold structure and Riemannian metric making $M\to H\backslash M$
  a Riemannian submersion.
  \source{petersen:page-0224}
\end{theorem}
\begin{proof}
  \uses{thm:pet-ch5-free-slice-theorem}
  By \cref{thm:pet-ch5-myers-steenrod-lie-group}, orbits are properly embedded
  copies of $H$ (free action, \cref{prop:pet-ch5-orbit-map-differential} with
  $\mathfrak h_p=0$ gives an immersion, and properness of the group action makes it
  an embedding). The $H$-invariant map $E:H\times T_pM\to M$,
  $(h,v)\mapsto h\exp_p(v)=\exp_{hp}(Dh|_pv)$, restricted to the normal bundle at
  $p$ gives $\exp^\perp:H\times T_p^\perp Hp\to M$; the natural trivialization
  $(h,v)\mapsto Dh|_p(v)$, $H\times T_p^\perp Hp\to T^\perp Hp$, is a fiberwise
  linear isometry identifying $\exp^\perp$ with the normal exponential map
  $T^\perp Hp\to M$. By \cref{cor:pet-ch5-tubular-neighborhood} this is a
  diffeomorphism from a neighborhood of the zero section of $H\times T_p^\perp Hp$
  onto a neighborhood of $Hp$ in $M$.

  To get a \emph{uniform} tube $\exp^\perp:H\times B(0,\epsilon)\to M$: choose
  $\epsilon>0$ so $\exp^\perp:U\times B(0,\epsilon)\to M$, $e\in U\subset H$, is an
  embedding, with closed $\epsilon$-balls in the image compact (so having compact
  intersection with orbits); shrink further so that $v\in B(0,\epsilon)$,
  $|(hp)\exp_p(v)|<\epsilon$ implies $h\in U$ — this makes $p$ the unique closest
  point of $Hp$ to $\exp_p(v)$, and by the first variation formula, any segment from
  $\exp_p(v)$ to $Hp$ meets $Hp$ perpendicularly and at a point $hp$, $h\in U$; two
  normal vectors mapping to the same point would then violate the choice of
  $\epsilon$, forcing injectivity of $\exp^\perp$ on $H\times B(0,\epsilon)$.
  Nonsingularity holds since it holds at $(e,v)$, $|v|<\epsilon$, and
  $\exp^\perp(h,v)=h\exp_p(v)$ with $h$ acting as a diffeomorphism. Properness of
  $\exp^\perp$ follows since $Hp$ is properly embedded; so $\exp^\perp:H\times
  B(0,\epsilon)\to M$ is a closed embedding.

  Thus $M\to H\backslash M$ is locally a trivial bundle: each slice $B(0,\epsilon)$
  maps homeomorphically to its image in $H\backslash M$, giving smooth charts on
  $H\backslash M$ with smooth transition maps. The Riemannian metric on
  $H\backslash M$ is induced by identifying $T_{Hp}(H\backslash M)$ with
  $T_p^\perp Hp$; since $Dh|_p$ isometrically identifies $T_p^\perp Hp$ with
  $T_{hp}^\perp Hp$ for every $h\in H$, all such normal spaces along the orbit are
  isometric, so this is well defined and makes $M\to H\backslash M$ a Riemannian
  submersion.
  \source{petersen:page-0224,petersen:page-0225}
\end{proof}

\begin{remark}[Remark 5.6.22 — homogeneous quotients $H/K$]
  \label{rem:pet-ch5-quotient-homogeneous}
  \lean{PetersenLib.remark_homogeneousQuotient}
  \uses{thm:pet-ch5-free-slice-theorem}
  For $K\subset H$ compact, averaging any metric on $H$ over $K$ (acting by
  $(k,x)\mapsto xk^{-1}$) makes it right-$K$-invariant, giving a free isometric
  action to which \cref{thm:pet-ch5-free-slice-theorem} applies, so $H/K$ is a
  manifold with a Riemannian submersion metric from $H$; if the metric on $H$ is
  also left-invariant, $H$ acts by isometries on $H/K$, making it homogeneous.
  \source{petersen:page-0225}
\end{remark}

\begin{remark}[closed subgroups and right-invariant metrics]
  \label{rem:pet-ch5-closed-subgroup-quotient}
  \lean{PetersenLib.remark_closedSubgroupRightInvariantQuotient}
  \uses{rem:pet-ch5-quotient-homogeneous}
  For $K\subset H$ merely closed, right multiplication by $K$ is still a proper
  action, made isometric by a right-$K$-invariant metric on $H$; but it may not be
  possible to simultaneously keep the metric on $H$ left-invariant, so $H$ need not
  act by isometries on $H/K$ in this more general case.
  \source{petersen:page-0226}
\end{remark}

\begin{corollary}[Corollary 5.6.23 — orbits of proper isometric actions are properly embedded]
  \label{cor:pet-ch5-proper-action-orbit-embedding}
  \lean{PetersenLib.properIsometricAction_orbitProperlyEmbedded}
  \uses{prop:pet-ch5-orbit-map-differential,def:pet-ch5-proper-action}
  For a proper isometric action $H\times M\to M$ and $p\in M$, the orbit map
  $H/H_p\to Hp$ is a proper embedding.
  \source{petersen:page-0226}
\end{corollary}
\begin{proof}
  \uses{cor:pet-ch5-proper-action-orbit-embedding}
  The orbit map is already known to be a proper injective map with $H/H_p$ a
  manifold; \cref{prop:pet-ch5-orbit-map-differential} shows its differential is
  injective (kernel trivial modulo $\mathfrak h_p$, which is exactly what is
  quotiented out in $H/H_p$), so it is an immersion, and a proper injective
  immersion is a (proper) embedding.
  \source{petersen:page-0226}
\end{proof}

\begin{definition}[slice representation]
  \label{def:pet-ch5-slice-representation}
  \lean{PetersenLib.sliceRepresentation}
  \uses{cor:pet-ch5-proper-action-orbit-embedding}
  The \emph{slice representation} of a proper isometric action $H\times M\to M$ at
  $p$ is the linear representation $H_p\times T_p^\perp Hp\to T_p^\perp Hp$,
  $(h,v)\mapsto Dh|_p(v)$. Letting $H_p$ act on $H$ on the right, it also acts on
  $H\times T_p^\perp Hp$, and by \cref{cor:pet-ch5-proper-action-orbit-embedding}
  the quotient $H\times_{H_p}T_p^\perp Hp$ has a natural manifold structure.
  \source{petersen:page-0226}
\end{definition}

\begin{theorem}[Thm. 5.6.24 — the slice theorem]
  \label{thm:pet-ch5-slice-theorem}
  \lean{PetersenLib.sliceTheorem}
  \uses{def:pet-ch5-slice-representation,thm:pet-ch5-free-slice-theorem}
  Let $H\subset\operatorname{Iso}(M,g)$ be closed. For $p\in M$ there is a map
  $\exp^\perp:H\times_{H_p}T_p^\perp Hp\to M$ that is a diffeomorphism on a uniform
  tubular neighborhood $H\times_{H_p}B(0,\epsilon)$ onto an $\epsilon$-neighborhood
  of $Hp$.
  \source{petersen:page-0226}
\end{theorem}
\begin{proof}
  \uses{thm:pet-ch5-slice-theorem}
  Proceed as in \cref{thm:pet-ch5-free-slice-theorem}, using
  \cref{cor:pet-ch5-proper-action-orbit-embedding} for the general (not necessarily
  free) case. For fixed $p$, the bundle map $H\times T_p^\perp Hp\to T^\perp Hp$,
  $(h,v)\mapsto Dh|_p(v)$, is $H$-invariant, an isomorphism on fibers, and sends
  $H\times\{0\}$ to the zero section (the orbit $Hp$). Since $D(h\circ
  k^{-1})|_p(Dk|_p(v))=Dh|_p(v)$ for $k\in H_p$, this descends to a bundle
  isomorphism $H\times_{H_p}T_p^\perp Hp\to T^\perp Hp$, via which
  $\exp^\perp:H\times_{H_p}T_p^\perp Hp\to M$ is defined as the normal exponential
  map. As in \cref{thm:pet-ch5-free-slice-theorem} one finds $\epsilon>0$ so that
  $\exp^\perp:H\times_{H_p}B(0,\epsilon)\to M$ is an $H$-invariant embedding.
  \source{petersen:page-0226}
\end{proof}

\begin{corollary}[Corollary 5.6.25 — isotropy near an orbit]
  \label{cor:pet-ch5-isotropy-near-orbit}
  \lean{PetersenLib.isotropyGroupNearOrbit}
  \uses{thm:pet-ch5-slice-theorem}
  For small $v\in T_p^\perp Hp$, $H_{\exp_p(v)}=\{h\in H_p\mid Dh|_pv=v\}$.
  \source{petersen:page-0227}
\end{corollary}
\begin{proof}
  \uses{cor:pet-ch5-isotropy-near-orbit}
  By \cref{thm:pet-ch5-slice-theorem}, $\exp^\perp$ is an $H$-equivariant
  diffeomorphism near $Hp$, so $h\cdot\exp_p(v)=\exp_p(v)$ for small $v$ iff
  $h\in H_p$ (fixing the base point of the slice) and $h$ fixes the slice
  coordinate $v$ under the linear slice representation, i.e.\ $Dh|_pv=v$.
  \source{petersen:page-0227}
\end{proof}

\begin{corollary}[Corollary 5.6.26 — constant isotropy type gives a Riemannian quotient]
  \label{cor:pet-ch5-conjugate-isotropy-manifold}
  \lean{PetersenLib.constantIsotropyType_riemannianQuotient}
  \uses{cor:pet-ch5-isotropy-near-orbit,thm:pet-ch5-free-slice-theorem}
  If $H\subset\operatorname{Iso}(M,g)$ is closed and all isotropy groups are
  conjugate, the orbit space is a Riemannian manifold and $M\to H\backslash M$ a
  Riemannian submersion.
  \source{petersen:page-0227}
\end{corollary}
\begin{proof}
  \uses{cor:pet-ch5-conjugate-isotropy-manifold}
  By \cref{cor:pet-ch5-isotropy-near-orbit}, when all isotropy groups $H_q$,
  $q\in M$, are conjugate to a fixed $H_p$, the slice representation of $H_p$ on
  $T_p^\perp Hp$ has trivial (or uniformly conjugate) stabilizers on a neighborhood
  of $0$, so the local model $H\times_{H_p}T_p^\perp Hp$ has the same isotropy type
  throughout, matching the free case of \cref{thm:pet-ch5-free-slice-theorem}
  applied fiberwise; the slice charts of \cref{thm:pet-ch5-slice-theorem} then patch
  to a manifold structure on $H\backslash M$ exactly as in the free case, with the
  induced metric making $M\to H\backslash M$ a Riemannian submersion.
  \source{petersen:page-0227}
\end{proof}

\section{Completeness}

\subsection{The Hopf–Rinow Theorem}

\begin{theorem}[Thm. 5.7.1 (Hopf and Rinow, 1931) — equivalence of completeness notions]
  \label{thm:pet-ch5-hopf-rinow}
  \lean{PetersenLib.hopfRinowTheorem}
  \uses{def:pet-ch5-geodesically-complete,def:pet-ch5-distance,cor:pet-ch5-metric-ball-image,cor:pet-ch5-regular-curve-approx,cor:pet-ch5-segments-are-geodesics}
  For a Riemannian manifold $(M,g)$ the following are equivalent:
  \begin{enumerate}
  \item[(1)] $M$ is geodesically complete.
  \item[(2)] $M$ is geodesically complete at some $p$ (all geodesics through $p$
  are defined for all time).
  \item[(3)] $M$ satisfies the Heine–Borel property: every closed bounded set is
  compact.
  \item[(4)] $M$ is metrically complete.
  \end{enumerate}
  \source{petersen:page-0227}
\end{theorem}
\begin{proof}
  \uses{thm:pet-ch5-hopf-rinow}
  $(1)\Rightarrow(2)$ and $(3)\Rightarrow(4)$ are immediate.

  $(4)\Rightarrow(1)$: a maximal geodesic $c:[0,b)\to M$ with $b<\infty$ must, by
  \cref{prop:pet-ch5-leaves-compact-set}, leave every compact set as $t\to b$; since
  $|\dot c|$ is constant, $c(t)$, $t\to b$, is then a Cauchy sequence with no limit
  in $M$, violating metric completeness. So every maximal geodesic has $b=\infty$
  (and symmetrically $a=-\infty$).

  $(2)\Rightarrow(3)$: fix $p$ as in (2). It suffices to show
  $\exp_p(\bar B(0,R))=\bar B(p,R)$ for all $R$ (the inclusion $\subset$ is
  automatic), i.e.\ every $q\in M$ is joined to $p$ by a segment. By
  \cref{cor:pet-ch5-metric-ball-image}, choose $\varepsilon>0$ with
  $\bar B(p,\varepsilon)=\exp_p(\bar B(0,\varepsilon))$ compact and every point
  there joined to $p$ by a minimal geodesic. Fix $q\in M$; let
  $c(t):[0,\infty)\to M$ be the unit-speed geodesic (defined for all $t$ by (2))
  with $c(0)=p$, and choose $p'\in\bar B(p,\varepsilon)-B(p,\varepsilon)$ closest to
  $q$ among such boundary points. We first check $|pq|=\varepsilon+|p'q|$: otherwise
  \cref{cor:pet-ch5-regular-curve-approx} gives a unit-speed $c'\in\Omega_{p,q}$
  with $L(c')<|pp'|+|p'q|=\varepsilon+|p'q|$; choosing $t$ with
  $c'(t)\in\bar B(p,\varepsilon)-B(p,\varepsilon)$, since
  $t+|c'(t)q|\le L(c')<\varepsilon+|p'q|$ and $t=\varepsilon$ (as $c'(t)$ is at
  distance $\varepsilon$ from $p$ along a curve of length $t$ up to that point, so
  $t\ge\varepsilon$, and equality holds because $c'(t)\in\bar B(p,\varepsilon)$), we
  get $|c'(t)q|<|p'q|$, contradicting the choice of $p'$. Hence with $p'=c(\varepsilon)$
  (adjusting $c$ so its initial segment realizes $\bar B(p,\varepsilon)$), $|pq|=
  \varepsilon+|c(\varepsilon)q|$.

  Let $A=\{t\in[0,|pq|]\mid|pq|=t+|c(t)q|\}$; then $0,\varepsilon\in A$. If $t\in A$,
  the triangle inequality
  $|pq|=t+|c(t)q|\ge|pc(t)|+|c(t)q|\ge|pq|$ forces $t=|pc(t)|$, i.e.\ $c|_{[0,t]}$
  is a segment. If $a\in A$ and $t<a$, then
  \[
    |pq|\le|pc(t)|+|c(t)q|\le|pc(t)|+|c(t)c(a)|+|c(a)q|\le t+(a-t)+|c(a)q|=a+|c(a)q|=|pq|,
  \]
  forcing equality throughout, so $t=|pc(t)|$ and $t\in A$: thus $[0,a]\subset A$
  whenever $a\in A$. Since $t\mapsto|c(t)q|$ is continuous, $A$ is closed. Finally,
  if $a\in A$, \cref{cor:pet-ch5-metric-ball-image} gives $\delta>0$ so every point
  of $\bar B(c(a),\delta)$ is joined to $c(a)$ by a segment; choosing $q'\in
  \bar B(c(a),\delta)-B(c(a),\delta)$ closest to $q$,
  \[
    |pq|=a+|c(a)q|=a+|c(a)q'|+|q'q|=a+\delta+|q'q|\ge|pq'|+|q'q|\ge|pq|,
  \]
  forcing $|pq'|=a+\delta$; then the piecewise geodesic $p\to c(a)\to q'$ is a
  segment, hence a genuine geodesic by \cref{cor:pet-ch5-segments-are-geodesics},
  so $q'=c(a+\delta)$, and $|pq|=(a+\delta)+|q'q|$ shows $a+\delta\in A$. Thus $A$
  is open in $[0,|pq|]$ near every point below $|pq|$, closed, and nonempty, so
  $A=[0,|pq|]$; taking $t=|pq|$, $c(|pq|)=q$ and $c|_{[0,|pq|]}$ is a segment from
  $p$ to $q$.
  \source{petersen:page-0227,petersen:page-0228,petersen:page-0229}
\end{proof}

\begin{corollary}[Corollary 5.7.2 — complete manifolds have segments everywhere]
  \label{cor:pet-ch5-complete-segment-exists}
  \lean{PetersenLib.completeManifold_allPointsJoinedBySegment}
  \uses{thm:pet-ch5-hopf-rinow}
  If $(M,g)$ is complete (in any of the equivalent senses), any two points of $M$
  are joined by a segment.
  \source{petersen:page-0229}
\end{corollary}
\begin{proof}
  \uses{cor:pet-ch5-complete-segment-exists}
  This is exactly what was shown in the proof of $(2)\Rightarrow(3)$ of
  \cref{thm:pet-ch5-hopf-rinow}.
  \source{petersen:page-0229}
\end{proof}

\begin{corollary}[Corollary 5.7.3 — a proper Lipschitz function forces completeness]
  \label{cor:pet-ch5-proper-lipschitz-complete}
  \lean{PetersenLib.properLipschitzFunction_impliesComplete}
  \leanok
  \uses{thm:pet-ch5-metric-topology}
  If $(M,g)$ admits a proper Lipschitz function $f:M\to\mathbb R$, then $M$ is
  complete.
  \source{petersen:page-0229}
\end{corollary}
\begin{proof}
  \leanok
  \uses{cor:pet-ch5-proper-lipschitz-complete}
  We verify the Heine–Borel property. Every closed metric ball $\bar B(x,r)$ maps
  under $f$ into $[f(x)-Kr,\,f(x)+Kr]$ (Lipschitz constant $K$), so it is contained
  in the compact set $f^{-1}([f(x)-Kr,f(x)+Kr])$ (properness); being closed, $\bar
  B(x,r)$ is itself compact. Thus every closed ball is compact, $M$ is a proper
  metric space, and a proper metric space is complete. (This is the elementary half
  of \cref{thm:pet-ch5-hopf-rinow}: no geodesic machinery is used, only the metric
  structure of \cref{thm:pet-ch5-metric-topology}.)
  \source{petersen:page-0229}
\end{proof}

\subsection{Warped Product Characterization}

\begin{theorem}[Thm. 5.7.4 (Tashiro, 1965) — global warped product from a conformal Hessian]
  \label{thm:pet-ch5-tashiro}
  \lean{PetersenLib.tashiroWarpedProductTheorem}
  \uses{thm:pet-ch5-hopf-rinow,def:pet-ch5-exponential-map}
  Let $(M,g)$ be a complete Riemannian $n$-manifold admitting a nontrivial $f$
  with $\operatorname{Hess}f=\lambda g$. Then $(M,g)$ is isometric to a complete
  warped product metric of one of the forms:
  \begin{enumerate}
  \item[(1)] $M=\mathbb R\times N$, $g=dr^2+\rho^2(r)g_N$;
  \item[(2)] $M=\mathbb R^n$, $g=dr^2+\rho^2(r)ds_{n-1}^2$, $r>0$;
  \item[(3)] $M=S^n$, $g=dr^2+\rho^2(r)ds_{n-1}^2$, $r\in[a,b]$.
  \end{enumerate}
  \source{petersen:page-0229,petersen:page-0230}
\end{theorem}
\begin{proof}
  \uses{thm:pet-ch5-tashiro}
  Locally (cited from the local warped-product characterization for
  Hessian-conformal functions), $|\nabla f|$ is locally constant on
  $\{f=f_0\}\cap\{df\ne0\}$ for each $f_0$, so connected components of this set are
  closed. As $f$ is nontrivial, $df\ne0$ somewhere; let $N$ be a nonempty connected
  component of $\{f=f_0\}\cap\{df\ne0\}$, a closed hypersurface. For $p\in N$ let
  $c_p(t)=\exp_p\!\big(t\,\nabla f|_p/|\nabla f|\big)$, the unit-speed geodesic
  through $p$ normal to $N$.

  Fix $a<0<b$ and let $C\subset N$ be the (open, by openness of the regular set) set
  of $p$ for which $f$ is regular at $c_p(t)$ for all $t\in[a,b]$. On
  $U=\{c_p(t)\mid p\in C,t\in[a,b]\}$ the local warped product structure gives
  $g|_U=dr^2+\rho^2(r)g_N$ with $r$ the signed distance to $N$,
  $\nabla r=\nabla f/|\nabla f|$, $f=\int\rho(r)\,dr$. For every $p\in N$,
  $\frac{d^2(f\circ c_p)}{dt^2}=\operatorname{Hess}f(\dot c_p,\dot c_p)=\lambda\circ c_p$
  with $(f\circ c_p)(0)=f_0$, $\frac{d(f\circ c_p)}{dt}(0)=|\nabla f|$; on $U$,
  $\lambda\circ c_p=\lambda(f\circ c_p)$ depends only on the value of
  $f\circ c_p$, so this second-order ODE with initial data independent of
  $p\in C$ forces $f\circ c_p(t)$ to be independent of $p\in C$, and likewise
  $\frac{d(f\circ c_p)}{dt}(t)=|\nabla f|_{c_p(t)}|=\rho(t)$ is independent of
  $p\in C$. Continuity of $|\nabla f|_{c_p(t)}|$ in $(p,t)$ and regularity of $f$
  on $U$ prevent $|\nabla f|_{c_p(t)}|$ from vanishing as $p\to\partial C$ (for
  fixed $t$), so $C$ is also closed in $N$; as $N$ is connected, $C=N$.

  Let $(a,b)\ni0$ be the maximal interval on which $c_p(t)$ is regular for all
  $p\in N$; the warped product structure $dr^2+\rho^2(r)g_N$ extends to
  $(a,b)\times N$. If $(a,b)=\mathbb R$ this is case (1). If, say, $b<\infty$, the
  level set $\{r=b\}$ consists of critical points of $f$, so $\rho=|\nabla f|\to0$
  as $t\to b$; the warped product structure then forces any two points of $N$ to
  converge to each other as $t\to b$, so $\{r=b\}$ is a single point $x$, and the
  set of unit vectors $\{\dot c_p(b)\}_{p\in N}\subset T_xM$ is both open and closed
  (as $N$ is a closed submanifold), hence all unit vectors at $x$: every geodesic
  $c_p(t)$ reaches $x$ at $t=b$ and, continuing past $x$, coincides with another
  such geodesic. Level sets of $f$ near $x$ are exactly distance spheres centered
  at $x$, and the argument identifying $g_N$ as the round metric on $S^{n-1}$
  proceeds as in the local characterization, giving $N\simeq S^{n-1}$ and
  $\{f_0\le f<f(x)\}\simeq[0,b)\times N$ with $x$ the unique critical point in
  $\{f_0\le f\le f(x)\}$. A symmetric argument applies if $a>-\infty$. This gives
  case (3) if both $a,b$ are finite, and case (2) if exactly one is.
  \source{petersen:page-0230,petersen:page-0231}
\end{proof}

\begin{theorem}[Thm. 5.7.5 — global classification for $\operatorname{Hess}f=(\alpha f+\beta)g$]
  \label{thm:pet-ch5-hessian-classification}
  \lean{PetersenLib.affineHessianGlobalClassification}
  \uses{thm:pet-ch5-tashiro}
  Let $(M,g)$ be complete of dimension $n$ with a nontrivial $f$,
  $\operatorname{Hess}f=(\alpha f+\beta)g$, $\alpha,\beta\in\mathbb R$. Then
  $(M,g)$ is one of:
  \begin{enumerate}
  \item[(a)] $(I\times S^{n-1},\,dr^2+\mathrm{sn}_k^2(r)ds_{n-1}^2)$: a constant
  curvature space form;
  \item[(b)] $(\mathbb R\times N,\,dr^2+g_N)$: a product metric;
  \item[(c)] $\big(\mathbb R\times N,\,dr^2+(A\exp(\sqrt\alpha r)+B\exp(-\sqrt\alpha r))^2h\big)$,
  $A,B\ge0$.
  \end{enumerate}
  \source{petersen:page-0231}
\end{theorem}
\begin{proof}
  \uses{thm:pet-ch5-hessian-classification}
  By \cref{thm:pet-ch5-tashiro}, $g=dr^2+\rho^2(r)g_N$, $r\in I$, with $f=f(r)$,
  $\rho(r)=f'(r)$, $\lambda=\alpha f+\beta=f''(r)$.

  If $\alpha=\beta=0$, $\rho$ is constant: case (b).

  If $\alpha=0$, $\beta\ne0$, $\rho=\beta r+\gamma$ is affine; $I$ is a half-line on
  which $\rho$ vanishes at the finite boundary point (a single critical point of
  $f$), and the resulting metric is the standard Euclidean metric — this is the
  degenerate case of (a) with $k=0$.

  If $\alpha\ne0$, replace $f$ by $f+\beta/\alpha$ (so
  $\operatorname{Hess}(f+\beta/\alpha)=\alpha(f+\beta/\alpha)g$), reducing to
  $\beta=0$. If $\alpha<0$, $\rho=A\sin(\sqrt{-\alpha}\,r+r_0)$: $I$ is a compact
  interval and the metric is a round sphere, case (a) with $k=-\alpha>0$. If
  $\alpha>0$, $\rho=A\exp(\sqrt\alpha r)+B\exp(-\sqrt\alpha r)$ with at least one of
  $A,B>0$; if $A,B\ge0$ this is case (c) directly, while if $A,B$ have opposite
  sign, $\rho=C\sinh(\sqrt\alpha r+r_0)$, $I$ a half-line, giving the constant
  negatively curved warped product, case (a) with $k=-\alpha<0$.
  \source{petersen:page-0231,petersen:page-0232}
\end{proof}

\begin{remark}[Remark 5.7.6 — critical points and transnormal functions]
  \label{rem:pet-ch5-obata-transnormal}
  \lean{PetersenLib.remark_obataTransnormalFunctions}
  \uses{thm:pet-ch5-hessian-classification}
  In case (a) $f$ has at least one critical point, while in (b), (c) it has none.
  Obata (1961) established the round-sphere case via $\operatorname{Hess}f=(1-f)g$.
  Transnormal functions on a complete manifold give a similar topological
  decomposition, with zero, one, or two critical values, all level sets (including
  critical ones) smooth submanifolds, and $\tilde f=\phi(r)$ for $r$ the signed
  distance to a fixed level set.
  \source{petersen:page-0232}
\end{remark}

\subsection{The Segment Domain}

\begin{definition}[segment domain and its star interior]
  \label{def:pet-ch5-segment-domain}
  \lean{PetersenLib.segmentDomain,PetersenLib.segmentDomainStarInterior}
  \uses{def:pet-ch5-segment,thm:pet-ch5-hopf-rinow}
  For a complete $(M,g)$ and $p\in M$, the \emph{segment domain} is
  $\operatorname{seg}(p)=\{v\in T_pM\mid\exp_p(tv):[0,1]\to M\text{ is a segment}\}$,
  a closed star-shaped set with $M=\exp_p(\operatorname{seg}(p))$ (by
  \cref{thm:pet-ch5-hopf-rinow}). Its \emph{star interior} is
  $\operatorname{seg}^0(p)=\{sv\mid s\in[0,1),\,v\in\operatorname{seg}(p)\}$.
  \source{petersen:page-0232,petersen:page-0233}
\end{definition}

\begin{proposition}[Prop. 5.7.7 — injectivity of $\exp_p$ on the star interior]
  \label{prop:pet-ch5-segment-domain-injective}
  \lean{PetersenLib.expMap_injectiveOnSegmentDomainStarInterior}
  \uses{def:pet-ch5-segment-domain,cor:pet-ch5-segments-are-geodesics}
  If $x\in\exp_p(\operatorname{seg}^0(p))$, it is joined to $p$ by a unique
  segment; in particular $\exp_p$ is injective on $\operatorname{seg}^0(p)$.
  \source{petersen:page-0233}
\end{proposition}
\begin{proof}
  \uses{prop:pet-ch5-segment-domain-injective}
  Let $\sigma:[0,1]\to M$ be a segment with $\sigma(0)=p$, $\sigma(t_0)=x$,
  $t_0<1$. If $\tilde\sigma:[0,t_0]\to M$ were a different segment from $p$ to $x$,
  the concatenation $c(s)=\tilde\sigma(s)$ for $s\in[0,t_0]$, $\sigma(s)$ for
  $s\in[t_0,1]$, would be a piecewise smooth segment (same total minimal length)
  that is not smooth at $t_0$ unless $\tilde\sigma=\sigma|_{[0,t_0]}$, contradicting
  \cref{cor:pet-ch5-segments-are-geodesics} (segments are geodesics, hence smooth,
  and geodesics are uniquely determined by initial data) unless the two agree.
  \source{petersen:page-0233}
\end{proof}

\begin{lemma}[Lemma 5.7.8 — nonsingularity of $\exp_p$ on the star interior]
  \label{lem:pet-ch5-exp-nonsingular-segment-domain}
  \lean{PetersenLib.expMap_nonsingularOnSegmentDomainStarInterior}
  \uses{prop:pet-ch5-segment-domain-injective,cor:pet-ch5-uniform-injectivity-radius-compact}
  $D\exp_p$ is nonsingular at every point of $\operatorname{seg}^0(p)$.
  \source{petersen:page-0233}
\end{lemma}
\begin{proof}
  \uses{lem:pet-ch5-exp-nonsingular-segment-domain}
  The set of singular points of $\exp_p$ in $\operatorname{seg}^0(p)$ is closed;
  suppose it is nonempty and let $v$ be singular with $\exp_p$ nonsingular at all
  $tv$, $t\in[0,1)$. Since $c(t)=\exp_p(tv)$ is then an embedding on $[0,1)$, there
  are neighborhoods $U\ni[0,1)v$ in $T_pM$ and $V\supset c([0,1))$ in $M$ with
  $\exp_p:U\to V$ a diffeomorphism; note $v\notin U$, $c(1)\notin V$.

  For $w\in T_vT_pM$, extend to a Jacobi field $J$ on $T_pM$ commuting with the
  radial field, $[\partial_r,J]=0$, and push forward via $\exp_p$ to a vector field
  (also $J$) on $V$ commuting with $\partial_r$ there. Since $D\exp_p|_vw=0$,
  $\lim_{t\to1}J|_{\exp_p(tv)}=\lim_{t\to1}D\exp_p(J)|_{\exp_p(tv)}=0$: $D\exp_p$ is
  singular at $v$ exactly when $\exp_p(v)$ is a conjugate point along $c$.

  Since $|J|^2\to0$ as $t\to1$, $\log|J|^2\to-\infty$, so there is a sequence
  $t_n\to1$ with $\partial_r|J|^2/|J|^2|_{\exp_p(t_nv)}\to-\infty$. By the radial
  Jacobi equation $\partial_r|J|^2=2\operatorname{Hess}r(J,J)$, so
  $\operatorname{Hess}r(J,J)/|J|^2|_{\exp_p(t_nv)}\to-\infty$.

  Since $c(t)=\exp_p(tv)$ is a segment on $[0,1+\varepsilon]$ for some
  $\varepsilon>0$ (by hypothesis $\exp_p$ nonsingular strictly beyond $t=1$ locally
  and $v\in\operatorname{seg}^0(p)$ means $c$ extends slightly past $t=1$ as a
  segment), choose $\varepsilon$ small enough, by
  \cref{cor:pet-ch5-uniform-injectivity-radius-compact}, that
  $\tilde r(x)=|xc(1+\varepsilon)|$ is smooth on $B(c(1+\varepsilon),2\varepsilon)$.
  Consider $e(x)=r(x)+\tilde r(x)$; by the triangle inequality
  $e(x)\ge1+\varepsilon=|pc(1+\varepsilon)|$, with equality along $c(t)$,
  $t\in[0,1+\varepsilon]$. So $e$ attains an absolute minimum along $c$ and has
  nonnegative Hessian there, i.e.\ $\operatorname{Hess}e(J,J)\ge0$. But
  \[
    \frac{\operatorname{Hess}e(J,J)}{|J|^2}\Big|_{\exp_p(t_nv)}
    =\frac{\operatorname{Hess}r(J,J)}{|J|^2}\Big|_{\exp_p(t_nv)}
    +\frac{\operatorname{Hess}\tilde r(J,J)}{|J|^2}\Big|_{\exp_p(t_nv)}\to-\infty,
  \]
  since $\operatorname{Hess}\tilde r$ is bounded near $c(1)$ while the
  $\operatorname{Hess}r$ term diverges to $-\infty$: a contradiction. So $\exp_p$
  is nonsingular throughout $\operatorname{seg}^0(p)$.
  \source{petersen:page-0233,petersen:page-0234,petersen:page-0235}
\end{proof}

\begin{lemma}[Lemma 5.7.9 — the star boundary of $\operatorname{seg}(p)$]
  \label{lem:pet-ch5-cut-locus-characterization}
  \lean{PetersenLib.starBoundaryOfSegmentDomain}
  \uses{lem:pet-ch5-exp-nonsingular-segment-domain}
  If $v\in\operatorname{seg}(p)-\operatorname{seg}^0(p)$, then either (1) there is
  $w\ne v$ in $\operatorname{seg}(p)$ with $\exp_p(v)=\exp_p(w)$, or (2) $D\exp_p$
  is singular at $v$.
  \source{petersen:page-0235}
\end{lemma}
\begin{proof}
  \uses{lem:pet-ch5-cut-locus-characterization}
  Let $c(t)=\exp_p(tv)$. For $t>1$, choose segments $\sigma_t:[0,1]\to M$ with
  $\sigma_t(0)=p$, $\sigma_t(1)=c(t)$; since $c|_{[0,1]}$ is not a segment beyond
  $t=1$, $\dot\sigma_t(0)$ is never a positive multiple of $\dot c(0)=v$. Pick
  $t_n\to1$ with $\dot\sigma_{t_n}(0)\to w\in T_pM$; since
  $L(\sigma_{t_n})=|\dot\sigma_{t_n}(0)|\to L(c|_{[0,1]})=|v|$, we get $|w|=|v|$.

  If $w\ne v$, it cannot be a positive multiple of $v$ (having the same norm), so
  (1) holds with this $w$. If $w=v$: were $D\exp_p$ nonsingular at $v$, $\exp_p$
  would embed a neighborhood of $v$; then $\dot\sigma_{t_n}(0)\to v$ together with
  $\exp_p(\dot\sigma_{t_n}(0))=\exp_p(t_n\dot c(0))$ forces
  $\dot\sigma_{t_n}(0)=t_n v$ for large $n$ (by injectivity near $v$), making
  $c|_{[0,t_n]}$, $t_n>1$, a segment — contradicting $v\notin\operatorname{seg}^0(p)$
  (in fact contradicting maximality of $\operatorname{seg}(p)$ at $v$). So
  $D\exp_p$ is singular at $v$, case (2).
  \source{petersen:page-0235}
\end{proof}

\begin{proposition}[Prop. 5.7.10 — the star interior is open]
  \label{prop:pet-ch5-segment-domain-open}
  \lean{PetersenLib.segmentDomainStarInterior_isOpen}
  \uses{lem:pet-ch5-cut-locus-characterization,lem:pet-ch5-exp-nonsingular-segment-domain}
  $\operatorname{seg}^0(p)$ is open in $T_pM$.
  \source{petersen:page-0235}
\end{proposition}
\begin{proof}
  \uses{prop:pet-ch5-segment-domain-open}
  Fix $v\in\operatorname{seg}^0(p)$ and a neighborhood $V\subset T_pM$ of $v$ on
  which $\exp_p$ is a diffeomorphism onto its image (by
  \cref{lem:pet-ch5-exp-nonsingular-segment-domain} and the inverse function
  theorem). Let $v_i\in V\to v$; for each $i$, choose $w_i\in\operatorname{seg}(p)$
  with $\exp_p(v_i)=\exp_p(w_i)$ (possible since $M=\exp_p(\operatorname{seg}(p))$).
  If $w_i$ has an accumulation point $w\ne v$, then $v\notin\operatorname{seg}^0(p)$
  by \cref{lem:pet-ch5-cut-locus-characterization} applied in the limit — contrary
  to hypothesis; so $w_i\to v$, hence $w_i\in V$ for large $i$, and injectivity of
  $\exp_p$ on $V$ forces $w_i=v_i$. Thus $v_i\in\operatorname{seg}(p)$ for large
  $i$; $\exp_p$ is nonsingular at $v_i$ (as $v_i\in V$), and since $w_i=v_i$,
  alternative (1) of \cref{lem:pet-ch5-cut-locus-characterization} cannot hold at
  $v_i$, so $v_i\notin\operatorname{seg}(p)-\operatorname{seg}^0(p)$, i.e.\
  $v_i\in\operatorname{seg}^0(p)$. Hence $V\cap\operatorname{seg}^0(p)$ is a
  neighborhood of $v$.
  \source{petersen:page-0235,petersen:page-0236}
\end{proof}

\begin{definition}[cut locus]
  \label{def:pet-ch5-cut-locus}
  \lean{PetersenLib.cutLocus}
  \uses{prop:pet-ch5-segment-domain-open}
  The set $\operatorname{seg}(p)-\operatorname{seg}^0(p)\subset T_pM$ is the
  \emph{cut locus} of $p$ (in $T_pM$). By
  \cref{prop:pet-ch5-segment-domain-open,lem:pet-ch5-exp-nonsingular-segment-domain},
  $r(x)=|px|$ is smooth on the open dense set $U_p=\{p\}\cup\exp_p(\operatorname{seg}^0(p))$
  and fails to be smooth off $U_p$.
  \source{petersen:page-0236}
\end{definition}

\begin{corollary}[Corollary 5.7.11 — smoothness of $r$ up to the cut point]
  \label{cor:pet-ch5-cut-locus-smoothness}
  \lean{PetersenLib.distanceFunction_smoothUpToCutPoint}
  \uses{def:pet-ch5-cut-locus,lem:pet-ch5-cut-locus-characterization}
  Let $c:[0,\infty)\to M$ be a geodesic, $c(0)=p$, and
  $\operatorname{cut}(\dot c(0))=\sup\{t\mid c|_{[0,t]}\text{ is a segment}\}$. Then
  $r$ is smooth at $c(t)$ for $t<\operatorname{cut}(\dot c(0))$ but not at
  $x=c(\operatorname{cut}(\dot c(0)))$, the failure being either that $\exp_p$ is
  not one-to-one at $x$ or that $x$ is a critical value of $\exp_p$.
  \source{petersen:page-0236}
\end{corollary}
\begin{proof}
  \uses{cor:pet-ch5-cut-locus-smoothness}
  By definition, $t_0v:=\dot c(0)\operatorname{cut}(\dot c(0))$ lies in
  $\operatorname{seg}(p)-\operatorname{seg}^0(p)$, while $tv\in
  \operatorname{seg}^0(p)$ for $t<\operatorname{cut}(\dot c(0))$: this is exactly
  the statement of \cref{def:pet-ch5-cut-locus} that $r$ is smooth on
  $\exp_p(\operatorname{seg}^0(p))\ni c(t)$ but not on the complement, in
  particular not at $x=\exp_p(t_0v)$. The dichotomy for the failure at $x$ is
  exactly \cref{lem:pet-ch5-cut-locus-characterization}.
  \source{petersen:page-0236}
\end{proof}

\subsection{The Injectivity Radius}

\begin{remark}[injectivity radius via the cut locus]
  \label{rem:pet-ch5-injectivity-radius-cut-locus}
  \lean{PetersenLib.injectivityRadius_viaCutLocus}
  \uses{def:pet-ch5-cut-locus,def:pet-ch5-injectivity-radius}
  In a complete manifold, $\operatorname{inj}(p)$ (the largest $\varepsilon$ with
  $\exp_p:B(0,\varepsilon)\to B(p,\varepsilon)$ a diffeomorphism) equals $|v|$ for
  $v$ the closest point of the cut locus
  $\operatorname{seg}(p)-\operatorname{seg}^0(p)$ to the origin.
  \source{petersen:page-0236}
\end{remark}

\begin{lemma}[Lemma 5.7.12 (Klingenberg) — characterization of the injectivity radius]
  \label{lem:pet-ch5-klingenberg}
  \lean{PetersenLib.klingenbergLemma}
  \uses{rem:pet-ch5-injectivity-radius-cut-locus,lem:pet-ch5-cut-locus-characterization}
  Suppose $v\in\operatorname{seg}(p)-\operatorname{seg}^0(p)$ with
  $|v|=\operatorname{inj}(p)$. Either
  \begin{enumerate}
  \item[(1)] there is exactly one other $v'$ with $\exp_p(v')=\exp_p(v)$, and $v'$
  satisfies $\frac{d}{dt}\big|_{t=1}\exp_p(tv')=-\frac{d}{dt}\big|_{t=1}\exp_p(tv)$
  (the two segments meet to form a smooth geodesic loop at $x=\exp_p(v)$), or
  \item[(2)] $x=\exp_p(v)$ is a critical value of $\exp_p|_{\operatorname{seg}(p)}$.
  \end{enumerate}
  \source{petersen:page-0236}
\end{lemma}
\begin{proof}
  \uses{lem:pet-ch5-klingenberg}
  Suppose $x$ is a regular value and $c_1,c_2:[0,1]\to M$ are segments from $p$ to
  $x=\exp_p(v)$ with $\dot c_1(1)\ne-\dot c_2(1)$. Choose $w\in T_pM$... rather,
  work at $x$: choose $w\in T_xM$ with $g(w,\dot c_1(1))<0$ and $g(w,\dot c_2(1))<0$
  (possible since the two vectors are not antipodal). Let $c(s)$ be a curve with
  $\dot c(0)=w$; since $D\exp_p$ is nonsingular at $\dot c_i(0)$ (regularity of $x$,
  using \cref{lem:pet-ch5-cut-locus-characterization}), there are unique curves
  $v_i(s)\in T_pM$ with $v_i(0)=\dot c_i(0)$ and $D\exp_p(v_i(s))=c(s)$ near
  $s=0$, well defined for $s$ near $0$ by the inverse function theorem. The
  first-variation formula (via $g(w,\dot c_i(1))<0$) gives that
  $s\mapsto|v_i(s)|=|pc(s)|$ is strictly decreasing at $s=0$, so for small $s>0$,
  $|v_i(s)|<|v|=|px|$. This exhibits two distinct points $v_1(s),v_2(s)$ near
  $v_1(0),v_2(0)$ (hence not equal to each other for generic small $s$, as
  $\dot c_1(1)\ne\dot c_2(1)$ forces the two curves $v_i(s)$ apart to first order)
  both mapping under $\exp_p$ to $c(s)$ with $|v_i(s)|<|v|=\operatorname{inj}(p)$: but
  $\exp_p$ is injective on $B(0,\operatorname{inj}(p))\subset\operatorname{seg}^0(p)$
  by definition of the injectivity radius, so having two preimages of $c(s)$ inside
  this ball is a contradiction. Hence $\dot c_1(1)=-\dot c_2(1)$ whenever $x$ is a
  regular value with two distinct segments to it, giving case (1); if $x$ has only
  one segment, injectivity of $\exp_p$ fails at the boundary in the alternative
  manner of \cref{lem:pet-ch5-cut-locus-characterization}, i.e.\ case (2) (critical
  value).
  \source{petersen:page-0236,petersen:page-0237}
\end{proof}

\section{Further Study}

Several textbooks on Riemannian geometry treat most of the basic material of this
chapter, uniformly emphasizing the variational approach as the fundamental
technique for proving geodesic-minimizing results; a dedicated treatment of the
variational method is also recommended for further reading.
\source{petersen:page-0237}

\section{Exercises}

\begin{remark}[Exercise 5.9.1 — fixed-time existence implies completeness]
  \label{rem:pet-ch5-ex-1}
  \lean{PetersenLib.exercise5_9_1}
  \uses{def:pet-ch5-geodesically-complete}
  If $(M,g)$ has the property that all unit-speed geodesics exist for a fixed time
  $\varepsilon>0$, show $(M,g)$ is geodesically complete.
  \source{petersen:page-0237}
\end{remark}

\begin{remark}[Exercise 5.9.2 — chain rule for velocity and acceleration]
  \label{rem:pet-ch5-ex-2}
  \lean{PetersenLib.exercise5_9_2}
  \uses{def:pet-ch5-acceleration}
  For $c:I\to M$ and $\phi:J\to I$, show
  $\frac{d(c\circ\phi)}{dt}=\dot c\circ\phi\cdot\frac{d\phi}{dt}$ and
  $\frac{d^2(c\circ\phi)}{dt^2}=\dot c\circ\phi\cdot\frac{d^2\phi}{dt^2}
  +\ddot c\circ\phi\cdot\big(\frac{d\phi}{dt}\big)^2$.
  \source{petersen:page-0237,petersen:page-0238}
\end{remark}

\begin{remark}[Exercise 5.9.3 — geodesic equation with orthogonal coordinates]
  \label{rem:pet-ch5-ex-3}
  \lean{PetersenLib.exercise5_9_3}
  \uses{def:pet-ch5-geodesic}
  If the coordinate vector fields of a chart are orthogonal ($g_{ij}=0$, $i\ne j$),
  show the geodesic equations become
  $\frac{d}{dt}\big(g_{ii}\frac{dc^i}{dt}\big)=\frac12\sum_j\frac{\partial g_{jj}}
  {\partial x^i}\big(\frac{dc^j}{dt}\big)^2$.
  \source{petersen:page-0238}
\end{remark}

\begin{remark}[Exercise 5.9.4 — reparametrization to a geodesic]
  \label{rem:pet-ch5-ex-4}
  \lean{PetersenLib.exercise5_9_4}
  \uses{def:pet-ch5-geodesic}
  Show a regular curve can be reparametrized to be a geodesic iff its acceleration
  is everywhere tangent to the curve.
  \source{petersen:page-0238}
\end{remark}

\begin{remark}[Exercise 5.9.5 — a complete open submanifold is everything]
  \label{rem:pet-ch5-ex-5}
  \lean{PetersenLib.exercise5_9_5}
  \uses{def:pet-ch5-geodesically-complete}
  If $O\subset(M,g)$ is open and $(O,g)$ is complete, show $O=M$.
  \source{petersen:page-0238}
\end{remark}

\begin{remark}[Exercise 5.9.6 — Misner completeness implies completeness]
  \label{rem:pet-ch5-ex-6}
  \lean{PetersenLib.exercise5_9_6}
  \uses{def:pet-ch5-geodesically-complete}
  Call $(M,g)$ \emph{Misner complete} if every geodesic $c:(a,b)\to M$ with
  $b-a<\infty$ lies in a compact set. Show Misner completeness implies
  completeness.
  \source{petersen:page-0238}
\end{remark}

\begin{remark}[Exercise 5.9.7 — length-minimizers reparametrize a segment]
  \label{rem:pet-ch5-ex-7}
  \lean{PetersenLib.exercise5_9_7}
  \uses{def:pet-ch5-segment}
  For $c\in\Omega_{p,q}$ with $L(c)=|pq|$: (1) show $L(c|_{[a,b]})=|c(a)c(b)|$ for
  all $a,b\in[0,1]$; (2) show there is a segment $\sigma\in\Omega_{p,q}$ and
  monotone $\varphi:[0,1]\to[0,1]$ (not necessarily everywhere smooth) with
  $c=\sigma\circ\varphi$.
  \source{petersen:page-0238}
\end{remark}

\begin{remark}[Exercise 5.9.8 — completeness under a larger metric]
  \label{rem:pet-ch5-ex-8}
  \lean{PetersenLib.exercise5_9_8}
  \uses{def:pet-ch5-geodesically-complete}
  If $(M,g)$ is metrically complete and $\bar g\ge g$ another metric on $M$, show
  $(M,\bar g)$ is metrically complete.
  \source{petersen:page-0238}
\end{remark}

\begin{remark}[Exercise 5.9.9 — homogeneous manifolds are complete]
  \label{rem:pet-ch5-ex-9}
  \lean{PetersenLib.exercise5_9_9}
  \uses{def:pet-ch5-geodesically-complete}
  A Riemannian manifold is \emph{homogeneous} if its isometry group acts
  transitively. Show homogeneous manifolds are geodesically complete.
  \source{petersen:page-0238}
\end{remark}

\begin{remark}[Exercise 5.9.10 — completeness of $\mathbb R\times N$ warped products]
  \label{rem:pet-ch5-ex-10}
  \lean{PetersenLib.exercise5_9_10}
  \uses{def:pet-ch5-geodesically-complete}
  For $(M,g)=(\mathbb R\times N,dr^2+g_r)$ with $(N,g_r)$ complete for all
  $r\in\mathbb R$ (e.g.\ $g_r=\rho^2(r)g_N$, $\rho:\mathbb R\to(0,\infty)$, $(N,g_N)$
  metrically complete), show $(M,g)$ is metrically complete.
  \source{petersen:page-0238}
\end{remark}

\begin{remark}[Exercise 5.9.11 — completeness on a half-line warped product]
  \label{rem:pet-ch5-ex-11}
  \lean{PetersenLib.exercise5_9_11}
  \uses{def:pet-ch5-geodesically-complete}
  For $(M,g)=((0,\infty)\times N,dr^2+\rho^2(r)g_N)$, $\rho:(0,\infty)\to
  (0,\infty)$, $(N,g_N)$ complete: give examples that are complete and examples
  that are not.
  \source{petersen:page-0238}
\end{remark}

\begin{remark}[Exercise 5.9.12 — splitting via a Hessian-flat submersion to $\mathbb R^k$]
  \label{rem:pet-ch5-ex-12}
  \lean{PetersenLib.exercise5_9_12}
  \uses{def:pet-ch5-geodesically-complete}
  If $F:(M,g)\to(\mathbb R^k,g_{\mathbb R^k})$ is a Riemannian submersion, $(M,g)$
  complete, and each component of $F$ has zero Hessian, show
  $(M,g)=(N,h)\times(\mathbb R^k,g_{\mathbb R^k})$.
  \source{petersen:page-0238}
\end{remark}

\begin{remark}[Exercise 5.9.13 — filling a gap in the submetry theorem]
  \label{rem:pet-ch5-ex-13}
  \lean{PetersenLib.exercise5_9_13}
  \uses{thm:pet-ch5-berestovskii-submetry}
  Find and fill in the gap in the proof of \cref{thm:pet-ch5-berestovskii-submetry}.
  \source{petersen:page-0239}
\end{remark}

\begin{remark}[Exercise 5.9.14 — isotropic implies homogeneous]
  \label{rem:pet-ch5-ex-14}
  \lean{PetersenLib.exercise5_9_14}
  \uses{def:pet-ch5-fixed-point-set}
  A manifold is \emph{isotropic} at $p$ if $\operatorname{Iso}_p$ acts transitively
  on the unit sphere in $T_pM$. Show isotropy at every point implies homogeneity.
  \source{petersen:page-0239}
\end{remark}

\begin{remark}[Exercise 5.9.15 — orthonormal coordinates are distance coordinates]
  \label{rem:pet-ch5-ex-15}
  \lean{PetersenLib.exercise5_9_15}
  \uses{rem:pet-ch5-distance-function-geodesics}
  If coordinates satisfy $g_{ii}=\delta_{ij}$ (i.e.\ $g_{i1}=\delta_{i1}$ for the
  first coordinate direction), show $x^1$ is a distance function.
  \source{petersen:page-0239}
\end{remark}

\begin{remark}[Exercise 5.9.16 — distance functions persist under a compatible metric change]
  \label{rem:pet-ch5-ex-16}
  \lean{PetersenLib.exercise5_9_16}
  \uses{rem:pet-ch5-distance-function-geodesics}
  Let $r:U\to\mathbb R$ be a distance function w.r.t.\ $g$. If $\bar g$ satisfies
  $\bar g(\nabla r,v)=g(\nabla r,v)$ for all $v$ (gradient taken w.r.t.\ $g$), show
  $r$ is also a distance function w.r.t.\ $\bar g$.
  \source{petersen:page-0239}
\end{remark}

\begin{remark}[Exercise 5.9.17 — projective models of $S^n(R)$ and $H^n(R)$]
  \label{rem:pet-ch5-ex-17}
  \lean{PetersenLib.exercise5_9_17}
  \uses{ex:pet-ch5-sphere-great-circles,ex:pet-ch5-hyperbolic-hyperbolas}
  For $x\in\mathbb R^{n+1}$, $x^{n+1}>0$, the projection to $x^{n+1}=R$ is
  $P(x)=R(x^1/x^{n+1},\dots,x^n/x^{n+1},1)$. (1) Verify this formula; (2) show
  geodesics of $S^n(R)$, $H^n(R)$ are intersections with $2$-planes through the
  origin; (3) show the upper hemisphere of $S^n(R)$ projects onto all of
  $x^{n+1}=R$; (4) show $H^n(R)$ projects onto an open disc of radius $R$; (5) show
  geodesics of $S^n(R)$, $H^n(R)$ project to straight lines.
  \source{petersen:page-0239}
\end{remark}

\begin{remark}[Exercise 5.9.18 — a conformal change making a manifold complete]
  \label{rem:pet-ch5-ex-18}
  \lean{PetersenLib.exercise5_9_18}
  \uses{cor:pet-ch5-proper-lipschitz-complete}
  Show any $(M,g)$ admits a conformal change $(M,\lambda^2g)$ that is complete;
  take $\lambda:M\to[1,\infty)$ proper and rapidly growing.
  \source{petersen:page-0239}
\end{remark}

\begin{remark}[Exercise 5.9.19 — induced versus ambient distance on subsets of $\mathbb R^n$]
  \label{rem:pet-ch5-ex-19}
  \lean{PetersenLib.exercise5_9_19}
  \uses{def:pet-ch5-distance}
  On $U\subset\mathbb R^n$ compare the Riemannian-induced distance to the ambient
  $\mathbb R^n$ distance: (1) give a non-convex $U$ where they agree; (2) give a
  convex $U$ where they do not.
  \source{petersen:page-0239}
\end{remark}

\begin{remark}[Exercise 5.9.20 — the $T$-tensor detects totally geodesic submanifolds]
  \label{rem:pet-ch5-ex-20}
  \lean{PetersenLib.exercise5_9_20}
  \uses{def:pet-ch5-fixed-point-set}
  For $M\subset(\bar M,\bar g)$, using the $T$-tensor, show $T=0$ on $M$ iff
  $M\subset(\bar M,\bar g)$ is totally geodesic.
  \source{petersen:page-0239}
\end{remark}

\begin{remark}[Exercise 5.9.21 — convexity along geodesics via the Hessian]
  \label{rem:pet-ch5-ex-21}
  \lean{PetersenLib.exercise5_9_21}
  \uses{def:pet-ch5-geodesic}
  For $f:(M,g)\to\mathbb R$ smooth: (1) compute the first and second derivatives of
  $f\circ c$ along a geodesic $c$; (2) deduce that at a local max/min of $f$ the
  gradient vanishes and the Hessian is nonpositive/nonnegative; (3) show $f$ has
  everywhere nonnegative Hessian iff $f\circ c$ is convex for every geodesic $c$.
  \source{petersen:page-0239,petersen:page-0240}
\end{remark}

\begin{remark}[Exercise 5.9.22 — leading term of the volume density]
  \label{rem:pet-ch5-ex-22}
  \lean{PetersenLib.exercise5_9_22}
  \uses{def:pet-ch5-radial-distance-function}
  If the volume form near a point is $\lambda(r,\theta)\,dr\wedge\mathrm{vol}_{n-1}$,
  show $\lambda(r,\theta)=r^{n-1}+O(r^{n+1})$.
  \source{petersen:page-0240}
\end{remark}

\begin{remark}[Exercise 5.9.23 — segments to a properly embedded submanifold]
  \label{rem:pet-ch5-ex-23}
  \lean{PetersenLib.exercise5_9_23}
  \uses{def:pet-ch5-segment,thm:pet-ch5-hopf-rinow}
  Let $N\subset M$ be properly embedded, $(M,g)$ complete, $|xN|=\inf\{|xp|\mid
  p\in N\}$, and $\sigma:[a,b]\to M$ a unit-speed segment from $N$ to $x$. (1) Show
  $\sigma$ is also a segment from $N$ to $\sigma(t)$, $t<b$, and $\dot\sigma(a)\perp
  N$; (2) show every point of $M$ has a segment to $N$ when $N$ is closed; (3) show
  there is always an open neighborhood of $N$ where points are joined to $N$ by
  segments; (4) show $r(x)=|xN|$ is smooth on a punctured neighborhood of $N$; (5)
  show integral curves of $\nabla r$ are the geodesics perpendicular to $N$.
  \source{petersen:page-0240}
\end{remark}

\begin{remark}[Exercise 5.9.24 — cut locus on the square torus]
  \label{rem:pet-ch5-ex-24}
  \lean{PetersenLib.exercise5_9_24}
  \uses{def:pet-ch5-cut-locus}
  Find the cut locus on the square torus $\mathbb R^2/\mathbb Z^2$.
  \source{petersen:page-0240}
\end{remark}

\begin{remark}[Exercise 5.9.25 — cut locus on the sphere and $\mathbb{RP}^n$]
  \label{rem:pet-ch5-ex-25}
  \lean{PetersenLib.exercise5_9_25}
  \uses{def:pet-ch5-cut-locus}
  Find the cut locus on the sphere and on real projective space with the constant
  curvature metrics.
  \source{petersen:page-0240}
\end{remark}

\begin{remark}[Exercise 5.9.26 — first-order comparison of exponential images]
  \label{rem:pet-ch5-ex-26}
  \lean{PetersenLib.exercise5_9_26}
  \uses{def:pet-ch5-exponential-map}
  Show $|\exp_p(v)\exp_p(w)|=|v-w|+O(r^2)$ for $|v|,|w|\le r$.
  \source{petersen:page-0240}
\end{remark}

\begin{remark}[Exercise 5.9.27 — second-order asymptotics of $\operatorname{Hess}r$]
  \label{rem:pet-ch5-ex-27}
  \lean{PetersenLib.exercise5_9_27}
  \uses{lem:pet-ch5-normal-coords-first-order}
  In exponential normal coordinates $x^i$ at $p$: (1) show
  $(\operatorname{Hess}x^i)_{kl}=\Gamma^i_{kl}=O(r)$; (2) using
  $\tfrac12r^2=\tfrac12\sum_i(x^i)^2$ and $g=\delta_{ij}+O(r^2)$, show
  $\operatorname{Hess}\tfrac12r^2=g+O(r^2)$; (3) show $\operatorname{Hess}r=
  \tfrac1rg_r+O(r)$.
  \source{petersen:page-0240,petersen:page-0241}
\end{remark}

\begin{remark}[Exercise 5.9.28 — one-sided directional derivatives of distance to a submanifold]
  \label{rem:pet-ch5-ex-28}
  \lean{PetersenLib.exercise5_9_28}
  \uses{def:pet-ch5-segment,lem:pet-ch5-first-variation-formula,rem:pet-ch5-ex-26}
  Let $(M,g)$ be complete, $K\subset M$ compact (or properly embedded), $r(x)=|xK|$.
  (1) Show that if $r$ is differentiable at $x\notin K$ then $\overline{xK}$
  contains only one vector; (2) for unit-speed $c$, show that if $f=r\circ c$ is
  differentiable at $t$, all vectors $\overline{c(t)K}$ form the same angle with
  $\dot c(t)$; (3) show the upper right Dini derivative satisfies
  $\overline D^+f(t_0)\le g(\dot c(t_0),-\overline{c(t_0)K})$, via the first
  variation formula for a variation toward $K$; (4) show
  $|c(t_0)c(t)|=h+O(h^2)$ for $h=t-t_0$ small; (5)–(6) via a comparison point on a
  segment to $K$ and \cref{rem:pet-ch5-ex-26}, show the lower right Dini derivative
  satisfies $D^+f(t_0)\ge\min_{\overline{c(t_0)K}}g(\dot c(t_0),-\overline{c(t_0)K})$;
  (7) conclude the one-sided derivatives of $f$ exist everywhere.
  \source{petersen:page-0241}
\end{remark}

\begin{remark}[Exercise 5.9.29 — metric-space length agrees with the Riemannian length]
  \label{rem:pet-ch5-ex-29}
  \lean{PetersenLib.exercise5_9_29}
  \uses{def:pet-ch5-curve-length,rem:pet-ch5-ex-26}
  In a metric space $(X,|\cdot|)$, define
  $L(c)=\sup\{\sum|c(t_i)c(t_{i+1})|\mid a=t_1\le\dots\le t_k=b\}$. (1) Show a
  curve of finite length is absolutely continuous; (2) show the Cantor step
  function is a counterexample to the converse; (3) show this definition recovers
  the Riemannian curve length for smooth (and absolutely continuous) curves, using
  \cref{rem:pet-ch5-ex-26} to approximate the Euclidean case; (4) show an
  absolutely continuous $c:[a,b]\to M$ of length $|c(a)c(b)|$ factors as
  $c=\sigma\circ\varphi$ for a segment $\sigma$ and monotone $\varphi$.
  \source{petersen:page-0242}
\end{remark}

\begin{remark}[Exercise 5.9.30 — a synthetic characterization of exponential coordinates]
  \label{rem:pet-ch5-ex-30}
  \lean{PetersenLib.exercise5_9_30}
  \uses{rem:pet-ch5-radial-isometry}
  If coordinates $x^i$ near $p$ satisfy $x^i(p)=0$ and $g_{ij}x^j=\delta_{ij}x^j$,
  show they must be exponential normal coordinates. Hint: set
  $r=\sqrt{\delta_{ij}x^ix^j}$, show it is a smooth distance function away from
  $p$, and that the gradient's integral curves are geodesics emanating from $p$.
  \source{petersen:page-0242}
\end{remark}

\begin{remark}[Exercise 5.9.31 — intersections of totally geodesic submanifolds]
  \label{rem:pet-ch5-ex-31}
  \lean{PetersenLib.exercise5_9_31}
  \uses{def:pet-ch5-fixed-point-set}
  If $N_1,N_2\subset M$ are totally geodesic, show each component of $N_1\cap N_2$
  is a totally geodesic submanifold, with tangent space at $p$ the Zariski tangent
  space $T_pN_1\cap T_pN_2$.
  \source{petersen:page-0242}
\end{remark}

\begin{remark}[Exercise 5.9.32 — an isometry with differential $-I$]
  \label{rem:pet-ch5-ex-32}
  \lean{PetersenLib.exercise5_9_32}
  \uses{prop:pet-ch5-isometry-uniqueness}
  Let $F:(M,g)\to(M,g)$ be an isometry fixing $p$. Show $DF|_p=-I$ on $T_pM$ iff
  $F^2=\mathrm{id}_M$ and $p$ is an isolated fixed point.
  \source{petersen:page-0242}
\end{remark}

\begin{remark}[Exercise 5.9.33 — functional distance equals distance when complete]
  \label{rem:pet-ch5-ex-33}
  \lean{PetersenLib.exercise5_9_33}
  \uses{def:pet-ch5-functional-distance,thm:pet-ch5-hopf-rinow}
  Show that for a complete manifold, the functional distance $d_F$ equals the
  Riemannian distance $d$.
  \source{petersen:page-0242}
\end{remark}

\begin{remark}[Exercise 5.9.34 — regular exponential images are local energy minima]
  \label{rem:pet-ch5-ex-34}
  \lean{PetersenLib.exercise5_9_34}
  \uses{thm:pet-ch5-energy-minima-geodesics,lem:pet-ch5-exp-nonsingular-segment-domain}
  Let $c:[0,1]\to M$ be a geodesic such that $\exp_{c(0)}$ is regular at $tc'(0)$
  for all $t\le1$. Show $c$ is a local minimum of the energy functional. Hint: the
  lift of $c$ via $\exp_{c(0)}$ is a minimizing geodesic in the pullback metric.
  \source{petersen:page-0242}
\end{remark}

\begin{remark}[Exercise 5.9.35 — bi-invariant metrics and the Lie exponential]
  \label{rem:pet-ch5-ex-35}
  \lean{PetersenLib.exercise5_9_35}
  \uses{def:pet-ch5-exponential-map,thm:pet-ch5-hopf-rinow}
  For a Lie group $G$ with bi-invariant (pseudo-)Riemannian metric: (1) show
  homomorphisms $\mathbb R\to G$ are precisely the integral curves of left-invariant
  vector fields through $e$; (2) show geodesics through $e$ are precisely such
  homomorphisms, so the Lie exponential coincides with the Riemannian exponential
  at $e$; (3) for a Riemannian (positive-definite) bi-invariant metric, show every
  $x\in G$ has a square root, using metric completeness; (4) show
  $\mathrm{SL}(n,\mathbb R)$ admits no bi-invariant Riemannian metric.
  \source{petersen:page-0242,petersen:page-0243}
\end{remark}

\begin{remark}[Exercise 5.9.36 — Riemannian submersions are submetries]
  \label{rem:pet-ch5-ex-36}
  \lean{PetersenLib.exercise5_9_36}
  \uses{def:pet-ch5-submetry}
  Show a Riemannian submersion is a submetry.
  \source{petersen:page-0243}
\end{remark}

\begin{remark}[Exercise 5.9.37 (Hermann) — completeness descends along submersions and gives a fibration]
  \label{rem:pet-ch5-ex-37}
  \lean{PetersenLib.exercise5_9_37}
  \uses{def:pet-ch5-submetry,thm:pet-ch5-hopf-rinow}
  Let $F:(M,g_M)\to(N,g_N)$ be a Riemannian submersion. (1) Show $(N,g_N)$ is
  complete if $(M,g_M)$ is; (2) show $F$ is a fibration (locally
  $F^{-1}(U)\cong U\times F^{-1}(p)$) if $(M,g_M)$ is complete, and give a
  counterexample when it is not.
  \source{petersen:page-0243}
\end{remark}

\begin{remark}[Exercise 5.9.38 — commuting isometries and even-codimension fixed sets]
  \label{rem:pet-ch5-ex-38}
  \lean{PetersenLib.exercise5_9_38}
  \uses{rem:pet-ch5-fixed-point-codimension}
  Let $S$ be a set of orientation-preserving isometries of $(M,g)$ that pairwise
  commute. Show each component of $\operatorname{Fix}(S)$ has even codimension.
  \source{petersen:page-0243}
\end{remark}

\begin{remark}[Exercise 5.9.39 — affine maps]
  \label{rem:pet-ch5-ex-39}
  \lean{PetersenLib.exercise5_9_39}
  \uses{def:pet-ch5-geodesic,prop:pet-ch5-isometry-uniqueness}
  A local diffeomorphism $F:(M,g_M)\to(N,g_N)$ is \emph{affine} if
  $F_*(\nabla^M_XY)=\nabla^N_{F_*X}F_*Y$ for all vector fields $X,Y$. (1) Show
  affine maps carry geodesics to geodesics; (2) show $F$ is determined by $F(p)$
  and $DF|_p$; (3) give an affine $\mathbb R^n\to\mathbb R^n$ that is not an
  isometry.
  \source{petersen:page-0243}
\end{remark}

\begin{remark}[Exercise 5.9.40 — projective linear groups and isometry groups of space forms]
  \label{rem:pet-ch5-ex-40}
  \lean{PetersenLib.exercise5_9_40}
  \uses{rem:pet-ch5-ex-39,ex:pet-ch5-sphere-great-circles,ex:pet-ch5-hyperbolic-hyperbolas}
  For $\mathbb{FP}^n$ ($\mathbb F=\mathbb R,\mathbb C$): (1) $\mathrm{GL}(n{+}1,
  \mathbb F)$ acts on $1$-dimensional subspaces of $\mathbb F^{n+1}$; (2) the
  kernel $H$ of this action is $\{\lambda I_{n+1}\}$, a normal subgroup; (3)
  $\mathrm{PGL}(n{+}1,\mathbb F)=\mathrm{GL}(n{+}1,\mathbb F)/H$ acts with each
  element determined by its value and differential at one point; (4) no
  Riemannian metric on $\mathbb{FP}^n$ makes this action isometric; (5) the action
  is affine (\cref{rem:pet-ch5-ex-39}) for the submersion metric; (6) the isometry
  group of $\mathbb{RP}^n$ is $\mathrm{PO}(n{+}1)$; (7) that of $\mathbb{CP}^n$ is
  $\mathrm{PU}(n{+}1)$; (8) that of $H^n(\mathbb R)$ is $\mathrm{PO}(n,1)$; (9) for
  $G=\{[\begin{smallmatrix}O&v\\0&1\end{smallmatrix}]\mid O\in O(n),v\in\mathbb R^n\}
  \subset\mathrm{GL}(n{+}1,\mathbb R)$ (as in Exercise 1.6.9), show $\mathrm{PG}=G$.
  \source{petersen:page-0243}
\end{remark}

\begin{remark}[Exercise 5.9.41 — the compact-open topology on $\operatorname{Iso}(M,g)$]
  \label{rem:pet-ch5-ex-41}
  \lean{PetersenLib.exercise5_9_41}
  \uses{thm:pet-ch5-myers-steenrod-lie-group,def:pet-ch5-proper-action}
  Let $\operatorname{Diff}(M;K,O)=\{F\in\operatorname{Diff}(M)\mid F(K)\subset O\}$.
  (1) Finite intersections of such sets ($K$ compact, $O$ open) form a topology
  base — the compact-open topology; (2) it is second countable; (3) on a
  Riemannian manifold, convergence in it is uniform convergence on compact sets;
  (4) a sequence in $\operatorname{Iso}(M,g)$ converges in it iff it converges
  pointwise (Arzelà–Ascoli); (5) $\operatorname{Iso}_0(M,g)$ is locally compact
  (Arzelà–Ascoli); (6) $\operatorname{Iso}_0(M,g)$ is compact; (7)
  $\operatorname{Iso}(M,g)$ acts properly on $M$; (8) $F\mapsto(F(p),DF|_p)$ is
  continuous on $\operatorname{Iso}(M,g)$; (9) this evaluation map is a
  homeomorphism onto its image.
  \source{petersen:page-0244}
\end{remark}

\begin{remark}[Exercise 5.9.42 — Bianchi identities and the curvature tensor in normal coordinates]
  \label{rem:pet-ch5-ex-42}
  \lean{PetersenLib.exercise5_9_42}
  \uses{lem:pet-ch5-normal-coords-first-order}
  In exponential normal coordinates at $p$, all computed at $p$: (1) taking three
  derivatives of $x^i=\sum_sg_{is}x^s$ (as in
  \cref{lem:pet-ch5-normal-coords-first-order}), show
  $\partial_i\partial_kg_{jl}+\partial_j\partial_ig_{kl}+\partial_k\partial_jg_{il}=0$;
  (2) using all four cyclic versions, show
  $\partial_i\partial_jg_{kl}=\partial_k\partial_lg_{ij}$; (3) using the curvature
  tensor formula in normal coordinates, show
  $R_{ikjl}=\partial_i\partial_lg_{kj}-\partial_i\partial_jg_{kl}$; (4) combine (1)
  and (3) to get $\partial_i\partial_jg_{kl}=\tfrac13(R_{ikjl}+R_{jkil})$.
  \source{petersen:page-0244,petersen:page-0245}
\end{remark}

\begin{remark}[Exercise 5.9.43 — Riemannian volume expansion (A. Gray)]
  \label{rem:pet-ch5-ex-43}
  \lean{PetersenLib.exercise5_9_43}
  \uses{rem:pet-ch5-ex-42}
  With notation as in \cref{rem:pet-ch5-ex-42}: (1) show
  $\sqrt{\det(g_{ij})}=1-\tfrac16\mathrm{Ric}_{ij}x^ix^j+O(|x|^3)$; (2) (A. Gray)
  show $\operatorname{vol}B(p,r)=\omega_nr^n\big(1-\tfrac{\mathrm{scal}(p)}{6(n+2)}r^2
  +O(r^3)\big)$, $\omega_n=\operatorname{vol}(B(0,1)\subset\mathbb R^n)$, by
  expanding the volume integral in polar coordinates using (1).
  \source{petersen:page-0246}
\end{remark}
